% Leçon 25 — Expression orale : exode rural — Chinois 1ère A — Cours signé OnBuch+
% Programme officiel MINESEC. Compiler avec : tectonic ZHA25-oral-exode-rural.tex
\newcommand{\DOCMATIERE}{Chinois}
\newcommand{\DOCNIVEAU}{1ère A}
\newcommand{\DOCMODULE}{Module 3 : Citoyenneté et ouverture au monde}
\newcommand{\DOCLECON}{ZHA25}
\newcommand{\DOCTITRE}{Leçon 25 — Expression orale : exode rural}
% OnBuch+ — préambule « cours signés » ENRICHI (compilé avec tectonic / XeLaTeX)
% Système visuel « Pop » d'OnBuch+ : fond crème (jamais blanc), bordures encre
% épaisses, ombres « dures » décalées, titres Archivo Black en capitales.
% Version MATHÉMATIQUES : graphes (pgfplots/tikz), boîtes pédagogiques
% (définition, propriété, méthode, attention, le savais-tu, exercice résolu,
% exercice, TP, bilan), page de garde et signature OnBuch+ ancrée en bas.
%
% Macros attendues dans main.tex AVANT \input{preamble} :
%   \DOCMATIERE  \DOCNIVEAU  \DOCMODULE  \DOCLECON (numéro)  \DOCTITRE
\documentclass[11pt]{article}

\usepackage{fontspec}
\usepackage[french]{babel} % français = langue principale ; le chinois via \zh{..} / textebox (xeCJK)
\usepackage{xeCJK}
\usepackage{pifont} % \ding{51} \ding{55} utilisés par les modèles
\usepackage[a4paper,margin=2cm,top=3.4cm,bottom=2.8cm,headheight=1.4cm,footskip=1.3cm]{geometry}
\usepackage{amsmath,amssymb,amsfonts}
\usepackage{mathtools} % \xmapsto, \coloneqq... souvent utilisés par les modèles
\usepackage{cancel} % simplifications barrées (bilans de forces)
% \abs{z} (module, valeur absolue) et \norm{u} (norme), utilisés par les modèles
\providecommand{\abs}{}\let\abs\relax\DeclarePairedDelimiter{\abs}{\lvert}{\rvert}
\providecommand{\norm}{}\let\norm\relax\DeclarePairedDelimiter{\norm}{\lVert}{\rVert}
\usepackage[table]{xcolor} % [table] : fournit aussi \rowcolors (alternance de lignes)
\usepackage{pagecolor}
\usepackage{tikz}
\usetikzlibrary{arrows.meta,positioning,calc,shapes.geometric,shapes.misc,shapes.arrows,decorations.pathmorphing,decorations.pathreplacing,decorations.markings,patterns,shadows,mindmap,trees,fit,backgrounds,matrix,intersections,angles,quotes}
\usepackage{pgfplots}
\usepackage[version=4]{mhchem} % équations nucléaires \ce{^{235}_{92}U}
\usepackage[european,siunitx]{circuitikz} % schémas électriques
\pgfplotsset{compat=1.18}
\usepgfplotslibrary{fillbetween,groupplots} % aires entre courbes (calcul intégral)
\usepackage{enumitem}
\usepackage{fancyhdr}
% [most] charge toutes les bibliothèques tcolorbox (listings, minted, raster,
% documentation...) et gonfle inutilement la mémoire TeX sur un document long
% (~300 boîtes) : on ne charge que ce qui est réellement utilisé.
\usepackage[skins,breakable]{tcolorbox}
\usepackage{graphicx}
\usepackage{colortbl}
\usepackage{booktabs}
\usepackage{tabularx}
\usepackage{diagbox} % case d'en-tête diagonale des tableaux à double entrée
% Colonnes extensibles centrée / gauche / droite (C, L, R), souvent utilisées par les modèles
\newcolumntype{C}{>{\centering\arraybackslash}X}
\newcolumntype{L}{>{\raggedright\arraybackslash}X}
\newcolumntype{R}{>{\raggedleft\arraybackslash}X}
\usepackage{array}
\usepackage{multicol}
\usepackage{multirow}
\usepackage{siunitx}
% parse-numbers=false : accepte aussi \SI{2,5 \times 10^{-3}}{...}
\sisetup{locale=FR,output-decimal-marker={,},group-minimum-digits=4,parse-numbers=false}
\DeclareSIUnit\ohm{\ensuremath{\symup{\Omega}}} % U+2126 absent de la police
\DeclareSIUnit\division{div} % réglages d'oscilloscope (ms/div, V/div)
\DeclareSIUnit\tr{tr} % tours (vitesse de rotation, ex. \SI{1500}{\tr\per\minute})
\DeclareSIUnit\molar{M} % concentration molaire (ex. \SI{3}{\molar}, \SI{1}{\micro\molar})
\DeclareSIUnit\an{an} % année (le modèle confond parfois avec \year, un compteur TeX) — ex. \SI{5}{\kilo\meter\per\an}
\DeclareSIUnit\FCFA{FCFA} % franc CFA (ex. \SI{500}{\FCFA\per\kilo\gram})
\DeclareSIUnit\degreeBrix{\textdegree Brix} % teneur en sucre d'un sirop/jus (réfractométrie)
\DeclareSIUnit\month{mois} % le modèle utilise parfois \month, un compteur TeX, comme unité de durée
\usepackage{qrcode}
% \degree : commande fréquemment utilisée par les modèles pour un angle ou une
% température en dehors de \SI{}{} (non fournie par les paquets chargés).
\providecommand{\degree}{^{\circ}}
% \qed (fin de démonstration), utilisé par les modèles sans amsthm
\providecommand{\qed}{\hfill$\square$}
% \barre : double barre des valeurs interdites dans un tableau de variations (tkz-tab)
\providecommand{\barre}{\ensuremath{\|}}
% environnement proof (amsthm non chargé) utilisé par les modèles
\ifdefined\proof\else\newenvironment{proof}[1][Démonstration]{\par\noindent\textit{#1.}\ }{\qed\par}\fi
\usepackage{lastpage}
\usepackage{hyperref}
\hypersetup{hidelinks}

% ============================================================
% Palette « Pop » OnBuch+ — mêmes hex que la classe P (plus_ui.dart)
% ============================================================
\definecolor{poporange}{HTML}{F59321}
\definecolor{popdark}{HTML}{E07A0C}
\definecolor{popcream}{HTML}{FFF6E8}
\definecolor{popink}{HTML}{1C1714}
\definecolor{poppurple}{HTML}{7A5AE0}
\definecolor{popgreen}{HTML}{2E9E6A}
\definecolor{poppink}{HTML}{E0457B}
\definecolor{popblue}{HTML}{2F7DE1}
\definecolor{popgold}{HTML}{C98A12}
\definecolor{popmuted}{HTML}{8A7F76}
\definecolor{popline}{HTML}{EADFD2}
\definecolor{popgray}{HTML}{8A7F76}
\colorlet{popgrey}{popgray}
% Alias tolérants (le modèle nomme parfois les couleurs différemment)
\colorlet{popwhite}{white}
\colorlet{popblack}{popink}
\colorlet{popred}{poppink}
\colorlet{popyellow}{popgold}
% Teintes claires (fonds de tableaux/encadrés) — le modèle nomme parfois
% les couleurs avec un suffixe L (ex. popblueL) sans les avoir définies.
\colorlet{poporangeL}{poporange!20}
\colorlet{popdarkL}{popdark!20}
\colorlet{poppurpleL}{poppurple!20}
\colorlet{popgreenL}{popgreen!20}
\colorlet{poppinkL}{poppink!20}
\colorlet{popblueL}{popblue!20}
\colorlet{popgoldL}{popgold!20}
\colorlet{popredL}{popred!20}
\colorlet{popgrayL}{popgray!20}
% Teintes claires pour fonds de tableaux / zones
\colorlet{popblueL}{popblue!10!white}
\colorlet{poporangeL}{poporange!14!white}
\colorlet{popgreenL}{popgreen!12!white}
\colorlet{poppurpleL}{poppurple!12!white}
\colorlet{poppinkL}{poppink!12!white}
\colorlet{popgoldL}{popgold!14!white}

\pagecolor{popcream}

% ============================================================
% Typographie
% ============================================================
\defaultfontfeatures{Ligatures=TeX}
\setmainfont{PlusJakartaSans-Medium.ttf}[Path=../../../fonts/,BoldFont=PlusJakartaSans-Bold.ttf]
% Symboles, API et emojis absents de la police principale : repli sur DejaVu Sans (caractères actifs)
\newfontface\symfont{DejaVuSans.ttf}[Path=../../../fonts/]
\makeatletter
\newcommand\symactive[1]{\catcode#1=13 \begingroup\lccode`\~=#1 \lowercase{\endgroup\def~}{{\symfont\char#1}}}
\newcommand\symblank[1]{\catcode#1=13 \begingroup\lccode`\~=#1 \lowercase{\endgroup\def~}{}}
\newcommand\symhyph[1]{\catcode#1=13 \begingroup\lccode`\~=#1 \lowercase{\endgroup\def~}{\mbox{-}}}
\newcount\symc
\newcommand\symrange[2]{\symc=#1 \@whilenum\symc<#2 \do{\expandafter\symactive\expandafter{\the\symc}\advance\symc by 1 }}
\newcommand\symblankrange[2]{\symc=#1 \@whilenum\symc<#2 \do{\expandafter\symblank\expandafter{\the\symc}\advance\symc by 1 }}
\makeatother
\symrange{"0250}{"02FF}\symactive{"2713}\symactive{"2714}\symactive{"2717}\symactive{"2718}\symactive{"2610}\symactive{"2611}\symactive{"2612}
\symrange{"2600}{"27BF}
\catcode"2705=13 \begingroup\lccode`\~="2705 \lowercase{\endgroup\def~}{{\symfont\char"2713}}\symblank{"26BD}\symblank{"26A1}
\symrange{"2460}{"257F}\symactive{"223C}\symactive{"2248}\symactive{"2260}
\symactive{"01F8}\symactive{"01F9}\symactive{"1E3E}\symactive{"1E3F}
\symhyph{"2011}\symhyph{"2E1A}\symblankrange{"1F300}{"1FAFF}\symblank{"FE0F}
% Caractères chinois : WenQuanYi Zen Hei (sous-ensemble GB2312 + pinyin), gras/italique simulés
\setCJKmainfont{WenQuanYiZenHei-subset.ttf}[Path=../../../fonts/,AutoFakeBold=3,AutoFakeSlant=0.2]
\xeCJKsetup{CJKspace=false}
% Pinyin : voyelles à caron (ǎ ǐ ǒ ǔ ǖ ǘ ǚ ǜ) absentes de la police principale -> caractères actifs
\newfontface\pyfont{WenQuanYiZenHei-subset.ttf}[Path=../../../fonts/]
\newcommand\pyactive[1]{\catcode#1=13 \begingroup\lccode`\~=#1 \lowercase{\endgroup\def~}{{\pyfont\char#1}}}
\pyactive{"01CD}\pyactive{"01CE}\pyactive{"01CF}\pyactive{"01D0}\pyactive{"01D1}\pyactive{"01D2}\pyactive{"01D3}\pyactive{"01D4}\pyactive{"01D5}\pyactive{"01D6}\pyactive{"01D7}\pyactive{"01D8}\pyactive{"01D9}\pyactive{"01DA}\pyactive{"01DB}\pyactive{"01DC}
% Maths en sans-serif (Fira Math) avec chiffres et lettres droites de Plus
% Jakarta, pour que les formules se fondent dans le texte.
\usepackage[math-style=ISO,bold-style=ISO]{unicode-math}
\setmathfont{FiraMath-Regular.otf}[Path=../../../fonts/]
\setmathfont{PlusJakartaSans-Medium.ttf}[Path=../../../fonts/,range={up/num,up/{Latin,latin},\mathup}]
\usepackage{icomma} % virgule décimale sans espace en mode math
% Caractères mathématiques Unicode absents des polices de texte (Plus Jakarta,
% Archivo Black) : rendus par leur équivalent mathématique, en texte comme en
% formule — sinon ils s'affichent en carré vide (ex. « Arithmétique dans ℤ »).
\usepackage{newunicodechar}
\newunicodechar{ℝ}{\ensuremath{\mathbb{R}}}
\newunicodechar{ℤ}{\ensuremath{\mathbb{Z}}}
\newunicodechar{ℂ}{\ensuremath{\mathbb{C}}}
\newunicodechar{ℕ}{\ensuremath{\mathbb{N}}}
\newunicodechar{ℚ}{\ensuremath{\mathbb{Q}}}
\newunicodechar{𝔻}{\ensuremath{\mathbb{D}}}
\newunicodechar{α}{\ensuremath{\alpha}}
\newunicodechar{β}{\ensuremath{\beta}}
\newunicodechar{γ}{\ensuremath{\gamma}}
\newunicodechar{δ}{\ensuremath{\delta}}
\newunicodechar{ε}{\ensuremath{\varepsilon}}
\newunicodechar{θ}{\ensuremath{\theta}}
\newunicodechar{λ}{\ensuremath{\lambda}}
\newunicodechar{μ}{\ensuremath{\mu}}
\newunicodechar{σ}{\ensuremath{\sigma}}
\newunicodechar{τ}{\ensuremath{\tau}}
\newunicodechar{φ}{\ensuremath{\varphi}}
\newunicodechar{ψ}{\ensuremath{\psi}}
\newunicodechar{ω}{\ensuremath{\omega}}
\newunicodechar{Σ}{\ensuremath{\Sigma}}
% majuscules produites par \MakeUppercase dans les titres (α-aminés -> Α)
\newunicodechar{Α}{\ensuremath{\alpha}}
\newunicodechar{Β}{\ensuremath{\beta}}
\newunicodechar{Γ}{\ensuremath{\Gamma}}
\newunicodechar{Δ}{\ensuremath{\Delta}}
\newunicodechar{Ε}{\ensuremath{\varepsilon}}
\newunicodechar{Θ}{\ensuremath{\Theta}}
\newunicodechar{Λ}{\ensuremath{\Lambda}}
\newunicodechar{Μ}{\ensuremath{\mu}}
\newunicodechar{Τ}{\ensuremath{\tau}}
\newunicodechar{Φ}{\ensuremath{\Phi}}
\newunicodechar{Ψ}{\ensuremath{\Psi}}
\newunicodechar{Ω}{\ensuremath{\Omega}}
\newunicodechar{↦}{\ensuremath{\mapsto}}
\newunicodechar{∈}{\ensuremath{\in}}
\newunicodechar{∉}{\ensuremath{\notin}}
\newunicodechar{∧}{\ensuremath{\wedge}}
\newunicodechar{⇔}{\ensuremath{\Leftrightarrow}}
\newunicodechar{⟺}{\ensuremath{\Longleftrightarrow}}
\newunicodechar{⇒}{\ensuremath{\Rightarrow}}
\newunicodechar{⟹}{\ensuremath{\Longrightarrow}}
\newunicodechar{∘}{\ensuremath{\circ}}
\newunicodechar{∪}{\ensuremath{\cup}}
\newunicodechar{∩}{\ensuremath{\cap}}
\newunicodechar{≡}{\ensuremath{\equiv}}
\newunicodechar{⊂}{\ensuremath{\subset}}
\newunicodechar{⊥}{\ensuremath{\perp}}
\newunicodechar{∃}{\ensuremath{\exists}}
\newunicodechar{∀}{\ensuremath{\forall}}
\newunicodechar{∛}{\ensuremath{\sqrt[3]{\,}}}
\newunicodechar{∜}{\ensuremath{\sqrt[4]{\,}}}
\newunicodechar{⃗}{\ensuremath{{}^{\to}}}
\newunicodechar{ⁿ}{\ifmmode^{n}\else\textsuperscript{\ensuremath{n}}\fi}
\newunicodechar{ˣ}{\ifmmode^{x}\else\textsuperscript{\ensuremath{x}}\fi}
\newunicodechar{ᵇ}{\ifmmode^{b}\else\textsuperscript{\ensuremath{b}}\fi}
\newunicodechar{ʳ}{\ifmmode^{r}\else\textsuperscript{\ensuremath{r}}\fi}
\newunicodechar{ᵘ}{\ifmmode^{u}\else\textsuperscript{\ensuremath{u}}\fi}
\newunicodechar{ʸ}{\ifmmode^{y}\else\textsuperscript{\ensuremath{y}}\fi}
\newunicodechar{ᵃ}{\ifmmode^{a}\else\textsuperscript{\ensuremath{a}}\fi}
\newunicodechar{ᵉ}{\ifmmode^{e}\else\textsuperscript{\ensuremath{e}}\fi}
\newunicodechar{ᵏ}{\ifmmode^{k}\else\textsuperscript{\ensuremath{k}}\fi}
\newunicodechar{ᵖ}{\ifmmode^{p}\else\textsuperscript{\ensuremath{p}}\fi}
\newunicodechar{ᴺ}{\ifmmode^{N}\else\textsuperscript{\ensuremath{N}}\fi}
\newunicodechar{ᶜ}{\ifmmode^{c}\else\textsuperscript{\ensuremath{c}}\fi}
\newunicodechar{ʰ}{\ifmmode^{h}\else\textsuperscript{\ensuremath{h}}\fi}
\newunicodechar{ᵠ}{\ifmmode^{q}\else\textsuperscript{\ensuremath{q}}\fi}
\newunicodechar{ₙ}{\ifmmode_{n}\else\textsubscript{\ensuremath{n}}\fi}
\newunicodechar{ₐ}{\ifmmode_{a}\else\textsubscript{\ensuremath{a}}\fi}
\newunicodechar{ᵢ}{\ifmmode_{i}\else\textsubscript{\ensuremath{i}}\fi}
\newunicodechar{ᵣ}{\ifmmode_{r}\else\textsubscript{\ensuremath{r}}\fi}
\newunicodechar{ₖ}{\ifmmode_{k}\else\textsubscript{\ensuremath{k}}\fi}
\newunicodechar{ᵦ}{\ifmmode_{\beta}\else\textsubscript{\ensuremath{\beta}}\fi}
\newunicodechar{ₓ}{\ifmmode_{x}\else\textsubscript{\ensuremath{x}}\fi}
\newfontfamily\popdisplay{ArchivoBlack-Regular.ttf}[Path=../../../fonts/]
\newfontfamily\poplogo{SpaceGrotesk-Bold.ttf}[Path=../../../fonts/]
\newfontfamily\popsemi{PlusJakartaSans-SemiBold.ttf}[Path=../../../fonts/]
\newfontfamily\popxbold{PlusJakartaSans-ExtraBold.ttf}[Path=../../../fonts/]
\newfontfamily\popxboldls{PlusJakartaSans-ExtraBold.ttf}[Path=../../../fonts/,LetterSpace=3.0]

\setlist[enumerate]{leftmargin=1.7em,itemsep=5pt,topsep=4pt}
\setlist[itemize]{leftmargin=1.7em,itemsep=4pt,topsep=4pt,label={\color{poporange}\textbullet}}
\setlength{\parindent}{0pt}
\setlength{\parskip}{5pt}
\renewcommand{\arraystretch}{1.25}

% pgfplots : style Pop par défaut
\pgfplotsset{
  every axis/.append style={axis line style={popink,line width=1pt},tick style={popink},
    grid style={popline,line width=0.5pt},label style={font=\small},tick label style={font=\footnotesize}},
  popcurve/.style={poporange,line width=1.8pt,no markers,smooth},
}
% Même style disponible dans un \draw TikZ simple (hors pgfplots)
\tikzset{popcurve/.style={poporange,line width=1.8pt}}
\tikzset{popgrid/.style={popline,very thin}}
\colorlet{popcurve}{poporange} % le nom du style est parfois utilisé comme couleur

% ============================================================
% Pill / squiggle
% ============================================================
\newcommand{\poppill}[3]{%
  \tikz[baseline=-0.55ex]\node[draw=popink,line width=1.2pt,rounded corners=2.6pt,%
    fill=#2,inner xsep=6.5pt,inner ysep=3pt]{%
    \popxboldls\fontsize{7}{7}\selectfont\color{#3}\MakeUppercase{#1}};%
}
\newcommand{\popsquiggle}[1]{%
  \tikz[baseline=-0.4ex]{
    \draw[#1,line width=2.6pt,line cap=round]
      plot [smooth] coordinates {(0,0.09) (0.34,0.01) (0.68,0.21) (1.02,0.01) (1.36,0.10)};
  }%
}
% Mot-clé surligné (marqueur orange sous le texte)
\newcommand{\cle}[1]{{\popxbold\color{popdark}#1}}

% ============================================================
% En-tête / pied
% ============================================================
\pagestyle{fancy}
\fancyhf{}
\renewcommand{\headrulewidth}{0pt}
\renewcommand{\footrulewidth}{0pt}
\lhead{\raisebox{-0.15ex}{\poplogo\fontsize{13}{13}\selectfont\color{popink}OnBuch\color{poporange}+}}
\rhead{\poppill{\DOCMATIERE\ \textbullet\ \DOCNIVEAU\ \textbullet\ Leçon \DOCLECON}{white}{popink}}
\lfoot{\popxbold\fontsize{7.6}{7.6}\selectfont\color{popmuted}\MakeUppercase{Cours signé OnBuch+}\ \textbullet\ {\popsemi\DOCTITRE}}
\rfoot{\tikz[baseline=-0.6ex]\node[rounded corners=8.5pt,fill=popink,minimum height=17pt,minimum width=17pt,inner xsep=5pt,inner sep=0]{\popxbold\fontsize{8}{8}\selectfont\color{popcream}\thepage\,/\,\pageref*{LastPage}};}

% ============================================================
% Titres
% ============================================================
\newcommand{\chaptitre}[2]{%
  \begin{tcolorbox}[enhanced,colback=poporange,colframe=popink,boxrule=2.6pt,arc=11pt,
      drop shadow=popink,left=15pt,right=15pt,top=12pt,bottom=11pt,before skip=3pt,after skip=18pt]
    \poppill{#2}{popink}{popcream}\par\vspace{9pt}
    {\raggedright\hyphenpenalty=10000\exhyphenpenalty=10000\popdisplay\fontsize{22}{24}\selectfont\color{popcream}\MakeUppercase{#1}\par}
  \end{tcolorbox}
}
% Section numérotée : pastille orange avec le numéro + titre Archivo
\newcounter{popsec}
\newcommand{\coursec}[1]{%
  \stepcounter{popsec}\setcounter{popsub}{0}%
  \par\vspace{16pt}\noindent
  \begin{minipage}{\linewidth}
  \tikz[baseline=-0.7ex]\node[draw=popink,line width=1.6pt,fill=poporange,rounded corners=4pt,minimum size=22pt,inner sep=0]{\popdisplay\fontsize{12}{12}\selectfont\color{popcream}\thepopsec};%
  \hspace{8pt}{\popdisplay\fontsize{15}{17}\selectfont\color{popink}\MakeUppercase{#1}}\par
  \vspace{4pt}\noindent{\color{poporange}\rule{36pt}{3.6pt}}
  \end{minipage}\par\nopagebreak\vspace{8pt}
}
\newcounter{popsub}
\newcommand{\courssub}[1]{%
  \stepcounter{popsub}%
  \par\vspace{9pt}\noindent{\popxbold\fontsize{11.5}{13}\selectfont\color{popink}{\color{poporange}\thepopsec.\thepopsub}\hspace{6pt}#1}\par\nopagebreak\vspace{4pt}
}

% ============================================================
% Boîtes pédagogiques — même recette : fond blanc, bordure épaisse colorée,
% ombre dure encre, badge plein. Titre optionnel [#1].
% ============================================================
\tcbset{popbase/.style={breakable,enhanced,colback=white,boxrule=2.3pt,arc=9pt,
  shadow={3pt}{-3pt}{0pt}{popink},left=14pt,right=14pt,top=11pt,bottom=11pt,before skip=11pt,after skip=15pt,
  fontupper=\popsemi\fontsize{10.6}{14.5}\selectfont\color{popink},
  fontlower=\popsemi\fontsize{10.6}{14.5}\selectfont\color{popink}}}
% Coche dessinée (le glyphe ✓ n'existe pas dans les polices du document)
\newcommand{\popcheck}[1]{\tikz[baseline=-0.5ex]\draw[#1,line width=1.6pt,line cap=round,line join=round](-0.13,0.0)--(-0.04,-0.09)--(0.13,0.1);}
\providecommand{\coche}{\popcheck{popgreen}}
\providecommand{\resultat}[1]{\par\smallskip\noindent\fcolorbox{popgreen}{popgreenL}{\parbox{\dimexpr\linewidth-2\fboxsep-2\fboxrule\relax}{\textbf{Résultat.} #1}}\par\smallskip}
\newcommand{\popboxhead}[3]{\poppill{#1}{#2}{white}\ifx\relax#3\relax\else\hspace{6pt}{\popxbold\fontsize{10.6}{12}\selectfont\color{popink}#3}\fi\par\vspace{7pt}}

\newtcolorbox{aretenir}[1][]{popbase,colframe=popgreen,before upper={\popboxhead{À retenir}{popgreen}{#1}}}
% Texte en chinois (document de lecture, dialogue, extrait) : cadre
% d'étude. Un titre optionnel (source, auteur) s'affiche dans l'en-tête.
\newtcolorbox{textebox}[1][]{popbase,colframe=popblue,before upper={\popboxhead{文本 · Texte}{popblue}{#1}},after upper={}}
% Marques ✓ / ✗ (forme correcte / fautive) souvent utilisées par les modèles
\providecommand{\cross}{\textcolor{poppink}{\ensuremath{\times}}}
\providecommand{\xmark}{\cross}\providecommand{\crossmark}{\cross}\providecommand{\cmark}{\ensuremath{\checkmark}}
% Ligne de réponse / justification à compléter (inventée par les modèles)
\providecommand{\justification}[1]{\par\noindent\textit{Justification~:} \dotfill\par}
% \textipa (package tipa non chargé) : la police ne couvre pas l'API -> simple passage
\providecommand{\textipa}[1]{#1}
% Soulignement (ulem non chargé) : \uline, \ul
\providecommand{\uline}[1]{\underline{#1}}\providecommand{\ul}[1]{\underline{#1}}
% \glossaire{\item ...} inventée par les modèles : liste de glossaire
\providecommand{\glossaire}[1]{\begin{itemize}#1\end{itemize}}
% \alert (beamer) inventée par les modèles : texte mis en évidence
\providecommand{\alert}[1]{\textbf{\textcolor{poppink}{#1}}}
% variantes ASCII d'ordinaux écrites en macro
\providecommand{\iere}{\textsuperscript{re}}\providecommand{\ieme}{\textsuperscript{e}}\providecommand{\ere}{\textsuperscript{re}}\providecommand{\eme}{\textsuperscript{e}}\providecommand{\er}{\textsuperscript{er}}\providecommand{\ie}{\textsuperscript{e}}
% Ordinaux français écrits en macro par les modèles : 1\ière, 3\ième
\newcommand{\ière}{\textsuperscript{re}}\newcommand{\ième}{\textsuperscript{e}}
% \exemple (inventée par les modèles) : étiquette « Exemple : »
\providecommand{\exemple}{\textbf{Exemple~:}\ }
% \square utilisable aussi hors mode math (cases à cocher)
\AtBeginDocument{\let\squareMath\square\renewcommand{\square}{\ensuremath{\squareMath}}}
% Mot ou phrase en chinois à l'intérieur d'un paragraphe français
\newcommand{\zh}[1]{\ifmmode\text{#1}\else{#1}\fi}
\let\eng\zh \let\esp\zh \let\all\zh % alias tolérés
% flèches d'intonation inventées par les modèles
\providecommand{\textsearrow}{}\renewcommand{\textsearrow}{\ensuremath{\searrow}}\providecommand{\textnearrow}{}\renewcommand{\textnearrow}{\ensuremath{\nearrow}}
% \fr{..} : mot français (inventé par les modèles)
\providecommand{\fr}[1]{\textit{#1}}
% zhbox : encadré de texte chinois inventé par les modèles
\newenvironment{zhbox}{\begin{quote}}{\end{quote}}
% \courssubsub : titre de niveau 3 inventé par les modèles
\providecommand{\courssubsub}[1]{\par\medskip\noindent\textbf{#1}\par\smallskip}
% \zhbf : texte chinois en gras inventé par les modèles
\providecommand{\zhbf}[1]{\textbf{#1}}
% \vs : « versus » inventé par les modèles
\providecommand{\vs}{\textit{vs}\ }
% \blank : case à compléter inventée par les modèles
\providecommand{\blank}[1][]{\underline{\hspace{1.6em}}}
\newtcolorbox{exemplebox}[1][]{popbase,colframe=poporange,before upper={\popboxhead{Exemple}{poporange}{#1}}}
\newtcolorbox{definition}[1][]{popbase,colframe=popblue,before upper={\popboxhead{Définition}{popblue}{#1}}}
\newtcolorbox{propriete}[1][]{popbase,colframe=poppurple,before upper={\popboxhead{Propriété}{poppurple}{#1}}}
\newtcolorbox{methode}[1][]{popbase,colframe=popgold,before upper={\popboxhead{Méthode}{popgold}{#1}}}
\newtcolorbox{attention}[1][]{popbase,colframe=poppink,before upper={\popboxhead{Attention}{poppink}{#1}}}
\newtcolorbox{experience}[1][]{popbase,colframe=popblue,colback=popblueL,before upper={\popboxhead{Expérience}{popblue}{#1}}}
% « Le savais-tu ? » : carte violette pleine (contexte, culture, Cameroun)
\newtcolorbox{savaistu}[1][]{popbase,colback=poppurple,colframe=popink,
  fontupper=\popsemi\fontsize{10.4}{14.2}\selectfont\color{white},
  before upper={\poppill{Le savais-tu\,?}{white}{poppurple}\ifx\relax#1\relax\else\hspace{6pt}{\popxbold\color{white}#1}\fi\par\vspace{7pt}}}
% Exercice résolu : énoncé en haut, \tcblower puis la solution sur fond crème
\newcounter{popexo}\newcounter{popexr}
\newtcolorbox{exoresolu}[1][]{popbase,colframe=poporange,colbacklower=popcream,
  segmentation style={popink,line width=1.2pt,dashed},
  before upper={\stepcounter{popexr}\popboxhead{Exercice résolu \thepopexr}{poporange}{#1}},
  before lower={\poppill{Solution}{popink}{popcream}\par\vspace{6pt}}}
% Exercice (non résolu en place) — la correction est dans la partie Corrigés
\newtcolorbox{exercice}[1][]{popbase,colframe=popink,
  before upper={\stepcounter{popexo}\popboxhead{Exercice \thepopexo}{popink}{#1}}}
% Corrigé
\newtcolorbox{corrige}[1][]{popbase,colframe=popgreen,colback=popgreenL,
  before upper={\popboxhead{Corrigé}{popgreen}{#1}}}
% Objectifs / prérequis / bilan
\newtcolorbox{objectifs}{popbase,colframe=popink,colback=poporangeL,
  before upper={\popboxhead{Objectifs de la leçon}{popink}{}}}
\newtcolorbox{prerequis}{popbase,colframe=popink,colback=popblueL,
  before upper={\popboxhead{Prérequis}{popblue}{}}}
\newtcolorbox{bilan}{popbase,colframe=popink,colback=white,boxrule=2.8pt,
  before upper={\popboxhead{Fiche bilan}{poporange}{L'essentiel à connaître pour l'examen}}}
% Figure encadrée avec légende
\newtcolorbox{popfigure}[1][]{enhanced,breakable=false,colback=white,colframe=popline,boxrule=1.4pt,arc=8pt,
  left=10pt,right=10pt,top=10pt,bottom=8pt,before skip=10pt,after skip=12pt,halign=center,
  lower separated=false,halign lower=center,
  fontlower=\popsemi\fontsize{9}{11.5}\selectfont\color{popmuted},
  #1}
\newcommand{\legende}[1]{\par\vspace{6pt}{\popsemi\fontsize{9}{11.5}\selectfont\color{popmuted}#1}}

% ============================================================
% Signature OnBuch+
% ============================================================
\newcommand{\poplogoinline}[1]{{\poplogo\fontsize{#1}{#1}\selectfont\color{popink}OnBuch\color{poporange}+}}
\newcommand{\signatureonbuch}{%
  \begin{tcolorbox}[enhanced,colback=popink,colframe=popink,boxrule=2.2pt,arc=10pt,
      drop shadow=poporange,left=14pt,right=14pt,top=10pt,bottom=10pt,before skip=0pt,after skip=0pt]
    \tikz[baseline=-0.7ex]\node[circle,fill=popgreen,draw=popcream,line width=1.3pt,minimum size=22pt,inner sep=0]{\popcheck{white}};%
    \hspace{9pt}\begin{minipage}[c]{0.62\linewidth}
      {\popxbold\fontsize{12}{13}\selectfont\color{popcream}Signé\ \textbullet\ {\poplogo OnBuch\color{poporange}+}}\par\vspace{2pt}
      {\raggedright\popsemi\fontsize{8}{10.5}\selectfont\color{popline}\MakeUppercase{Équipe pédagogique OnBuch+}\\\MakeUppercase{Conforme au programme officiel MINESEC}\par}
    \end{minipage}\hfill
    {\poplogo\fontsize{22}{22}\selectfont\color{popcream}OnBuch\color{poporange}+}
  \end{tcolorbox}%
}

% ============================================================
% Page de garde
% #1 titre · #2 matière · #3 niveau · #4 module · #5 n° leçon · #6 durée ·
% #7 description / objectifs (texte ou itemize)
% ============================================================
\newcommand{\pagegarde}[7]{%
  \thispagestyle{empty}
  \newgeometry{a4paper,margin=1.7cm,top=1.6cm,bottom=1.4cm}%
  \begin{tikzpicture}[remember picture,overlay]
    \fill[poporange!18] ([xshift=-3.2cm,yshift=-3.6cm]current page.north east) circle (4.6cm);
    \fill[poppurple!14] ([xshift=2.4cm,yshift=7.6cm]current page.south west) circle (3.8cm);
  \end{tikzpicture}%
  {\poplogo\fontsize{30}{30}\selectfont\color{popink}OnBuch\color{poporange}+}\hfill
  \raisebox{0.6ex}{\poppill{Cours signé \textbullet\ Premium}{popink}{popcream}}\par
  \vspace{1pt}\popsquiggle{poppurple}\par
  \vspace{8pt}
  \begin{tcolorbox}[enhanced,colback=poporange,colframe=popink,boxrule=2.8pt,arc=13pt,
      drop shadow=popink,left=18pt,right=18pt,top=12pt,bottom=13pt,after skip=10pt]
    \poppill{#2 \textbullet\ #3}{popink}{popcream}\hspace{5pt}\poppill{#4}{white}{popink}\par\vspace{8pt}
    {\popdisplay\fontsize{14}{14}\selectfont\color{popink}LEÇON #5}\par\vspace{4pt}
    {\raggedright\hyphenpenalty=10000\exhyphenpenalty=10000\popdisplay\fontsize{25}{27}\selectfont\color{popcream}\MakeUppercase{#1}\par}
    \vspace{8pt}
    {\popxbold\fontsize{9.5}{9.5}\selectfont\color{popink}\MakeUppercase{Durée conseillée : #6}}
  \end{tcolorbox}
  \begin{tcolorbox}[enhanced,colback=white,colframe=popink,boxrule=2.2pt,arc=10pt,
      drop shadow=popink,left=16pt,right=16pt,top=10pt,bottom=10pt,after skip=8pt]
    \poppill{Dans cette leçon}{poppurple}{white}\par\vspace{5pt}
    \popsemi\fontsize{9.3}{12.6}\selectfont\color{popink}#7
  \end{tcolorbox}
  \noindent\begin{minipage}[t]{0.487\linewidth}
    \begin{tcolorbox}[enhanced,colback=popgreen,colframe=popink,boxrule=2.2pt,arc=10pt,
        drop shadow=popink,left=11pt,right=11pt,top=10pt,bottom=10pt,height=3.1cm,valign=center]
      \tcbox[colback=white,colframe=popink,boxrule=1.3pt,arc=5pt,boxsep=0pt,nobeforeafter,
        left=3pt,right=3pt,top=3pt,bottom=3pt,valign=center]{\qrcode[height=1.9cm]{https://play.google.com/store/apps/details?id=com.luvvix.onbuch&hl=fr}}%
      \hspace{8pt}\begin{minipage}[c]{0.5\linewidth}
        {\popdisplay\fontsize{11}{12.5}\selectfont\color{white}\MakeUppercase{Télécharge OnBuch}}\par\vspace{4pt}
        {\raggedright\popsemi\fontsize{8.4}{11}\selectfont\color{white}Gratuit \textbullet\ Google Play\\Tuteur IA, annales, concours, résultats\par}
      \end{minipage}
    \end{tcolorbox}
  \end{minipage}\hfill
  \begin{minipage}[t]{0.487\linewidth}
    \begin{tcolorbox}[enhanced,colback=poppurple,colframe=popink,boxrule=2.2pt,arc=10pt,
        drop shadow=popink,left=11pt,right=11pt,top=10pt,bottom=10pt,height=3.1cm,valign=center]
      \tikz[baseline=-0.7ex]\node[circle,fill=white,draw=popink,line width=1.3pt,minimum size=1.6cm,inner sep=0]{\popdisplay\fontsize{24}{24}\selectfont\color{poppurple}+};%
      \hspace{8pt}\begin{minipage}[c]{0.6\linewidth}
        {\popdisplay\fontsize{11}{12.5}\selectfont\color{white}\MakeUppercase{Passe à OnBuch+}}\par\vspace{4pt}
        {\raggedright\popsemi\fontsize{8.4}{11}\selectfont\color{white}Tous les cours signés, Tuteur IA illimité, fiches PDF — onglet OnBuch+ de l'app\par}
      \end{minipage}
    \end{tcolorbox}
  \end{minipage}
  \vspace{8pt}
  \popcontenu
  \vfill
  \signatureonbuch
  \restoregeometry
  \newpage
}

% Rangée « Ce cours contient » (page de garde)
\newcommand{\poptile}[3]{% #1 couleur #2 gros texte #3 légende
  \begin{tcolorbox}[enhanced,colback=#1,colframe=popink,boxrule=2pt,arc=9pt,shadow={2.5pt}{-2.5pt}{0pt}{popink},
      left=6pt,right=6pt,top=9pt,bottom=9pt,halign=center,valign=center,height=1.75cm,width=\linewidth,nobeforeafter]
    {\popdisplay\fontsize{12.5}{13}\selectfont\color{white}\MakeUppercase{#2}}\par\vspace{3pt}
    {\centering\popsemi\fontsize{7.8}{9.5}\selectfont\color{white}#3\par}
  \end{tcolorbox}}
\newcommand{\popcontenu}{%
  \par\noindent{\popxboldls\fontsize{7.5}{7.5}\selectfont\color{popmuted}CE COURS SIGNÉ CONTIENT}\par\vspace{7pt}
  \noindent\begin{minipage}[t]{0.235\linewidth}\poptile{popblue}{Cours}{complet, illustré, pas à pas}\end{minipage}\hfill
  \begin{minipage}[t]{0.235\linewidth}\poptile{poporange}{Méthodes}{et exercices résolus}\end{minipage}\hfill
  \begin{minipage}[t]{0.235\linewidth}\poptile{poppink}{Exercices}{type examen + corrigés}\end{minipage}\hfill
  \begin{minipage}[t]{0.235\linewidth}\poptile{popgreen}{Fiche bilan}{l'essentiel à retenir}\end{minipage}\par}

% Sommaire
\newtcolorbox{sommaire}{enhanced,colback=white,colframe=popink,boxrule=2.2pt,arc=10pt,
  drop shadow=popink,left=18pt,right=18pt,top=13pt,bottom=13pt,before skip=6pt,after skip=20pt,
  before upper={\poppill{Sommaire}{poppurple}{white}\par\vspace{9pt}\popsemi\fontsize{10.6}{14.5}\selectfont\color{popink}}}

% Signature de fin de cours — poussée tout en bas de la dernière page
\newcommand{\findecours}{%
  \par\vfill\nopagebreak
  \begin{center}\popsquiggle{poporange}\end{center}\vspace{4pt}
  \signatureonbuch
}
\AtBeginDocument{\def\llbracket{\mathopen{[\mkern-3.5mu[}}\def\rrbracket{\mathclose{]\mkern-3.5mu]}}} % intervalles d'entiers (glyphe ⟦ absent de la police)
% Glyphes absents de Fira Math : redéfinis à partir de glyphes disponibles
\AtBeginDocument{%
  \def\setminus{\mathbin{\backslash}}\let\smallsetminus\setminus
  \def\vdots{\mathord{\vbox{\baselineskip3.2pt\lineskiplimit0pt\kern4pt\hbox{.}\hbox{.}\hbox{.}}}}%
  \def\ddots{\mathinner{\mkern1mu\raise6.5pt\hbox{.}\mkern2mu\raise3.6pt\hbox{.}\mkern2mu\raise0.7pt\hbox{.}\mkern1mu}}%
  \def\triangle{\mathord{\scalebox{0.9}{$\symup{\Delta}$}}}%
  \def\nabla{\mathord{\raisebox{\depth}{\scalebox{1}[-1]{$\symup{\Delta}$}}}}%
  \def\checkmark{\popcheck{popdark}}%
  \def\star{\mathbin{\ast}}\def\bigstar{\mathord{\ast}}%
  \def\diamond{\mathbin{\scalebox{0.6}{$\Diamond$}}}%
  \def\bigcup{\mathop{\mathchoice{\raisebox{-0.3em}{\scalebox{1.7}{$\cup$}}}{\scalebox{1.25}{$\cup$}}{\cup}{\cup}}\displaylimits}%
  \def\bigcap{\mathop{\mathchoice{\raisebox{-0.3em}{\scalebox{1.7}{$\cap$}}}{\scalebox{1.25}{$\cap$}}{\cap}{\cap}}\displaylimits}%
  \def\lesseqqgtr{\mathrel{\lessgtr}}\def\bigoplus{\mathop{\oplus}\displaylimits}%
}
\providecommand{\ssi}{si et seulement si} % abréviation fréquente
\AtBeginDocument{\providecommand{\wideparen}{\overparen}} % arc de cercle

%% FIGFIX-BEGIN v5 — correctifs globaux des figures (docs/onbuch-generation/tools/figfix)
\usepackage{adjustbox}\usepackage{etoolbox}\tcbuselibrary{raster}
% toute figure TikZ/pgfplots plus large que la ligne est réduite à la largeur disponible
\BeforeBeginEnvironment{tikzpicture}{\begin{adjustbox}{max width=\linewidth}}
\AfterEndEnvironment{tikzpicture}{\end{adjustbox}}
% légendes pgfplots lisibles par-dessus les courbes
\pgfplotsset{every axis/.append style={legend style={fill=white,fill opacity=0.92,text opacity=1,draw=black!15,inner sep=3pt}}}
% étiquettes d'arêtes (syntaxe edge["..."]) avec fond blanc
\tikzset{every edge quotes/.append style={fill=white,inner sep=1.2pt}}
%% FIGFIX-END
\begin{document}

\pagegarde{Leçon 25 — Expression orale : exode rural}{Chinois}{1ère A}{Module 3 : Citoyenneté et ouverture au monde}{ZHA25}{2 h}{Leçon d'expression orale consacrée au thème de l'exode rural : les élèves apprennent le vocabulaire clé, les formules pour présenter un exposé, donner leur avis et relancer un débat, puis s'entraînent par dialogues, jeux de rôle et débats guidés, avec un travail spécifique sur les tons.
\vspace{4pt}
\begin{multicols}{2}\small
\begin{itemize}
\item À la fin de cette leçon, je sais utiliser le vocabulaire de base sur la ville, la campagne et la migration
\item À la fin de cette leçon, je sais présenter oralement un exposé simple et structuré en chinois
\item À la fin de cette leçon, je sais donner mon avis avec 我觉得 / 我认为 et le justifier avec 因为
\item À la fin de cette leçon, je sais être d'accord ou pas d'accord poliment (我同意 / 我不同意)
\item À la fin de cette leçon, je sais relancer un débat en posant des questions (你呢？你觉得怎么样？)
\item À la fin de cette leçon, je sais prononcer correctement les tons des mots clés de la leçon
\end{itemize}\end{multicols}}

\begin{sommaire}
\begin{multicols}{2}
\begin{enumerate}
\item Vocabulaire clé : ville, campagne, exode rural
\item Formules fonctionnelles : présenter, opiner, réagir, relancer
\item Dialogues modèles avec pinyin
\item Travail sur les tons et la prononciation
\item Jeux de rôle et débat guidé
\item Bilan et préparation de l'exposé évalué
\item Activité d'intégration
\item Méthodes et exercices résolus
\item Exercices
\item Corrigés des exercices
\item Fiche bilan
\end{enumerate}
\end{multicols}
\end{sommaire}

\begin{objectifs}
\begin{itemize}
\item À la fin de cette leçon, je sais utiliser le vocabulaire de base sur la ville, la campagne et la migration
\item À la fin de cette leçon, je sais présenter oralement un exposé simple et structuré en chinois
\item À la fin de cette leçon, je sais donner mon avis avec 我觉得 / 我认为 et le justifier avec 因为
\item À la fin de cette leçon, je sais être d'accord ou pas d'accord poliment (我同意 / 我不同意)
\item À la fin de cette leçon, je sais relancer un débat en posant des questions (你呢？你觉得怎么样？)
\item À la fin de cette leçon, je sais prononcer correctement les tons des mots clés de la leçon
\item À la fin de cette leçon, je sais comparer l'exode rural au Cameroun et en Chine en phrases simples
\item À la fin de cette leçon, je sais conclure un exposé avec une formule de clôture (谢谢大家！)
\end{itemize}
\end{objectifs}

\begin{prerequis}
\begin{itemize}
\item Se présenter et parler de sa famille et de sa ville (Seconde)
\item Les structures de base : 是, 有, 在, verbes modaux 想 / 要
\item Exprimer une opinion simple : 我觉得 + phrase
\item Les questions avec 吗, 呢, 为什么
\item Les quatre tons et les règles de sandhi du 3e ton
\end{itemize}
\end{prerequis}

\newpage

\coursec{Vocabulaire clé : ville, campagne, exode rural}

\textbf{Situation d'accroche.} À Yaoundé comme à Douala, de nombreux jeunes quittent leur village de l'Ouest ou du Nord pour chercher du travail en ville. On les croise à la gare routière de Mvan ou au carrefour Warda, un sac sur le dos, pleins d'espoir. En Chine aussi, des millions de personnes ont quitté les campagnes pour Shanghai ou Pékin : c'est le plus grand exode rural de l'histoire. Comment parler de ce phénomène en chinois ? Comment structurer un exposé, donner son avis et débattre avec ses camarades ? Cette leçon vous donne le lexique, les structures grammaticales, les formules fonctionnelles et la prononciation pour en discuter à l'oral.

\textbf{Problématique.} Quel vocabulaire et quelles structures faut-il maîtriser pour décrire le départ de la campagne vers la ville, ses causes, ses conséquences, et pour mener un échange oral structuré ?

\courssub{Observer : une phrase support à lire à voix haute}

Avant d'apprendre les mots un par un, lisons ensemble une phrase qui résume le phénomène. Écoutez le professeur, répétez en chœur, puis lisez seul.

\begin{textebox}[Phrase support de la leçon]
现在很多年轻人离开农村去城市，因为城市里有更多的工作机会。\\
Xiànzài hěn duō niánqīngrén líkāi nóngcūn qù chéngshì, yīnwèi chéngshì lǐ yǒu gèng duō de gōngzuò jīhuì.\\
« Aujourd'hui, beaucoup de jeunes quittent la campagne pour la ville, parce qu'en ville il y a plus d'opportunités de travail. »
\end{textebox}

Observons cette phrase. Elle contient presque tout le lexique de la leçon : \zh{年轻人} (les jeunes), \zh{离开} (quitter), \zh{农村} (la campagne), \zh{城市} (la ville), \zh{工作机会} (les opportunités de travail), et la structure de cause \zh{因为} (parce que). Nous allons maintenant démonter chaque pièce.

\courssub{Le lexique de base : la campagne et la ville}

\begin{center}
\begin{tabularx}{\linewidth}{|X|X|X|X|}
\hline
\rowcolor{popblueL} Caractères & Pinyin & Nature & Traduction \\
\hline
农村 & nóngcūn & nom & la campagne, les zones rurales \\
\hline
城市 & chéngshì & nom & la ville \\
\hline
乡下 & xiāngxia & nom & la campagne (familier) \\
\hline
农村人口 & nóngcūn rénkǒu & locution nominale & la population rurale \\
\hline
年轻人 & niánqīngrén & nom & les jeunes, les jeunes gens \\
\hline
农民工 & nóngmíngōng & nom & l'ouvrier migrant \\
\hline
\end{tabularx}
\end{center}

\begin{definition}[农村, 城市, 乡下 : trois mots pour deux réalités]
\zh{农村} \emph{nóngcūn} est le mot standard et neutre pour « la campagne » : il est composé de \zh{农} (agriculture, clé \zh{农}) et \zh{村} (village, clé \zh{木}). \zh{城市} \emph{chéngshì} désigne « la ville » : \zh{城} (ville fortifiée, clé \zh{土}) + \zh{市} (marché, clé \zh{巾}) : la ville, c'est historiquement le lieu du marché ! \zh{乡下} \emph{xiāngxia} signifie aussi « la campagne », mais dans un registre \cle{familier}, souvent avec une nuance affective (« mon bled », « mon village natal »). À l'oral, on dit volontiers : \zh{我老家在乡下。} \emph{Wǒ lǎojiā zài xiāngxia.} « Mon village natal est à la campagne. »
\end{definition}

\begin{exemplebox}[Premiers exemples contrastés]
\zh{农村很安静。} \emph{Nóngcūn hěn ānjìng.} « La campagne est calme. »\\
\zh{城市的生活很方便。} \emph{Chéngshì de shēnghuó hěn fāngbiàn.} « La vie en ville est pratique. »\\
\zh{农村人口越来越少。} \emph{Nóngcūn rénkǒu yuè lái yuè shǎo.} « La population rurale est de moins en moins nombreuse. »
\end{exemplebox}

\begin{attention}[Ne confondez pas 农村 et 乡下 à l'écrit formel]
Dans un exposé ou une copie, préférez \zh{农村}, plus neutre. \zh{乡下} convient très bien à l'oral et aux dialogues. Autre piège : ne dites pas \zh{*村} seul pour « la campagne » ; \zh{村} \emph{cūn} signifie « un village » précis, pas la campagne en général.
\end{attention}

\courssub{Le verbe 离开 : quitter un lieu}

\begin{propriete}[Construction de 离开 líkāi]
Structure : \cle{Sujet + 离开 + lieu (+ 去 + destination)}.\\
Le verbe \zh{离开} se construit \textbf{directement} avec le lieu qu'on quitte, \textbf{sans préposition}, contrairement au français « quitter DE » ou « partir DE ». Pour indiquer la destination, on ajoute \zh{去} + lieu.
\end{propriete}

\begin{exemplebox}[离开 en action]
\zh{他离开了农村。} \emph{Tā líkāi le nóngcūn.} « Il a quitté la campagne. »\\
\zh{很多年轻人离开农村去城市。} \emph{Hěn duō niánqīngrén líkāi nóngcūn qù chéngshì.} « Beaucoup de jeunes quittent la campagne pour aller en ville. »\\
\zh{她明年要离开雅温得去北京。} \emph{Tā míngnián yào líkāi Yǎwēndé qù Běijīng.} « L'année prochaine, elle va quitter Yaoundé pour Pékin. »
\end{exemplebox}

\begin{attention}[Erreur typique du francophone]
Ne traduisez pas « quitter de » par une préposition : on ne dit jamais \zh{*从农村离开} dans ce sens (cette phrase signifierait « partir depuis la campagne », avec un autre point de vue). La bonne formule est : \zh{离开农村}。 De même, « quitter la maison » = \zh{离开家}, « quitter le pays » = \zh{离开国家}.
\end{attention}

\courssub{Travail, opportunités et conditions de vie}

\begin{center}
\begin{tabularx}{\linewidth}{|X|X|X|X|}
\hline
\rowcolor{popblueL} Caractères & Pinyin & Nature & Traduction \\
\hline
找工作 & zhǎo gōngzuò & verbe + compl. & chercher du travail \\
\hline
机会 & jīhuì & nom & l'opportunité, la chance \\
\hline
生活 & shēnghuó & nom / verbe & la vie ; vivre \\
\hline
方便 & fāngbiàn & adjectif & pratique, commode \\
\hline
困难 & kùnnan & adj. / nom & difficile ; la difficulté \\
\hline
安静 & ānjìng & adjectif & calme, tranquille \\
\hline
贵 & guì & adjectif & cher \\
\hline
挤 & jǐ & adjectif & bondé, serré \\
\hline
想家 & xiǎng jiā & verbe + compl. & avoir le mal du pays \\
\hline
\end{tabularx}
\end{center}

\begin{exemplebox}[Pourquoi part-on ? Pourquoi hésite-t-on ?]
\zh{他去城市找工作。} \emph{Tā qù chéngshì zhǎo gōngzuò.} « Il va en ville chercher du travail. »\\
\zh{城市里有很多机会。} \emph{Chéngshì lǐ yǒu hěn duō jīhuì.} « En ville, il y a beaucoup d'opportunités. »\\
\zh{城市的生活很方便，可是很贵。} \emph{Chéngshì de shēnghuó hěn fāngbiàn, kěshì hěn guì.} « La vie en ville est pratique, mais chère. »\\
\zh{在农村，生活有时候很困难。} \emph{Zài nóngcūn, shēnghuó yǒu shíhou hěn kùnnan.} « À la campagne, la vie est parfois difficile. »\\
\zh{地铁很挤，他很想家。} \emph{Dìtiě hěn jǐ, tā hěn xiǎng jiā.} « Le métro est bondé, il a le mal du pays. »
\end{exemplebox}

\begin{definition}[农民工 nóngmíngōng : le mot clé du contexte chinois]
\zh{农民工} désigne l'\cle{ouvrier migrant} : un paysan (\zh{农}) devenu ouvrier (\zh{工}) en ville, tout en restant officiellement rattaché à son village (système \zh{户口} \emph{hùkǒu}). C'est le mot central pour parler de l'exode rural en Chine : \zh{农民工在城市工作，但是他们的家在农村。} \emph{Nóngmíngōng zài chéngshì gōngzuò, dànshì tāmen de jiā zài nóngcūn.} « Les ouvriers migrants travaillent en ville, mais leur famille est à la campagne. »
\end{definition}

\courssub{Deux structures grammaticales utiles}

\begin{propriete}[越来越 + adjectif : « de plus en plus »]
Structure : \cle{越来越 + adjectif}. Exprime une évolution progressive. L'adjectif après \zh{越来越} ne prend \textbf{jamais} \zh{很}.\\
\zh{越来越多} \emph{yuè lái yuè duō} « de plus en plus nombreux » ; \zh{越来越贵} \emph{yuè lái yuè guì} « de plus en plus cher » ; \zh{越来越方便} \emph{yuè lái yuè fāngbiàn} « de plus en plus pratique ».
\end{propriete}

\begin{propriete}[因为...所以... : exprimer la cause et la conséquence]
Structure : \cle{因为 + cause，所以 + conséquence}. Contrairement au français, qui interdit « parce que... donc... » dans la même phrase, le chinois \textbf{encourage} la double marque. À l'oral rapide, on peut garder \zh{因为} seul, mais dans un exposé soigné, utilisez les deux.
\end{propriete}

\begin{exemplebox}[Combiner les deux structures]
\zh{因为城市里有更多的工作机会，所以越来越多的年轻人离开农村。} \emph{Yīnwèi chéngshì lǐ yǒu gèng duō de gōngzuò jīhuì, suǒyǐ yuè lái yuè duō de niánqīngrén líkāi nóngcūn.} « Parce qu'en ville il y a plus d'opportunités de travail, de plus en plus de jeunes quittent la campagne. »\\
\zh{因为农村的生活很困难，所以他想搬家。} \emph{Yīnwèi nóngcūn de shēnghuó hěn kùnnan, suǒyǐ tā xiǎng bānjiā.} « Parce que la vie à la campagne est difficile, il veut déménager. »
\end{exemplebox}

\begin{attention}[Pièges à éviter]
1) Jamais \zh{*越来越很...} : dites \zh{越来越多}, pas \zh{*越来越很多}. 2) En français on ne dit pas « parce que... donc... », mais en chinois \zh{因为...所以...} est la forme la plus naturelle : ne supprimez pas \zh{所以} par calque du français. 3) \zh{越来越} se place devant l'adjectif, jamais après. 4) Attention au changement de ton (sandhi) sur \zh{所以} : \zh{suǒ} (3e ton) + \zh{yǐ} (3e ton) $\rightarrow$ \emph{suóyǐ} (2e + 3e ton).
\end{attention}

\courssub{Formules fonctionnelles pour l'exposé et le débat}

Pour réussir votre expression orale (exposé + débat), vous devez maîtriser ces « briques » communicatives. Apprenez-les par cœur.

\begin{center}
\begin{tabularx}{\linewidth}{|X|X|X|}
\hline
\rowcolor{popblueL} Fonction & Chinois (Pinyin) & Français \\
\hline
\multicolumn{3}{|c|}{\textbf{Ouvrir / Présenter le sujet}} \\ \hline
Accrocher l'auditoire & 大家好，今天我想谈谈…… \emph{Dàjiā hǎo, jīntiān wǒ xiǎng tántan……} & Bonjour à tous, aujourd'hui je veux parler de… \\
Annoncer le thème & 我的题目是：农村青年进城。 \emph{Wǒ de tímù shì: Nóngcūn qīngnián jìn chéng.} & Mon sujet est : les jeunes ruraux vont en ville. \\
\hline
\multicolumn{3}{|c|}{\textbf{Structurer / Enchaîner}} \\ \hline
Premièrement & 首先 \emph{shǒuxiān} & Premièrement \\
Deuxièmement & 其次 \emph{qícì} & Deuxièmement / En outre \\
Enfin & 最后 \emph{zuìhòu} & Enfin \\
D'un côté... de l'autre & 一方面……另一方面…… \emph{yī fāngmiàn…… lìng yī fāngmiàn……} & D'un côté… de l'autre côté… \\
\hline
\multicolumn{3}{|c|}{\textbf{Donner son avis / Argumenter}} \\ \hline
Je pense que & 我觉得 / 我认为 \emph{Wǒ juéde / Wǒ rènwéi} & Je pense que / À mon avis \\
Selon moi & 按我的看法 \emph{Àn wǒ de kànfǎ} & Selon mon point de vue \\
La raison est que & 原因是…… \emph{Yuányīn shì……} & La raison est que… \\
Parce que... donc... & 因为……所以…… \emph{Yīnwèi…… suǒyǐ……} & Parce que… donc… \\
\hline
\multicolumn{3}{|c|}{\textbf{Réagir / Relancer (Débat)}} \\ \hline
Je suis d'accord & 我同意。 \emph{Wǒ tóngyì.} / 对。 \emph{Duì.} & Je suis d'accord. / Exact. \\
Je ne suis pas d'accord & 我不同意。 \emph{Wǒ bù tóngyì.} / 我不觉得。 \emph{Wǒ bù juéde.} & Je ne suis pas d'accord. \\
Pourquoi ? & 为什么？ \emph{Wèishénme?} & Pourquoi ? \\
Et toi ? / Et vous ? & 你呢？ \emph{Nǐ ne?} / 您呢？ \emph{Nín ne?} & Et toi ? / Et vous ? \\
Je ne comprends pas & 我不太懂。 \emph{Wǒ bù tài dǒng.} & Je ne comprends pas bien. \\
Peux-tu répéter ? & 请再说一遍。 \emph{Qǐng zài shuō yī biàn.} & Pouvez-vous répéter ? \\
\hline
\multicolumn{3}{|c|}{\textbf{Conclure}} \\ \hline
En résumé & 总之 \emph{Zǒngzhī} / 综上所述 \emph{Zōngshàng shùshuō} & En résumé / En conclusion \\
Merci & 谢谢大家。 \emph{Xièxie dàjiā.} & Merci à tous. \\
\hline
\end{tabularx}
\end{center}

\courssub{Dialogue modèle : débat guidé}

Lisez ce dialogue à deux. Repérez les formules du tableau ci-dessus. Travaillez l'intonation montante pour les questions, descendante pour les affirmations.

\begin{textebox}[Dialogue : L'exode rural, une chance ou un problème ?]
\textbf{A :} 大家好。今天我们讨论农村青年进城。我觉得这是好事，因为城市有更多机会。\\
\emph{Dàjiā hǎo. Jīntiān wǒmen tǎolùn nóngcūn qīngnián jìn chéng. Wǒ juéde zhè shì hǎoshì, yīnwèi chéngshì yǒu gèng duō jīhuì.}\\
« Bonjour. Aujourd'hui nous discutons des jeunes ruraux qui vont en ville. Je pense que c'est une bonne chose, car la ville offre plus d'opportunités. »

\textbf{B :} 我不同意。城市生活很贵，而且很挤。很多农民工很想家。\\
\emph{Wǒ bù tóngyì. Chéngshì shēnghuó hěn guì, érqiě hěn jǐ. Hěn duō nóngmíngōng hěn xiǎng jiā.}\\
« Je ne suis pas d'accord. La vie en ville est chère, et en plus bondée. Beaucoup d'ouvriers migrants ont le mal du pays. »

\textbf{A :} 可是，城市医院好，学校也好。一方面方便，另一方面有发展。\\
\emph{Kěshì, chéngshì yīyuàn hǎo, xuéxiào yě hǎo. Yī fāngmiàn fāngbiàn, lìng yī fāngmiàn yǒu fāzhǎn.}\\
« Mais les hôpitaux et les écoles en ville sont bons. D'un côté c'est pratique, de l'autre il y a du développement. »

\textbf{B :} 你的意思是：利大于弊？\\
\emph{Nǐ de yìsi shì: lì dà yú bì?}\\
« Tu veux dire : les avantages l'emportent sur les inconvénients ? »

\textbf{A :} 对。总之，越来越多年轻人会离开农村。谢谢大家。\\
\emph{Duì. Zǒngzhī, yuè lái yuè duō niánqīngrén huì líkāi nóngcūn. Xièxie dàjiā.}\\
« Exact. En résumé, de plus en plus de jeunes quitteront la campagne. Merci à tous. »
\end{textebox}

\courssub{Travail sur les tons et la prononciation}

\begin{methode}[Entraînement phonétique ciblé (10 min)]
\textbf{Étape 1 : Sandhi du 3e ton (obligatoire).} Repérez les séquences 3e + 3e ton dans le dialogue et appliquez la règle : le premier devient 2e ton.
\begin{itemize}
    \item \zh{所以} \emph{suǒyǐ} $\rightarrow$ \emph{suóyǐ} (2-3)
    \item \zh{很好} \emph{hěn hǎo} $\rightarrow$ \emph{hén hǎo} (2-3) — \textit{rappel HSK1}
    \item \zh{想家} \emph{xiǎng jiā} $\rightarrow$ \emph{xiáng jiā} (2-1) — \zh{想} 3e ton, \zh{家} 1er ton $\rightarrow$ pas de sandhi ici, mais \zh{很想} \emph{hěn xiǎng} $\rightarrow$ \emph{hén xiǎng}.
\end{itemize}
\textbf{Étape 2 : Tons neutres et légers.} \zh{机会} \emph{jīhu\`i} (1 - ton neutre léger sur \zh{会}) ; \zh{农民工} \emph{nóngmíngōng} (2-2-1) ; \zh{乡下} \emph{xiāngxia} (1 - ton neutre).
\textbf{Étape 3 : Intonation de la phrase.}
\begin{itemize}
    \item Déclarative (\zh{我觉得……}) : courbe descendante globale.
    \item Question ouverte (\zh{为什么？}) : montante finale.
    \item Question fermée (\zh{对吗？}) : montante forte sur la particule \zh{吗}.
\end{itemize}
\textbf{Étape 4 : Exercice « Miroir ».} Le professeur lit une réplique du dialogue, la classe répète en chœur en imitant exactement la mélodie (prosodie), pas seulement les tons isolés.
\end{methode}

\begin{experience}[Atelier prononciation : paires minimales]
Travaillez par paires. L'élève A lit une phrase, l'élève B identifie le mot entendu (différence de ton).
\begin{enumerate}
    \item \zh{农村} (nóngcūn, 2-1) vs \zh{弄春} (nòng chūn, 4-1) — \textit{inventé pour l'exercice}
    \item \zh{城市} (chéngshì, 2-4) vs \zh{成事} (chéngshì, 2-4) — \textit{mêmes tons, sens diff.}
    \item \zh{机会} (jīhuì, 1-4) vs \zh{鸡灰} (jī huī, 1-1) — \textit{différence ton 4 / ton 1 sur 2e syllabe}
    \item \zh{离开} (líkāi, 2-1) vs \zh{里开} (lǐ kāi, 3-1) — \textit{différence ton 2 / ton 3 sur 1re syllabe}
\end{enumerate}
\emph{Correction collective au tableau : écriture du pinyin avec traits de tons.}
\end{experience}

\courssub{Carte mentale et schéma du phénomène}

\begin{popfigure}
\begin{tikzpicture}[mindmap, grow cyclic, every node/.style={concept, align=center}, concept color=popblue!40, root concept/.append style={font=\small\bfseries}, level 1/.append style={level distance=3.2cm, sibling angle=120, font=\small}, level 2/.append style={level distance=2.2cm, sibling angle=50, font=\scriptsize}]
\node[root concept, align=center]{农村 $\rightarrow$ 城市\\ \scriptsize exode rural}
  child[concept color=popgreen!40] { node[align=center]{原因\\ causes}
    child { node[align=center]{找工作\\ chercher travail} }
    child { node[align=center]{上学\\ études} }
    child { node[align=center]{机会多\\ opportunités} }
  }
  child[concept color=poporange!40] { node[align=center]{Avantages ville\\ 优点}
    child { node[align=center]{方便\\ pratique} }
    child { node[align=center]{医院好\\ bons hôpitaux} }
    child { node[align=center]{学校好\\ bonnes écoles} }
  }
  child[concept color=poppink!50] { node[align=center]{Problèmes\\ 缺点}
    child { node[align=center]{贵\\ cher} }
    child { node[align=center]{挤\\ bondé} }
    child { node[align=center]{想家\\ mal du pays} }
  };
\end{tikzpicture}
\legende{Carte mentale du lexique de l'exode rural : causes, avantages de la ville, problèmes.}
\end{popfigure}

\begin{popfigure}
\begin{tikzpicture}[node distance=2.6cm, >=Stealth, every node/.style={align=center}]
\node[draw=popgreen, fill=popgreenL, rounded corners, thick, minimum width=2.2cm, minimum height=1cm, align=center] (village){农村\\ \scriptsize nóngcūn};
\node[draw=popblue, fill=popblueL, rounded corners, thick, minimum width=2.2cm, minimum height=1cm, right=of village, align=center] (ville){城市\\ \scriptsize chéngshì};
\node[draw=poporange, fill=poporangeL, rounded corners, thick, minimum width=2.2cm, minimum height=1cm, right=of ville, align=center] (job){工作 / 机会\\ \scriptsize gōngzuò / jīhuì};
\draw[->, very thick, popdark] (village) -- node[above, font=\small]{离开 líkāi} (ville);
\draw[->, very thick, popdark] (ville) -- node[above, font=\small]{找 zhǎo / 有 yǒu} (job);
\end{tikzpicture}
\legende{Schéma fléché : on quitte la campagne, on va en ville, on cherche/trouve du travail.}
\end{popfigure}

\courssub{Comparaison des structures et entraînement écrit}

\begin{center}
\begin{tabularx}{\linewidth}{|X|X|X|}
\hline
\rowcolor{popblueL} Structure & Exemple en chinois + pinyin & Traduction \\
\hline
离开 + lieu (sans prép.) & 离开农村 (líkāi nóngcūn) & quitter la campagne \\
\hline
越来越 + adj. & 越来越多 (yuè lái yuè duō) & de plus en plus nombreux \\
\hline
因为...所以... & 因为有机会，所以他们离开。 (Yīnwèi yǒu jīhuì, suǒyǐ tāmen líkāi.) & Parce qu'il y a des opportunités, ils partent. \\
\hline
A 比 B + adj. (rappel) & 城市比农村方便。 (Chéngshì bǐ nóngcūn fāngbiàn.) & La ville est plus pratique que la campagne. \\
\hline
一方面...另一方面... & 一方面方便，另一方面贵。 (Yī fāngmiàn fāngbiàn, lìng yī fāngmiàn guì.) & D'un côté pratique, de l'autre cher. \\
\hline
\end{tabularx}
\end{center}

\begin{exoresolu}[Traduire en chinois (niveau exposé soigné)]
Énoncé : Traduisez en chinois : « De plus en plus de jeunes quittent la campagne parce que la vie en ville est pratique. »
\tcblower
Solution détaillée : 
1) « De plus en plus de jeunes » = \zh{越来越多的年轻人} (structure \zh{越来越} + adj \zh{多} + \zh{的} + nom). 
2) « Quittent la campagne » = \zh{离开农村} (verbe direct, pas de préposition). 
3) La cause : structure formelle \zh{因为...所以...}. « Parce que la vie en ville est pratique » = \zh{因为城市的生活很方便} ; conséquence = \zh{所以越来越多的年轻人离开农村}.
Phrase complète (ordre chinois : Cause $\rightarrow$ Conséquence) : 
\zh{因为城市的生活很方便，所以越来越多的年轻人离开农村。} 
\emph{Yīnwèi chéngshì de shēnghuó hěn fāngbiàn, suǒyǐ yuè lái yuè duō de niánqīngrén líkāi nóngcūn.}
\emph{Note :} L'ordre « Conséquence, parce que Cause » (\zh{越来越多的年轻人离开农村，因为城市的生活很方便。}) est possible à l'oral, mais l'ordre « Parce que..., donc... » est la norme de l'écrit et de l'exposé structuré.
\end{exoresolu}

\begin{exercice}[Entraînement écrit : vocabulaire et structures]
1) Traduisez en chinois : « Il quitte son village pour chercher du travail à Douala. »
2) Complétez avec \zh{越来越} : 城市的人口\_\_\_\_\_\_多。
3) Construisez une phrase avec \zh{因为...所以...} sur le thème de l'exode rural au Cameroun (utilisez \zh{雅温得} ou \zh{杜阿拉}).
4) Mettez les mots dans l'ordre : 离开 / 农村 / 很多 / 年轻人 / 去 / 城市 / 因为 / 机会 / 多 / 所以 / 。 (Deux phrases possibles : ordre français ou ordre chinois formel).
\end{exercice}

\begin{corrige}[Exercice « Entraînement écrit »]
1) \zh{他离开他的村子去杜阿拉找工作。} \emph{Tā líkāi tā de cūnzi qù Dù'ālā zhǎo gōngzuò.} (Précision : \zh{村子} \emph{cūnzi} pour « son village » précis, pas \zh{农村} qui est général).
2) \zh{城市的人口越来越多。} \emph{Chéngshì de rénkǒu yuè lái yuè duō.}
3) Exemple attendu : \zh{因为农村的工作很少，所以很多年轻人去雅温得。} \emph{Yīnwèi nóngcūn de gōngzuò hěn shǎo, suǒyǐ hěn duō niánqīngrén qù Yǎwēndé.}
4) Ordre formel (exposé) : \zh{因为机会多，所以很多年轻人离开农村去城市。} \emph{Yīnwèi jīhuì duō, suǒyǐ hěn duō niánqīngrén líkāi nóngcūn qù chéngshì.} \\
Ordre oral (sujet + verbe + parce que) : \zh{很多年轻人离开农村去城市，因为机会多。} \emph{Hěn duō niánqīngrén líkāi nóngcūn qù chéngshì, yīnwèi jīhuì duō.}
\end{corrige}

\courssub{Jeu de rôle et débat guidé : mise en situation}

\begin{experience}[Débat : « Faut-il encourager les jeunes à rester au village ? » (20 min)]
\textbf{Consigne :} Par groupes de 4. Deux élèves sont « Pour le départ » (rôle A), deux sont « Pour le retour / le maintien » (rôle B). Utilisez \textbf{obligatoirement} 3 formules du tableau « Formules fonctionnelles » et \textbf{2 structures grammaticales} (\zh{因为...所以...}, \zh{越来越}, \zh{离开}, \zh{比}).

\textbf{Rôle A (Partir) :} Arguments : travail, argent, études, hôpital, \zh{方便}.
\textbf{Rôle B (Rester/Revenir) :} Arguments : \zh{贵}, \zh{挤}, \zh{想家}, \zh{安静}, famille, terrain, agriculture.

\textbf{Déroulement :}
1. \textbf{Préparation (5 min) :} Chaque duo écrit 3 arguments en chinois (caractères + pinyin) sur une fiche.
2. \textbf{Débat (10 min) :} 
   - A1 ouvre : \zh{大家好，我觉得离开农村去城市好，因为……} 
   - B1 répond : \zh{我不同意。我觉得……因为……所以……} 
   - A2 relance : \zh{可是……一方面……另一方面……你呢？} 
   - B2 conclut : \zh{总之……谢谢大家。}
3. \textbf{Restitution (5 min) :} Un groupe volontaire joue la scène devant la classe. La classe note sur une grille : « Tons corrects ? / Formules utilisées ? / Structures utilisées ? / Fluidité ? ».
\end{experience}

\courssub{Préparation de l'exposé évalué (évaluation formative)}

\begin{methode}[Comment préparer votre exposé oral noté (2-3 minutes)]
\textbf{Étape 1 : Choisir un angle précis.} Ne dites pas tout. Exemples : « Les femmes et l'exode rural », « L'exode rural à Bafoussam vs Yaoundé », « Le retour au village pour créer une entreprise ».
\textbf{Étape 2 : Structurer en 4 parties (utilisez les connecteurs).}
\begin{enumerate}
    \item \textbf{Introduction (30\%)} : \zh{大家好，今天我的题目是……} + \zh{首先，我想说……}
    \item \textbf{Causes (25\%)} : \zh{因为……所以……} + \zh{越来越多……}
    \item \textbf{Avantages / Inconvénients (25\%)} : \zh{一方面……另一方面……} + \zh{比} (comparatif).
    \item \textbf{Conclusion / Avis personnel (20\%)} : \zh{总之，我认为……} + \zh{谢谢大家。}
\end{enumerate}
\textbf{Étape 3 : Rédiger le script (caractères + pinyin).} Vérifiez : ordre des mots, \zh{的/得/地}, tons (surlignez les 3e tons qui changent).
\textbf{Étape 4 : S'entraîner à voix haute.} Minutez-vous. Enregistrez-vous (téléphone). Écoutez : les tons sont-ils stables ? Le débit est-il naturel ? Les formules de liaison (\zh{首先}, \zh{其次}, \zh{最后}) sont-elles fluides ?
\textbf{Étape 5 : Fiche de secours (autorisée).} Une carte Bristol format A6 avec : \textbf{Mots-clés en caractères} + \textbf{Connecteurs} + \textbf{2 phrases difficiles en pinyin}. \textbf{Interdit :} phrase complète en caractères ou en français.
\end{methode}

\begin{savaistu}[Le plus grand exode rural de l'histoire : Chine vs Cameroun]
En Chine, plus de 290 millions de \zh{农民工} travaillent en ville, loin de leur village : c'est le plus grand mouvement de population jamais observé. Chaque année, au moment du Nouvel An chinois (\zh{春节} \emph{Chūnjié}), des centaines de millions de migrants prennent le train pour retrouver leur famille : on parle de \zh{春运} \emph{chūnyùn}, « le transport de printemps », le plus grand déplacement humain annuel du monde. Au Cameroun, le phénomène est tout aussi visible : Yaoundé et Douala attirent chaque année des milliers de jeunes des régions rurales (Ouest, Nord, Adamaoua), à la recherche d'un emploi ou d'études. En chinois, vous savez maintenant le dire : \zh{越来越多的年轻人离开农村去城市。} \emph{Yuè lái yuè duō de niánqīngrén líkāi nóngcūn qù chéngshì.}
\end{savaistu}

\coursec{Vocabulaire clé : ville, campagne, exode rural}

\begin{textebox}[Contexte : reportage radio à Yaoundé]
记者：现在越来越多的喀麦隆年轻人离开农村去城市。\\
Jìzhě: Xiànzài yuè lái yuè duō de Kāmàilóng niánqīngrén líkāi nóngcūn qù chéngshì.\\
« Journaliste : De plus en plus de jeunes Camerounais quittent la campagne pour la ville. »\\[4pt]
农民：在村里找工作太难了，城市的生活条件好。\\
Nóngmín: Zài cūn lǐ zhǎo gōngzuò tài nán le, chéngshì de shēnghuó tiáojiàn hǎo.\\
« Paysan : Au village, trouver du travail est trop dur, les conditions de vie en ville sont bonnes. »\\[4pt]
老师：这叫农村人口外流。村里剩下很多留守儿童和留守老人。\\
Lǎoshī: Zhè jiào nóngcūn rénkǒu wàiliú. Cūn lǐ shèng xià hěn duō liúshǒu értóng hé liúshǒu lǎorén.\\
« Professeur : Cela s'appelle l'exode rural. Il reste au village beaucoup d'enfants et de personnes âgées "restés sur place". »
\end{textebox}

\begin{definition}[Mots clés du thème]
\begin{itemize}
\item \zh{城市} (chéngshì) : ville, zone urbaine.
\item \zh{农村} (nóngcūn) : campagne, zone rurale (litt. « village agricole »).
\item \zh{农村人口外流} (nóngcūn rénkǒu wàiliú) : exode rural (litt. « sortie de la population rurale »).
\item \zh{找工作} (zhǎo gōngzuò) : chercher du travail / un emploi.
\item \zh{生活条件} (shēnghuó tiáojiàn) : conditions de vie.
\item \zh{留守儿童} (liúshǒu értóng) : enfants « restés sur place » (parents partis travailler en ville).
\item \zh{留守老人} (liúshǒu lǎorén) : personnes âgées « restées sur place ».
\item \zh{城市化} (chéngshìhuà) : urbanisation.
\item \zh{发展} (fāzhǎn) : développer, développement.
\item \zh{差距} (chājù) : écart, disparité (ex: \zh{城乡差距} écart ville-campagne).
\end{itemize}
\end{definition}

\begin{center}
\begin{tabularx}{\linewidth}{|X|X|X|X|}
\hline
\rowcolor{popblueL} Caractères & Pinyin & Nature & Traduction \\
\hline
城市 & chéngshì & nom & ville \\
\hline
农村 & nóngcūn & nom & campagne \\
\hline
农村人口外流 & nóngcūn rénkǒu wàiliú & nom composé & exode rural \\
\hline
找工作 & zhǎo gōngzuò & verbe + objet & chercher du travail \\
\hline
生活条件 & shēnghuó tiáojiàn & nom & conditions de vie \\
\hline
留守儿童 & liúshǒu értóng & nom composé & enfants laissés au village \\
\hline
留守老人 & liúshǒu lǎorén & nom composé & personnes âgées laissées au village \\
\hline
城市化 & chéngshìhuà & nom / verbe & urbanisation / urbaniser \\
\hline
发展 & fāzhǎn & verbe / nom & développer / développement \\
\hline
差距 & chājù & nom & écart, disparité \\
\hline
\end{tabularx}
\end{center}

\begin{attention}[Pièges de prononciation et d'écriture]
\begin{itemize}
\item \zh{农} (nóng, 2\textsuperscript{e} ton) $\neq$ \zh{浓} (nóng, épais/concentré) : même son, clé différente (\zh{辶} marche vs \zh{氵}eau).
\item \zh{外流} (wàiliú) : \zh{外} (wài, dehors) 4\textsuperscript{e} ton ; \zh{流} (liú, couler/partir) 2\textsuperscript{e} ton. Ne pas confondre avec \zh{外留} (wàiliú, « retenir à l'extérieur », rare).
\item \zh{留守} (liúshǒu) : \zh{留} (liú, rester) 2\textsuperscript{e} ton ; \zh{守} (shǒu, garder) 3\textsuperscript{e} ton. \zh{留守儿童} est un terme sociologique précis, ne pas dire \zh{剩下的孩子} (enfants restants) dans un exposé formel.
\item \zh{城市化} : suffixe \zh{化} (huà, -isation) : \zh{现代化} (modernisation), \zh{工业化} (industrialisation).
\end{itemize}
\end{attention}

\coursec{Formules fonctionnelles : présenter, opiner, réagir, relancer}

Dans la section précédente, nous avons appris le vocabulaire clé pour parler de la ville, de la campagne et de l'exode rural : \zh{城市} (chéngshì), \zh{农村} (nóngcūn), \zh{农村人口外流} (nóngcūn rénkǒu wàiliú). Mais un bon vocabulaire ne suffit pas : pour faire un exposé ou participer à un débat en chinois, il faut maîtriser les \cle{formules fonctionnelles}, c'est-à-dire les phrases toutes faites qui structurent la prise de parole.

\courssub{Observer d'abord : un mini-exposé modèle}

\begin{textebox}[Un élève de Première présente l'exode rural]
大家好！今天我要谈一谈喀麦隆的农村人口外流问题。\\
Dàjiā hǎo! Jīntiān wǒ yào tántan Kāmàilóng de nóngcūn rénkǒu wàiliú wèntí.\\
« Bonjour à tous ! Aujourd'hui, je vais parler du problème de l'exode rural au Cameroun. »\\[4pt]
首先，农村的工作很少；然后，城市的学校比较好；最后，年轻人想要新的生活。\\
Shǒuxiān, nóngcūn de gōngzuò hěn shǎo; ránhòu, chéngshì de xuéxiào bǐjiào hǎo; zuìhòu, niánqīngrén xiǎng yào xīn de shēnghuó.\\
« D'abord, il y a peu de travail à la campagne ; ensuite, les écoles en ville sont meilleures ; enfin, les jeunes veulent une nouvelle vie. »\\[4pt]
我觉得政府应该帮助农村。总之，城市和农村都很重要。我的讲话完了，谢谢大家！\\
Wǒ juéde zhèngfǔ yīnggāi bāngzhù nóngcūn. Zǒngzhī, chéngshì hé nóngcūn dōu hěn zhòngyào. Wǒ de jiǎnghuà wán le, xièxie dàjiā!\\
« Je pense que le gouvernement devrait aider la campagne. Bref, la ville et la campagne sont toutes les deux importantes. Mon exposé est terminé, merci à tous ! »
\end{textebox}

Observez ce texte : il est court, mais il contient \textbf{toutes} les formules d'un exposé réussi. Repérons-les : salutation (\zh{大家好}), annonce du sujet (\zh{今天我要谈一谈...}), plan en trois temps (\zh{首先... 然后... 最后...}), avis personnel (\zh{我觉得...}), justification (\zh{因为...所以...}), conclusion (\zh{总之...}) et remerciement (\zh{谢谢大家}).

\courssub{Présenter un exposé et annoncer le plan}

\begin{propriete}[Saluer et annoncer le sujet]
Structure : \zh{大家好！今天我要谈一谈 + sujet。}\\
\emph{Dàjiā hǎo! Jīntiān wǒ yào tántan + sujet.} « Bonjour à tous ! Aujourd'hui, je vais parler de... »
\begin{itemize}
\item \zh{大家} (dàjiā) = « tout le monde » ; \zh{大家好} est la salutation standard devant un public.
\item \zh{要} (yào) exprime l'intention proche : « je vais ».
\item \zh{谈一谈} (tántan) : forme redoublée de \zh{谈} (parler de) ; le redoublement (verbe + 一 + verbe) adoucit l'action : « parler un peu de ». Variante simple : \zh{说一说} (shuōshuo).
\end{itemize}
Exemple : \zh{今天我要谈一谈农村的生活。} \emph{Jīntiān wǒ yào tántan nóngcūn de shēnghuó.} « Aujourd'hui, je vais parler de la vie à la campagne. »
\end{propriete}

\begin{propriete}[Annoncer le plan : 首先... 然后... 最后...]
Structure : \zh{首先 + point 1；然后 + point 2；最后 + point 3。}\\
\zh{首先} (shǒuxiān) = d'abord ; \zh{然后} (ránhòu) = ensuite ; \zh{最后} (zuìhòu) = enfin. Ces trois adverbes se placent \textbf{en tête de phrase}, suivis d'une virgule chinoise (，). On peut répéter \zh{然后} pour plus de points : \zh{首先... 然后... 然后... 最后...}.
\end{propriete}

\begin{exemplebox}[Annoncer un plan]
\zh{首先，很多年轻人离开农村；然后，他们去了大城市；最后，农村只剩下老人和孩子。}\\
\emph{Shǒuxiān, hěn duō niánqīngrén líkāi nóngcūn; ránhòu, tāmen qù le dà chéngshì; zuìhòu, nóngcūn zhǐ shèng xià lǎorén hé háizi.}\\
« D'abord, beaucoup de jeunes quittent la campagne ; ensuite, ils vont dans les grandes villes ; enfin, il ne reste à la campagne que les personnes âgées et les enfants. »
\end{exemplebox}

\courssub{Donner son avis et justifier}

\begin{propriete}[Trois façons de donner son avis]
Structure commune : \textbf{formule d'opinion + phrase complète} (sujet + verbe). Pas de « que » en chinois.
\begin{itemize}
\item \zh{我觉得...} (wǒ juéde...) : « je trouve que... » — la plus courante à l'oral, avis personnel, subjectif.
\item \zh{我认为...} (wǒ rènwéi...) : « je pense / j'estime que... » — plus formel, avis réfléchi, argumenté.
\item \zh{在我看来，...} (zài wǒ kànlái, ...) : « à mon avis, ... » — introduit un point de vue global ; \textbf{virgule obligatoire} après \zh{看来}.
\end{itemize}
\end{propriete}

\begin{exemplebox}[Donner son avis]
\zh{我觉得城市的生活比较方便。}\\
\emph{Wǒ juéde chéngshì de shēnghuó bǐjiào fāngbiàn.} « Je trouve que la vie en ville est plus pratique. »\\[4pt]
\zh{我认为年轻人应该学习农业。}\\
\emph{Wǒ rènwéi niánqīngrén yīnggāi xuéxí nóngyè.} « Je pense que les jeunes devraient étudier l'agriculture. »\\[4pt]
\zh{在我看来，农村的空气比城市的空气好。}\\
\emph{Zài wǒ kànlái, nóngcūn de kōngqì bǐ chéngshì de kōngqì hǎo.} « À mon avis, l'air de la campagne est meilleur que celui de la ville. »
\end{exemplebox}

\begin{propriete}[Justifier : 因为...所以... et 由于...]
Structure : \zh{因为 + cause，所以 + conséquence。}\\
\emph{Yīnwèi... suǒyǐ...} « parce que... donc... ».
\textbf{Différence majeure avec le français} : en chinois, les deux mots s'emploient \textbf{souvent ensemble} dans la même phrase. C'est la norme. On peut utiliser seul \zh{因为} (cause) ou seul \zh{所以} (conséquence), mais la structure complète est préférée à l'oral pour la clarté.
Variante formelle (écrit / discours officiel) : \zh{由于 + cause，+ conséquence} (yóuyú, « en raison de »), souvent sans \zh{所以}.
\end{propriete}

\begin{exemplebox}[Justifier]
\zh{因为农村的工作很少，所以很多年轻人去城市。}\\
\emph{Yīnwèi nóngcūn de gōngzuò hěn shǎo, suǒyǐ hěn duō niánqīngrén qù chéngshì.}\\
« Parce qu'il y a peu de travail à la campagne, beaucoup de jeunes vont en ville. »\\[4pt]
\zh{由于交通不方便，农村发展慢。}\\
\emph{Yóuyú jiāotōng bù fāngbiàn, nóngcūn fāzhǎn màn.}\\
« En raison des transports peu pratiques, la campagne se développe lentement. »
\end{exemplebox}

\begin{attention}[因为...所以... : le piège du français]
En français, « parce que... donc... » dans la même phrase est une faute de syntaxe (pléonasme). En chinois, c'est la structure \textbf{standard et recommandée} ! Ne traduisez jamais mot à mot : gardez \zh{因为...所以...} en chinois, mais n'en traduisez qu'un seul en français (généralement « parce que » ou « c'est pourquoi »). Erreur typique : oublier \zh{所以} $\rightarrow$ phrase inachevée pour un Chinois.
\end{attention}

\courssub{Réagir : être d'accord ou pas, poliment}

\begin{exemplebox}[Être d'accord]
\zh{我同意（你的看法）。} \emph{Wǒ tóngyì (nǐ de kànfǎ).} « Je suis d'accord (avec ton point de vue). »\\[4pt]
\zh{你说得对！} \emph{Nǐ shuō de duì!} « Tu as raison ! » (litt. « ce que tu dis est juste » ; structure V + \zh{得} + adjectif).
\end{exemplebox}

\begin{exemplebox}[Ne pas être d'accord poliment]
\zh{我不同意你的看法。}\\
\emph{Wǒ bù tóngyì nǐ de kànfǎ.} « Je ne suis pas d'accord avec ton point de vue. »\\[4pt]
\zh{我不这么认为。} \emph{Wǒ bù zhème rènwéi.} « Je ne le pense pas ainsi. » (doux : \zh{这么} = « ainsi / de cette façon »)\\[4pt]
\zh{可是，农村也有好处。} \emph{Kěshì, nóngcūn yě yǒu hǎochù.} « Mais la campagne a aussi des avantages. » (\zh{可是} = mais, introduit une objection).
\end{exemplebox}

\begin{attention}[La politesse du désaccord \& Sandhi de \zh{不}]
En Chine comme au Cameroun, contredire frontalement est mal perçu. Préférez \zh{我不这么认为} ou \zh{你说得对，可是...} à \zh{你错了 !}.
\textbf{Prononciation :} \zh{不} (bù, 4\textsuperscript{e} ton) devient \textbf{2\textsuperscript{e} ton (bú)} devant un 4\textsuperscript{e} ton. \zh{不同意} $\rightarrow$ \emph{Wǒ \textbf{bú} tóngyì}. \zh{不对} $\rightarrow$ \emph{\textbf{bú} duì}. C'est obligatoire à l'oral.
\end{attention}

\courssub{Relancer le débat et conclure}

\begin{propriete}[呢 : la particule qui relance]
Structure : \textbf{nom/pronom + 呢？}\\
\zh{呢} (ne, ton neutre) « renvoie » la question à l'interlocuteur sans répéter le verbe. \zh{你呢？} (Nǐ ne?) = « Et toi ? ». \zh{农村呢？} « Et la campagne ? ».
Autres relances : \zh{你觉得怎么样？} (Nǐ juéde zěnmeyàng?) « Qu'en penses-tu ? » ; \zh{你为什么这么想？} (Nǐ wèishénme zhème xiǎng?) « Pourquoi penses-tu cela ? »
\end{propriete}

\begin{propriete}[Conclure un exposé]
\zh{总之，...} (zǒngzhī) = « bref, en résumé » : en tête de phrase + virgule.
Formule de clôture rituelle : \zh{我的讲话完了，谢谢大家！} \emph{Wǒ de jiǎnghuà wán le, xièxie dàjiā!} « Mon exposé est terminé, merci à tous ! » \zh{了} marque l'accomplissement (\zh{完了} = « est fini »).
\end{propriete}

\begin{methode}[Un bon exposé oral en 5 étapes]
\begin{enumerate}
\item \textbf{Saluer} : \zh{大家好！} (regardez le public, tons : dà-jiā-hǎo, 4-1-3).
\item \textbf{Annoncer le sujet} : \zh{今天我要谈一谈...} (lentement, c'est la phrase clé).
\item \textbf{Développer} avec \zh{首先... 然后... 最后...} : une idée par étape, phrases courtes.
\item \textbf{Donner son avis} : \zh{我觉得...} / \zh{在我看来，...}, justifié par \zh{因为...所以...}.
\item \textbf{Conclure et remercier} : \zh{总之... 我的讲话完了，谢谢大家！}
\end{enumerate}
\end{methode}

\begin{popfigure}
\begin{tikzpicture}[node distance=0.55cm, every node/.style={font=\small, align=center}]
\node[draw=popblue, fill=popblueL, rounded corners, minimum width=4.8cm, align=center] (a){\zh{大家好！}\\ \textbf{Saluer}};
\node[draw=poppurple, fill=poppurpleL, rounded corners, minimum width=4.8cm, below=of a, align=center] (b){\zh{今天我要谈一谈...}\\ \textbf{Annoncer le sujet}};
\node[draw=poporange, fill=poporangeL, rounded corners, minimum width=4.8cm, below=of b, align=center] (c){\zh{首先... 然后... 最后...}\\ \textbf{Développer}};
\node[draw=popgreen, fill=popgreenL, rounded corners, minimum width=4.8cm, below=of c, align=center] (d){\zh{总之...}\\ \textbf{Conclure}};
\node[draw=poppink, fill=poppinkL, rounded corners, minimum width=4.8cm, below=of d, align=center] (e){\zh{谢谢大家！}\\ \textbf{Remercier}};
\draw[-{Stealth}, thick, popdark] (a) -- (b);
\draw[-{Stealth}, thick, popdark] (b) -- (c);
\draw[-{Stealth}, thick, popdark] (c) -- (d);
\draw[-{Stealth}, thick, popdark] (d) -- (e);
\end{tikzpicture}
\legende{Organigramme : structure d'un exposé oral en chinois.}
\end{popfigure}

\begin{center}
\begin{tabularx}{\linewidth}{|X|X|X|X|}
\hline
\rowcolor{popblueL} Caractères & Pinyin & Nature & Traduction \\
\hline
大家好！ & Dàjiā hǎo! & locution (interjection) & Bonjour à tous ! \\
\hline
今天我要谈一谈... & Jīntiān wǒ yào tántan... & structure verbale & Aujourd'hui je vais parler de... \\
\hline
首先 / 然后 / 最后 & shǒuxiān / ránhòu / zuìhòu & adverbes de séquence & d'abord / ensuite / enfin \\
\hline
我觉得... & Wǒ juéde... & verbe mental + clause & je trouve que... \\
\hline
我认为... & Wǒ rènwéi... & verbe mental + clause & je pense que... (formel) \\
\hline
在我看来，... & Zài wǒ kànlái,... & préposition + nom + particule & à mon avis,... \\
\hline
因为...所以... & yīnwèi... suǒyǐ... & conjonctions corrélatives & parce que... (donc)... \\
\hline
由于... & yóuyú... & préposition (cause) & en raison de... \\
\hline
我同意... & Wǒ tóngyì... & verbe + objet & je suis d'accord avec... \\
\hline
你说得对！ & Nǐ shuō de duì! & phrase (V + 得 + Adj) & Tu as raison ! \\
\hline
我不这么认为。 & Wǒ bù zhème rènwéi. & phrase (négation + adv + V) & Je ne le pense pas ainsi. \\
\hline
可是... & kěshì... & conjonction (adversative) & mais... \\
\hline
你呢？ & Nǐ ne? & pronom + particule & et toi ? \\
\hline
你觉得怎么样？ & Nǐ juéde zěnmeyàng? & phrase interrogative & Qu'en penses-tu ? \\
\hline
总之，... & Zǒngzhī,... & adverbe de conclusion & bref,... / en résumé,... \\
\hline
我的讲话完了，谢谢大家！ & Wǒ de jiǎnghuà wán le, xièxie dàjiā! & phrase rituelle & Mon exposé est fini, merci ! \\
\hline
\end{tabularx}
\end{center}

\courssub{Travail sur les tons et la prononciation}

\begin{methode}[Exercices de tons ciblés pour le débat]
\begin{enumerate}
\item \textbf{Tons isolés (mots-clés) :} Répétez en exagérant la courbe.
\begin{itemize}
\item \zh{农村} (nóng-cūn) : 2-1 (montant - haut).
\item \zh{城市} (chéng-shì) : 2-4 (montant - descendant).
\item \zh{外流} (wài-liú) : 4-2 (descendant - montant).
\item \zh{首先} (shǒu-xiān) : 3-1 (bas - haut). \textit{Attention : 3\textsuperscript{e} ton bas (demi-ton) avant 1\textsuperscript{e} ton.}
\item \zh{同意} (tóng-yì) : 2-4.
\item \zh{不同意} (bú-tóng-yì) : \textbf{2-2-4} (sandhi \zh{不} $\rightarrow$ \zh{ú} devant 4\textsuperscript{e} ton).
\item \zh{觉得} (jué-de) : 2-neutre. \zh{认为} (rèn-wéi) : 4-2.
\item \zh{因为} (yīn-wèi) : 1-4. \zh{所以} (suǒ-yǐ) : 3-3 (souvent 2-3 à l'oral).
\end{itemize}
\item \textbf{Chaînes tonales (phrases types) :} Travaillez le rythme (groupes de sens).
\begin{itemize}
\item \zh{今天我要谈一谈 / 农村人口外流 / 问题。} (4-1-3 | 2-2-4-2 | 4-2)
\item \zh{首先 / 农村的工作很少。} (3-1 | 2-1-4-3-3)
\item \zh{因为 / 城市的工作多 / 所以 / 生活方便。} (1-4 | 2-4-4-1 | 3-3 | 1-4)
\item \zh{我不同意 / 你的看法。} (3-2-2-4 | 3-4-3)
\item \zh{总之 / 城市和农村 / 都很重要。} (3-1 | 2-4-2-1 | 1-3-4)
\end{itemize}
\item \textbf{Intonation finale :} Question ouverte (\zh{你呢？} $\nearrow$), Affirmation (\zh{我觉得好。} $\searrow$), Exclamation (\zh{你说得对！} $\searrow$ fort).
\end{enumerate}
\end{methode}

\begin{experience}[Atelier prononciation : "Le bon ton change le sens"]
En binôme. L'élève A lit une phrase avec un ton volontairement faux (ex: \zh{农村} \emph{nǒng cǔn} 3-3 au lieu de 2-1). L'élève B doit identifier l'erreur et la corriger. Échangez les rôles.
\begin{center}
\begin{tabularx}{\linewidth}{|X|X|}
\hline
\rowcolor{poporangeL} Phrase à corriger & Correction attendue \\
\hline
\zh{我要谈一谈 农村 (nǒng cǔn) 问题。} & \zh{农村 (nóng cūn)} 2-1 \\
\hline
\zh{我 不 (bù) 同意。} & \zh{不 (bú)} 2\textsuperscript{e} ton (devant 4\textsuperscript{e}) \\
\hline
\zh{首先 (shǒu xiǎn) 3-3...} & \zh{首先 (shǒu xiān)} 3-1 \\
\hline
\zh{城市 (chéng shǐ) 2-3...} & \zh{城市 (chéng shì)} 2-4 \\
\hline
\end{tabularx}
\end{center}
\end{experience}

\courssub{Jeux de rôle et débat guidé}

\begin{experience}[Jeu de rôle : "Le conseil municipal de Buea"]
\textbf{Contexte :} Le maire de Buea organise une réunion publique sur l'exode rural. 4 rôles.
\begin{enumerate}
\item \textbf{Le Maire} : Ouvre la séance, annonce le sujet (\zh{今天我们讨论...}), invite les avis.
\item \textbf{Un jeune parti à Douala} : Explique pourquoi il est parti (\zh{因为...所以...}), donne son avis (\zh{我觉得城市好}).
\item \textbf{Un paysan resté au village} : Désaccord poli (\zh{我不这么认为}), arguments pour la campagne (\zh{空气好，成本低}), relance (\zh{你呢？}).
\item \textbf{Un enseignant} : Synthétise (\zh{总之...}), propose une solution (\zh{政府应该...}), clôture.
\end{enumerate}
\textbf{Contraintes linguistiques obligatoires :} Chaque intervenant doit utiliser au moins 3 formules du tableau récapitulatif. Notez les tons \zh{不} (bú/bù) et \zh{一} (yī/yí/yì).
\end{experience}

\begin{experience}[Débat guidé : "Faut-il freiner l'exode rural ?"]
\textbf{Phase 1 (Préparation 5 min) :} En groupes de 4. Listez 3 arguments POUR freiner (ex: \zh{保护农业}, \zh{减少城市压力}) et 3 arguments CONTRE (ex: \zh{自由迁徙}, \zh{城市发展快}). Écrivez les mots-clés en caractères.
\textbf{Phase 2 (Débat 10 min) :} Tour de table. Structure imposée :
\zh{我觉得... 因为... 所以... 你呢？}
L'enseignant note : justesse des tons, fluidité \zh{因为...所以}, politesse désaccord (\zh{可是/我不这么认为}), usage de \zh{呢}.
\textbf{Phase 3 (Bilan 5 min) :} Un rapporteur par groupe résume en 30 sec : \zh{总之，我们组认为...}
\end{experience}

\courssub{Bilan et préparation de l'exposé évalué}

\begin{aretenir}[Check-list pour l'exposé évalué (2-3 minutes)]
\begin{itemize}
\item[\ding{51}] \textbf{Structure} : Salutation $\rightarrow$ Sujet ($\zh{今天我要谈一谈...}$) $\rightarrow$ Plan ($\zh{首先/然后/最后}$) $\rightarrow$ Avis + Justification ($\zh{我觉得...因为...所以...}$) $\rightarrow$ Conclusion ($\zh{总之...}$) $\rightarrow$ Remerciement.
\item[\ding{51}] \textbf{Vocabulaire} : Minimum 6 mots de la liste (ex: \zh{农村人口外流}, \zh{留守儿童}, \zh{城市化}, \zh{差距}, \zh{发展}, \zh{生活条件}).
\item[\ding{51}] \textbf{Grammaire} : Au moins une justification \zh{因为...所以...} correcte ; un désaccord ou une nuance (\zh{可是/虽然...但是...}).
\item[\ding{51}] \textbf{Prononciation} : Tons des mots-clés stables (\zh{农村, 城市, 首先, 同意, 因为, 所以}) ; sandhi \zh{不} (bú) devant 4\textsuperscript{e} ton ; \zh{一} (yí/yi) selon contexte ; intonation finale descendante sur les affirmations.
\item[\ding{51}] \textbf{Présence} : Regard vers le public, débit maîtrisé, pas de lecture mot à mot (fiches autorisées avec mots-clés seulement).
\end{itemize}
\end{aretenir}

\begin{savaistu}[Exode rural : regards croisés Cameroun - Chine]
\begin{itemize}
\item \textbf{Cameroun :} Taux d'urbanisation $\approx$ 58\% (2023). Principaux pôles : Yaoundé (politique), Douala (économique), Buea, Bamenda. Causes : manque d'infrastructures rurales (routes, électricité, santé), attractivité du secteur informel urbain, insécurité dans certaines zones (NosO). Conséquences : bidonvilles (\zh{棚户区} pénghùqū) à Douala/Yaoundé, vieillissement des villages, \zh{留守儿童} scolarisés par grands-parents.
\item \textbf{Chine :} Taux d'urbanisation $\approx$ 66\% (2023). Phénomène massif depuis les années 1980 : \zh{农民工} (nóngmíngōng, travailleurs migrants ruraux), \zh{留守儿童} (61 millions en 2020), \zh{留守老人}. Système \zh{户口} (hùkǒu, registre de résidence) : frein historique à l'installation permanente en ville, réformes en cours. Politique actuelle : \zh{新型城镇化} (nouvelle urbanisation) + \zh{乡村振兴} (revitalisation rurale) pour réduire l'écart \zh{城乡差距}.
\item \textbf{Comparaison :} Les deux pays font face au défi de l'équilibre territorial. La Chine a une planification étatique forte (\zh{五年规划} plans quinquennaux) ; le Cameroun mise sur la décentralisation et l'agriculture de seconde génération. Le vocabulaire chinois (\zh{留守儿童}, \zh{空心村} « village creux/vide ») est désormais utilisé dans les rapports ONG au Cameroun.
\end{itemize}
\end{savaistu}

\begin{exoresolu}[Compléter un débat avec les formules fonctionnelles]
Énoncé : Complétez ce dialogue entre deux élèves de Yaoundé avec les formules étudiées (\zh{我觉得}, \zh{因为...所以}, \zh{我不同意}, \zh{你呢}, \zh{总之}).\\
A : \zh{城市好还是农村好？\_\_\_\_\_\_ 城市好。}\\
B : \zh{为什么？}\\
A : \zh{\_\_\_\_\_\_ 城市的工作多，\_\_\_\_\_\_ 生活方便。}\\
B : \zh{\_\_\_\_\_\_ 你的看法。农村的空气好，东西也便宜。\_\_\_\_\_\_ ？}\\
A : \zh{我觉得两个都好。}\\
B : \zh{\_\_\_\_\_\_ ，城市和农村都有自己的好处。}
\tcblower
Solution détaillée :\\
A : \zh{城市好还是农村好？\textbf{我觉得}城市好。} — Avis personnel avec \zh{我觉得}.\\
A : \zh{\textbf{因为}城市的工作多，\textbf{所以}生活方便。} — Justification complète \zh{因为...所以...} (norme chinoise).\\
B : \zh{\textbf{我不同意}你的看法。农村的空气好，东西也便宜。\textbf{你呢}？} — Désaccord direct mais courtois entre pairs, relance avec \zh{呢}.\\
B : \zh{\textbf{总之}，城市和农村都有自己的好处。} — Conclusion synthétique avec \zh{总之}.\\
Traduction : « Ville ou campagne, laquelle est mieux ? — Je trouve la ville mieux. — Pourquoi ? — Parce qu'il y a plus d'emplois en ville, donc la vie est pratique. — Je ne suis pas d'accord avec ton avis. L'air est bon à la campagne et les choses sont moins chères. Et toi ? — Je trouve les deux bien. — Bref, ville et campagne ont chacune leurs avantages. »
\end{exoresolu}

\begin{attention}[Erreurs fréquentes des francophones — Récapitulatif]
\begin{itemize}
\item \textbf{Verbe \zh{觉得} + adjectif} : \zh{我觉得城市方便} (sans \zh{是}). \zh{是} ne s'emploie que devant un nom (\zh{我觉得城市是个好地方}).
\item \textbf{\zh{谈} vs \zh{关于}} : \zh{谈一谈这个问题} (verbe transitif direct). Pas de \zh{关于} après \zh{谈}. Si \zh{关于}, utiliser \zh{说一说关于...的问题}.
\item \textbf{Tons critiques} : \zh{首先} (shǒuxiān 3-1), \zh{同意} (tóngyì 2-4), \zh{不同意} (bú tóngyì 2-2-4), \zh{城市} (chéngshì 2-4), \zh{农村} (nóngcūn 2-1).
\item \textbf{\zh{你呢} seul} : Ne pas ajouter \zh{怎么样} (\zh{你呢怎么样？} $\times$).
\item \textbf{Ordre \zh{因为...所以...}} : Ne pas inverser (Cause $\rightarrow$ Conséquence). Ne pas omettre \zh{所以} à l'oral formel.
\end{itemize}
\end{attention}

\begin{exercice}[À vous de parler — Expression orale évaluée]
En binôme, préparez un mini-débat de 6 à 8 répliques sur le thème : \zh{年轻人应该留在农村还是去城市？} (Les jeunes devraient-ils rester à la campagne ou partir en ville ?).
\textbf{Contraintes obligatoires :}
\begin{enumerate}
\item Annonce de sujet claire (\zh{今天我们讨论...}).
\item Deux avis distincts (\zh{我觉得} / \zh{我认为} / \zh{在我看来}).
\item Une justification complète \zh{因为...所以...}.
\item Un désaccord poli (\zh{我不这么认为} ou \zh{可是...}).
\item Une relance avec \zh{呢} ou \zh{你觉得怎么样？}.
\item Conclusion commune avec \zh{总之}.
\end{enumerate}
Enregistrez-vous (audio/vidéo) et auto-évaluez avec la check-list de la boîte « À retenir ».
\end{exercice}

Ces formules sont désormais vos outils : dans le prochain module, nous les réinvestirons pour l'expression écrite (rédaction d'un article d'opinion) et la compréhension de l'écrit (article de presse chinois sur l'\zh{乡村振兴}).

\coursec{Dialogues modèles avec pinyin}

Dans la section précédente, nous avons appris les \cle{formules fonctionnelles} pour présenter un sujet (\zh{我今天要谈……}), donner son avis (\zh{我觉得……}), réagir (\zh{我同意 / 我不同意}) et relancer la parole (\zh{你呢？你怎么看？}). Il est temps de mettre ces outils en action : nous allons maintenant étudier, phrase par phrase, deux \cle{dialogues modèles} sur l'exode rural. Ces dialogues sont le cœur de l'évaluation orale : vous devrez les connaître, les prononcer avec les bons tons, puis les transformer grâce à la technique de la substitution.

\courssub{Dialogue 1 : une interview à la sortie du marché}

\begin{textebox}[Dialogue 1 — Interview sur l'exode rural]
A：你为什么离开农村？\\
\emph{Nǐ wèishénme líkāi nóngcūn?}\\
« Pourquoi as-tu quitté la campagne ? »\\[4pt]
B：因为农村的工作很少，我想在城市找工作。\\
\emph{Yīnwèi nóngcūn de gōngzuò hěn shǎo, wǒ xiǎng zài chéngshì zhǎo gōngzuò.}\\
« Parce qu'il y a très peu de travail à la campagne, je veux chercher du travail en ville. »\\[4pt]
A：你觉得城市的生活怎么样？\\
\emph{Nǐ juéde chéngshì de shēnghuó zěnmeyàng?}\\
« Comment trouves-tu la vie en ville ? »\\[4pt]
B：很方便，可是也很贵。\\
\emph{Hěn fāngbiàn, kěshì yě hěn guì.}\\
« C'est très pratique, mais c'est aussi très cher. »
\end{textebox}

\textbf{Observation.} Lisez d'abord le dialogue en silence, puis repérez les éléments suivants :

\begin{itemize}
\item La question s'ouvre sur le mot interrogatif \zh{为什么} \emph{wèishénme} « pourquoi », placé \textbf{après le sujet}, jamais en tête de phrase comme en français.
\item La réponse commence par \zh{因为} \emph{yīnwèi} « parce que » : en chinois, on répond presque toujours à « pourquoi » par « parce que », mot à mot.
\item La deuxième question utilise la structure d'opinion \zh{你觉得……怎么样？} « que penses-tu de… ? », vue dans la section précédente.
\item La réponse finale oppose deux idées avec \zh{可是} \emph{kěshì} « mais ».
\end{itemize}

\begin{propriete}[Interroger avec 为什么 et répondre avec 因为]
\textbf{Structure de la question :} Sujet + \cle{为什么} + verbe ?\\
\zh{你为什么离开农村？} \emph{Nǐ wèishénme líkāi nóngcūn?} — « Pourquoi quittes-tu la campagne ? »\\
\textbf{Structure de la réponse :} \cle{因为} + cause, (所以) + conséquence.\\
\zh{因为农村的工作很少，所以我想去城市。} \emph{Yīnwèi nóngcūn de gōngzuò hěn shǎo, suǒyǐ wǒ xiǎng qù chéngshì.} — « Parce qu'il y a peu de travail à la campagne, (donc) je veux aller en ville. »\\
\textbf{Différence majeure avec le français :} en chinois, \zh{因为……所以……} peuvent se combiner dans la même phrase (« parce que… donc… »), ce qui est une faute en français. À l'oral, on peut n'utiliser que \zh{因为}.
\end{propriete}

\begin{attention}[Ne dites pas « wèishénme nǐ… »]
Les francophones transposent l'ordre français « Pourquoi \emph{tu} pars ? » et produisent la faute \zh{为什么你离开农村？}. En chinois, le mot interrogatif reste \textbf{à la place de la réponse}, après le sujet : \zh{你为什么离开农村？}. De même, on dit \zh{你想买什么？} « Que veux-tu acheter ? » et jamais \zh{什么你想买？}.
\end{attention}

\courssub{Dialogue 2 : un petit débat entre camarades}

\begin{textebox}[Dialogue 2 — Débat : rester ou partir ?]
A：我觉得年轻人应该留在农村。\\
\emph{Wǒ juéde niánqīngrén yīnggāi liú zài nóngcūn.}\\
« Je pense que les jeunes devraient rester à la campagne. »\\[4pt]
B：我不同意。城市有更多的机会。\\
\emph{Wǒ bù tóngyì. Chéngshì yǒu gèng duō de jīhuì.}\\
« Je ne suis pas d'accord. La ville offre plus d'opportunités. »\\[4pt]
A：可是农村的空气很好，生活很安静。\\
\emph{Kěshì nóngcūn de kōngqì hěn hǎo, shēnghuó hěn ānjìng.}\\
« Mais l'air de la campagne est bon, et la vie y est calme. »\\[4pt]
B：你说得对，但是农村的学校和医院不太好。\\
\emph{Nǐ shuō de duì, dànshì nóngcūn de xuéxiào hé yīyuàn bù tài hǎo.}\\
« Tu as raison, mais les écoles et les hôpitaux de la campagne ne sont pas très bons. »
\end{textebox}

\textbf{Observation.} Ce dialogue illustre la mécanique complète du débat en quatre temps :

\begin{enumerate}
\item \textbf{Opinion} : \zh{我觉得……} + proposition.
\item \textbf{Désaccord} : \zh{我不同意。} + contre-argument.
\item \textbf{Concession} : \zh{可是……} + nouvel argument.
\item \textbf{Accord partiel puis réplique} : \zh{你说得对，但是……} « tu as raison, mais… ».
\end{enumerate}

\begin{attention}[La place de 应该 yīnggāi]
\cle{应该} \emph{yīnggāi} « devoir, il faudrait » est un verbe auxiliaire : il se place \textbf{avant le verbe principal}, comme 想, 会, 可以.\\
\zh{年轻人应该留在农村。} \emph{Niánqīngrén yīnggāi liú zài nóngcūn.} — « Les jeunes devraient rester à la campagne. »\\
Ne dites jamais \zh{年轻人留应该在农村}. La négation se met devant l'auxiliaire : \zh{我们不应该忘记农村。} \emph{Wǒmen bù yīnggāi wàngjì nóngcūn.} — « Nous ne devrions pas oublier la campagne. »
\end{attention}

\begin{exemplebox}[Deux phrases à réutiliser dans vos débats]
\zh{城市有更多的机会。} \emph{Chéngshì yǒu gèng duō de jīhuì.} — « La ville offre plus d'opportunités. » (structure : lieu + \zh{有} + \zh{更} + adjectif + \zh{的} + nom)\\[4pt]
\zh{农村的空气很新鲜。} \emph{Nóngcūn de kōngqì hěn xīnxiān.} — « L'air de la campagne est frais. » (structure : lieu + \zh{的} + nom + \zh{很} + adjectif)
\end{exemplebox}

\begin{center}
\begin{tabularx}{\linewidth}{|X|X|X|X|}
\hline
\rowcolor{popblueL} Caractères & Pinyin & Nature & Traduction \\
\hline
离开 & líkāi & verbe & quitter \\
\hline
农村 & nóngcūn & nom & la campagne, les villages \\
\hline
城市 & chéngshì & nom & la ville \\
\hline
机会 & jīhuì & nom & opportunité, chance \\
\hline
空气 & kōngqì & nom & l'air \\
\hline
新鲜 & xīnxiān & adjectif & frais \\
\hline
安静 & ānjìng & adjectif & calme, silencieux \\
\hline
方便 & fāngbiàn & adjectif & pratique, commode \\
\hline
贵 & guì & adjectif & cher \\
\hline
应该 & yīnggāi & auxiliaire & devoir, il faudrait \\
\hline
留 & liú & verbe & rester \\
\hline
同意 & tóngyì & verbe & être d'accord \\
\hline
更 & gèng & adverbe & davantage, plus \\
\hline
\end{tabularx}
\end{center}

\courssub{Pour la ville ou pour la campagne ? La banque d'arguments}

La figure suivante classe les arguments des deux dialogues. Apprenez-la : elle vous servira de plan pour votre propre exposé.

\begin{popfigure}
\begin{center}
\begin{tikzpicture}
\node[draw=popdark, fill=popgoldL, rounded corners, font=\bfseries, align=center] (c) at (0,0) {农村还是城市？\\ \emph{Nóngcūn háishì chéngshì?}\\ Campagne ou ville ?};
\node[draw=popblue, fill=popblueL, rounded corners, align=center, text width=5.2cm] (ville) at (-4.6,-3.2) {\textbf{Pour la ville}\\ 工作多 \emph{gōngzuò duō} : beaucoup de travail\\ 机会更多 \emph{jīhuì gèng duō} : plus d'opportunités\\ 生活很方便 \emph{shēnghuó hěn fāngbiàn} : vie pratique\\ 学校和医院好 \emph{xuéxiào hé yīyuàn hǎo} : bonnes écoles et hôpitaux};
\node[draw=popgreen, fill=popgreenL, rounded corners, align=center, text width=5.2cm] (camp) at (4.6,-3.2) {\textbf{Pour la campagne}\\ 空气很新鲜 \emph{kōngqì hěn xīnxiān} : air frais\\ 生活很安静 \emph{shēnghuó hěn ānjìng} : vie calme\\ 东西不贵 \emph{dōngxi bù guì} : la vie coûte moins cher\\ 家人在一起 \emph{jiārén zài yìqǐ} : la famille réunie};
\draw[-{Stealth}, thick, popblue] (c) -- (ville);
\draw[-{Stealth}, thick, popgreen] (c) -- (camp);
\end{tikzpicture}
\end{center}
\legende{Carte mentale des arguments : pour la ville / pour la campagne}
\end{popfigure}

\courssub{Travail de lecture : de l'écoute à la lecture à deux}

\begin{methode}[Les quatre étapes du travail de lecture en classe]
\begin{enumerate}
\item \textbf{Écoute sans texte.} Le professeur lit le dialogue deux fois, livre fermé. Écoutez les tons, pas seulement les mots. Levez la main quand vous entendez un mot interrogatif (\zh{为什么}, \zh{怎么样}).
\item \textbf{Répétition en chœur.} Toute la classe répète phrase par phrase, en imitant le rythme et les tons du professeur. Le professeur corrige immédiatement les tons fautifs.
\item \textbf{Lecture à deux.} En binôme, l'un joue A, l'autre B, puis on échange les rôles. Le texte reste ouvert : l'objectif est la fluidité.
\item \textbf{Lecture sans pinyin.} Cachez le pinyin et ne gardez que les caractères : c'est la version de l'évaluation.
\end{enumerate}
\end{methode}

\begin{attention}[Les tons qui changent le sens dans ces dialogues]
\begin{itemize}
\item \zh{买} \emph{mǎi} (3\up{e} ton) « acheter » et \zh{卖} \emph{mài} (4\up{e} ton) « vendre » : au marché de Mokolo, la confusion est spectaculaire !
\item \zh{问} \emph{wèn} (4\up{e} ton) « demander » et \zh{闻} \emph{wén} (2\up{e} ton) « sentir ».
\item \zh{工作} \emph{gōngzuò} : deux tons (1\up{er} + 4\up{e}) ; ne dites pas \emph{gòngzuò}.
\item \zh{觉得} \emph{juéde} : le second caractère est au ton neutre, léger et court.
\end{itemize}
\end{attention}

\courssub{La technique de la substitution}

\begin{methode}[Substituer : changer un seul élément à la fois]
La substitution consiste à \textbf{garder exactement la structure} du dialogue modèle et à ne remplacer \textbf{qu'un seul mot} (ou groupe de mots) par un autre élément du vocabulaire. Procédez ainsi :
\begin{enumerate}
\item Choisissez la phrase modèle : \zh{因为农村的工作很少。}
\item Identifiez l'élément variable : le nom \zh{工作} « travail ».
\item Remplacez-le par un autre argument de la carte mentale : \zh{学校} « écoles », \zh{医院} « hôpitaux », \zh{机会} « opportunités ».
\item Adaptez si nécessaire l'adjectif : \zh{很少} « très peu nombreux » devient \zh{不太好} « pas très bonnes » pour les écoles.
\item Vérifiez la prononciation du nouveau mot, puis dites la phrase entière à voix haute.
\end{enumerate}
\textbf{Exemples de substitutions réussies :}\\
\zh{因为农村的学校不太好，我想去城市学习。} \emph{Yīnwèi nóngcūn de xuéxiào bù tài hǎo, wǒ xiǎng qù chéngshì xuéxí.} — « Parce que les écoles de la campagne ne sont pas très bonnes, je veux aller étudier en ville. »\\
\zh{农村的空气很新鲜，可是工作很少。} \emph{Nóngcūn de kōngqì hěn xīnxiān, kěshì gōngzuò hěn shǎo.} — « L'air de la campagne est frais, mais il y a peu de travail. »
\end{methode}

\begin{experience}[Atelier guidé : substitution en chaîne]
En groupe de 4, formez un cercle. Le premier élève dit une phrase du dialogue (ex. \zh{城市有更多的机会。}). Le suivant change \textbf{un seul élément} (lieu, adjectif, nom) et répète la phrase modifiée. On continue jusqu'à épuisement des arguments de la carte mentale. Le professeur note les phrases au tableau et corrige la prononciation en temps réel.
\end{experience}

\begin{exoresolu}[Substitution guidée]
Énoncé : à partir de la phrase modèle \zh{城市有更多的机会。}, produisez deux nouvelles phrases en remplaçant d'abord \zh{机会} par \zh{工作}, puis \zh{城市} par \zh{农村} avec l'adjectif contraire.
\tcblower
Solution détaillée :
\begin{enumerate}
\item On garde la structure : lieu + \zh{有} + \zh{更} + adjectif + \zh{的} + nom. On remplace \zh{机会} par \zh{工作} ; comme \zh{工作} se quantifie avec \zh{多}, on obtient : \zh{城市有更多的工作。} \emph{Chéngshì yǒu gèng duō de gōngzuò.} — « La ville offre plus de travail. »
\item On change le lieu : \zh{农村}. Mais la campagne a \emph{moins} d'opportunités : on remplace \zh{更多} par la négation \zh{没有太多} : \zh{农村没有太多的机会。} \emph{Nóngcūn méiyǒu tài duō de jīhuì.} — « La campagne n'offre pas beaucoup d'opportunités. » Attention : avec la négation, on utilise \zh{没有} et non \zh{不有}.
\end{enumerate}
\end{exoresolu}

\begin{exercice}[À vous de parler]
En binôme, reconstruisez le Dialogue 2 en remplaçant : (a) \zh{学校} par \zh{医院} dans la dernière réplique ; (b) \zh{空气} par \zh{生活} dans la troisième réplique ; (c) terminez par votre propre opinion avec \zh{我觉得……}. Présentez votre nouveau dialogue devant la classe.
\end{exercice}

\begin{corrige}[Exercice — À vous de parler]
Corrigé indicatif : (a) \zh{你说得对，但是农村的医院不太好。} \emph{Nǐ shuō de duì, dànshì nóngcūn de yīyuàn bù tài hǎo.} — « Tu as raison, mais les hôpitaux de la campagne ne sont pas très bons. » (b) \zh{可是农村的生活很安静。} \emph{Kěshì nóngcūn de shēnghuó hěn ānjìng.} — « Mais la vie à la campagne est très calme. » (c) Exemple : \zh{我觉得年轻人可以先在城市工作，以后回农村。} \emph{Wǒ juéde niánqīngrén kěyǐ xiān zài chéngshì gōngzuò, yǐhòu huí nóngcūn.} — « Je pense que les jeunes peuvent d'abord travailler en ville, puis retourner au village. » Toute phrase correcte dans sa structure et ses tons est acceptée.
\end{corrige}

\courssub{Préparer l'exposé évalué : structure et entraînement}

\begin{methode}[Construire son exposé de 1 à 2 minutes]
Suivez ce canevas pour rédiger votre brouillon (caractères + pinyin), puis entraînez-vous à le dire sans notes.
\begin{enumerate}
\item \textbf{Accroche} (1 phrase) : \zh{大家好，我今天要谈谈农村青年进城的问题。} \emph{Dàjiā hǎo, wǒ jīntiān yào tántan nóngcūn qīngnián jìn chéng de wèntí.}
\item \textbf{Constat} (1 phrase avec \zh{因为……}) : \zh{因为农村工作少，很多年轻人离开家乡。} \emph{Yīnwèi nóngcūn gōngzuò shǎo, hěnduō niánqīngrén líkāi jiāxiāng.}
\item \textbf{Argument ville} (1 phrase) : \zh{城市有更多的机会，生活方便。} \emph{Chéngshì yǒu gèng duō de jīhuì, shēnghuó fāngbiàn.}
\item \textbf{Argument campagne} (1 phrase avec \zh{可是}) : \zh{可是农村空气新鲜，家人在一起。} \emph{Kěshì nóngcūn kōngqì xīnxiān, jiārén zài yìqǐ.}
\item \textbf{Opinion personnelle} (1 phrase avec \zh{我觉得……应该……}) : \zh{我觉得年轻人应该先去城市学本领，以后回乡建设家乡。} \emph{Wǒ juéde niánqīngrén yīnggāi xiān qù chéngshì xué běnlǐng, yǐhòu huí xiāng jiànshè jiāxiāng.}
\item \textbf{Conclusion} (relance) : \zh{你们怎么看？} \emph{Nǐmen zěnme kàn?}
\end{enumerate}
\textbf{Conseil :} Chronométrez-vous. Visez 90 secondes. Enregistrez-vous sur votre téléphone pour vérifier vos tons (surtout \zh{工作}, \zh{觉得}, \zh{应该}).
\end{methode}

\begin{savaistu}[春运 chūnyùn : le plus grand déplacement humain du monde]
Chaque année, autour du Nouvel An chinois, des centaines de millions de travailleurs migrants quittent les grandes villes pour retourner dans leur village : c'est la \cle{春运} \emph{chūnyùn}, « le transport de la fête du Printemps ». Les trains et les bus sont bondés pendant plusieurs semaines. Ce phénomène montre que l'exode rural n'est pas un aller simple : les liens avec le village restent très forts en Chine, comme au Cameroun où beaucoup de citadins retournent « au village » pour les fêtes de fin d'année.
\end{savaistu}

\begin{aretenir}[L'essentiel de la section]
\begin{itemize}
\item Question : Sujet + \zh{为什么} + verbe ? Réponse : \zh{因为} + cause.
\item Opinion : \zh{我觉得……} ; désaccord : \zh{我不同意。} ; concession : \zh{你说得对，但是……}
\item \zh{应该} se place avant le verbe ; la négation de \zh{有} est \zh{没有}.
\item Technique de la substitution : une seule modification à la fois, structure intacte.
\item Progression de la lecture : écoute, chœur, binôme, lecture sans pinyin.
\item Exposé type : Accroche → Constat (\zh{因为}) → Arguments ville/campagne (\zh{可是}) → Avis (\zh{我觉得……应该……}) → Relance.
\end{itemize}
\end{aretenir}

\coursec{Travail sur les tons et la prononciation}

Dans la section précédente, nous avons lu et compris les dialogues modèles sur l'exode rural. Chaque phrase y était écrite en caractères, en pinyin et en français. Mais le pinyin n'est pas une décoration : les petits accents au-dessus des voyelles sont la \cle{partition musicale} de la langue chinoise. Un mot mal prononcé en chinois n'est pas un mot « mal dit », c'est souvent \emph{un autre mot}. Cette section est donc consacrée à l'entraînement de la bouche et de l'oreille : nous allons travailler les tons des mots clés de notre thème, la fameuse règle du sandhi du 3\up{e} ton, l'intonation des questions, puis nous finirons par des exercices de discrimination et un choral drill.

\courssub{Observer d'abord : les tons des mots clés de la leçon}

Reprenons les quatre mots essentiels de notre thème « exode rural » et écoutons-les attentivement (le professeur les prononce, les élèves répètent) :

\begin{center}
\begin{tabularx}{\linewidth}{|X|X|X|X|}
\hline
\rowcolor{popblueL} Caractères & Pinyin & Tons & Traduction \\
\hline
农村 & nóngcūn & 2 -- 1 & la campagne, les zones rurales \\
\hline
城市 & chéngshì & 2 -- 4 & la ville \\
\hline
离开 & líkāi & 2 -- 1 & quitter, partir de \\
\hline
机会 & jīhuì & 1 -- 4 & l'occasion, la chance, l'opportunité \\
\hline
\end{tabularx}
\end{center}

Observez : aucun de ces mots n'a deux tons identiques. Le chinois aime les \emph{contrastes mélodiques} à l'intérieur d'un mot : on monte puis on reste haut (nóngcūn), on monte puis on tombe (chéngshì), on reste haut puis on tombe (jīhuì). C'est ce contraste qui rend la parole chinoise si « chantante » aux oreilles des francophones.

\begin{exemplebox}[Écouter et répéter]
\zh{很多年轻人离开农村去城市找工作。}\\
\emph{Hěn duō niánqīngrén líkāi nóngcūn qù chéngshì zhǎo gōngzuò.}\\
« Beaucoup de jeunes quittent la campagne pour aller chercher du travail en ville. »\\
Répétez cette phrase trois fois : d'abord lentement mot par mot, puis par groupes de sens (很多年轻人 / 离开农村 / 去城市 / 找工作), puis à vitesse normale.
\end{exemplebox}

\courssub{Les quatre tons et le ton neutre : le schéma de référence}

Rappelons les cinq « mélodies » du chinois mandarin, avec un mot de notre leçon pour chacune :

\begin{popfigure}
\begin{tikzpicture}[scale=1.0]
% Ton 1
\begin{scope}[shift={(0,0)}]
\draw[popline, ->] (0,0) -- (0,2.2);
\draw[popline, ->] (0,0) -- (2.4,0);
\draw[popblue, very thick] (0.3,1.7) -- (2.1,1.7);
\node[align=center, popdark] at (1.2,-0.5) {\textbf{1\ier{} ton} : haut et plat};
\node[align=center, popblue] at (1.2,2.5) {\zh{机} jī (dans 机会)};
\end{scope}
% Ton 2
\begin{scope}[shift={(3.4,0)}]
\draw[popline, ->] (0,0) -- (0,2.2);
\draw[popline, ->] (0,0) -- (2.4,0);
\draw[poporange, very thick] (0.3,0.5) -- (2.1,1.7);
\node[align=center, popdark] at (1.2,-0.5) {\textbf{2\up{e} ton} : montant};
\node[align=center, poporange] at (1.2,2.5) {\zh{农} nóng (dans 农村)};
\end{scope}
% Ton 3
\begin{scope}[shift={(6.8,0)}]
\draw[popline, ->] (0,0) -- (0,2.2);
\draw[popline, ->] (0,0) -- (2.4,0);
\draw[popgreen, very thick] (0.3,1.4) .. controls (0.9,0.3) and (1.5,0.3) .. (2.1,1.2);
\node[align=center, popdark] at (1.2,-0.5) {\textbf{3\up{e} ton} : descendant-montant};
\node[align=center, popgreen] at (1.2,2.5) {\zh{我} wǒ (dans 我想)};
\end{scope}
% Ton 4
\begin{scope}[shift={(10.2,0)}]
\draw[popline, ->] (0,0) -- (0,2.2);
\draw[popline, ->] (0,0) -- (2.4,0);
\draw[poppink, very thick] (0.3,1.7) -- (2.1,0.4);
\node[align=center, popdark] at (1.2,-0.5) {\textbf{4\up{e} ton} : tombant};
\node[align=center, poppink] at (1.2,2.5) {\zh{市} shì (dans 城市)};
\end{scope}
% Ton neutre
\begin{scope}[shift={(13.6,0)}]
\draw[popline, ->] (0,0) -- (0,2.2);
\draw[popline, ->] (0,0) -- (2.4,0);
\fill[poppurple] (1.2,1.0) circle (0.08);
\node[align=center, popdark] at (1.2,-0.5) {\textbf{Ton neutre} : court et léger};
\node[align=center, poppurple] at (1.2,2.5) {\zh{吗} ma (question)};
\end{scope}
\end{tikzpicture}
\legende{Les quatre tons et le ton neutre, avec un mot exemple de la leçon pour chacun.}
\end{popfigure}

\begin{aretenir}[Les cinq mélodies en une phrase]
1\ier{} ton : la voix reste \textbf{haute et plate}, comme quand le médecin vous demande de dire « aaaa ». 2\up{e} ton : la voix \textbf{monte}, comme dans le « hein ? » français de surprise. 3\up{e} ton : la voix \textbf{descend puis remonte} (en réalité, dans la parole rapide, elle reste surtout basse). 4\up{e} ton : la voix \textbf{tombe brusquement}, comme un « non ! » sec et décidé. Ton neutre : syllabe \textbf{courte, légère, sans mélodie propre}, comme le « le » français dans « je le vois ».
\end{aretenir}

\begin{methode}[La gestuelle des tons : parler avec la main]
Pour ancrer les tons dans la mémoire du corps, dessinez le contour tonal avec la main pendant que vous parlez :
\begin{enumerate}
\item 1\ier{} ton : la main trace une ligne horizontale devant soi, à hauteur d'épaule.
\item 2\up{e} ton : la main monte en diagonale, comme pour poser une question.
\item 3\up{e} ton : la main descend puis remonte, en creux, comme un « V » arrondi.
\item 4\up{e} ton : la main tombe d'un coup sec, comme pour couper.
\item Ton neutre : petit geste bref du poignet, sans ampleur.
\end{enumerate}
Appliquez cette gestuelle à nos mots clés : \zh{农村} nóngcūn (la main monte, puis trace un plat haut), \zh{城市} chéngshì (la main monte, puis tombe net), \zh{离开} líkāi (monte, puis plat haut), \zh{机会} jīhuì (plat haut, puis tombe net). Après quelques séances, la main deviendra inutile : la mélodie sera dans la voix.
\end{methode}

\courssub{Paires minimales : entraîner l'oreille à distinguer}

Une \cle{paire minimale} est un couple de mots qui ne diffèrent que par un seul trait sonore — ici, les tons. Si l'oreille entend la différence, la bouche finira par la produire.

\begin{exemplebox}[Paires minimales de la leçon]
\textbf{Paire 1 :} \zh{工作} gōngzuò (1--4, « le travail ») et \zh{中国} Zhōngguó (1--2, « la Chine »).\\
Même première syllabe haute et plate ; tout se joue sur la deuxième : zuò \emph{tombe} (4\up{e} ton), guó \emph{monte} (2\up{e} ton).\\
\textbf{Paire 2 :} \zh{很少} hěn shǎo (3--3, « très peu ») et \zh{很少} prononcé hén shǎo (2--3).\\
Surprise : les deux sont « corrects » ! La deuxième version est même la seule prononciation naturelle, à cause de la règle du sandhi que nous allons voir. À l'écrit, on note toujours hěn shǎo ; à l'oral, on dit hén shǎo.\\
\textbf{Paire 3 (bonus) :} \zh{城市} chéngshì (la ville) et une mauvaise prononciation chéngshī avec deux tons plats — qui ferait perdre le sens et pourrait faire penser à \zh{成诗} chéngshī (« devenir un poème »).
\end{exemplebox}

\begin{attention}[Erreur fréquente des francophones : 城市 avec deux tons plats]
Le français n'a pas de tons : notre tendance naturelle est de prononcer toutes les syllabes sur le même niveau, « à plat ». Beaucoup d'élèves disent donc \emph{chéngshī} au lieu de \emph{chéngshì}. Or le 4\up{e} ton de \zh{市} doit \textbf{tomber net}, comme un « voilà ! » décidé. Testez-vous : dites « la VILLE ! » en français en appuyant sur la fin — c'est exactement l'énergie du 4\up{e} ton. Même vigilance pour \zh{机会} jīhuì : le \zh{会} huì doit tomber, sinon « l'occasion » devient un son mou et méconnaissable.
\end{attention}

\courssub{La règle du sandhi du 3\up{e} ton}

\begin{propriete}[Deux 3\up{e} tons consécutifs : le premier devient un 2\up{e} ton]
\textbf{Règle :} quand deux syllabes au 3\up{e} ton se suivent, la première se prononce comme un 2\up{e} ton (montant). La deuxième garde son 3\up{e} ton.\\
\textbf{Schéma :} 3 + 3 $\rightarrow$ 2 + 3.\\
\textbf{Important :} l'écriture pinyin ne change pas ! On écrit toujours les tons d'origine ; seule la prononciation change.\\
\textbf{Exemples :}
\begin{itemize}
\item \zh{很好} hěn hǎo se prononce \emph{hén hǎo} « très bien » ;
\item \zh{我想} wǒ xiǎng se prononce \emph{wó xiǎng} « je pense, je voudrais » ;
\item \zh{很少} hěn shǎo se prononce \emph{hén shǎo} « très peu » ;
\item \zh{你好} nǐ hǎo se prononce \emph{ní hǎo} « bonjour ».
\end{itemize}
\textbf{Cas particulier :} devant tout autre ton (1\ier{}, 2\up{e}, 4\up{e} ou neutre), le 3\up{e} ton ne fait pas sa courbe complète : il devient un \cle{demi-3\up{e} ton}, c'est-à-dire qu'on ne prononce que la partie basse, sans la remontée.
\end{propriete}

\begin{exemplebox}[Le demi-3\up{e} ton en action]
\zh{我觉得农村的生活很安静。}\\
\emph{Wǒ juéde nóngcūn de shēnghuó hěn ānjìng.}\\
« Je trouve que la vie à la campagne est très calme. »\\
Dans \zh{我觉得} wǒ juéde, le \zh{我} wǒ est suivi d'un 2\up{e} ton (jué) : il se prononce donc en \textbf{demi-3\up{e} ton} — la voix descend légèrement et reste basse, sans remonter. Mais dans \zh{我想} wǒ xiǎng, le \zh{我} est suivi d'un autre 3\up{e} ton : il devient alors un vrai 2\up{e} ton (wó xiǎng). Même caractère, deux comportements différents selon le voisin !
\end{exemplebox}

\begin{attention}[Piège : ne pas écrire le sandhi]
Dans vos rédactions et vos transcriptions, n'écrivez \textbf{jamais} « hén hǎo » ou « ní hǎo » : le pinyin officiel conserve les tons d'origine (hěn hǎo, nǐ hǎo). Le sandhi est une règle de \emph{prononciation}, pas d'orthographe. À l'examen, écrire « wó xiǎng » est une faute de transcription, même si c'est exactement ce qu'on entend !
\end{attention}

\courssub{L'intonation des questions : 吗 monte, 为什么 descend}

Les tons régissent les syllabes ; l'\cle{intonation} régit la phrase entière. En chinois comme en français, la phrase a une mélodie globale — mais attention, elle ne suit pas toujours les mêmes habitudes.

\begin{center}
\begin{tabularx}{\linewidth}{|X|X|X|}
\hline
\rowcolor{popblueL} Type de question & Exemple (caractères + pinyin) & Intonation et traduction \\
\hline
Question avec 吗 & \zh{你想离开农村吗？} Nǐ xiǎng líkāi nóngcūn ma ? & Montante sur 吗, léger et court. « Veux-tu quitter la campagne ? » \\
\hline
Question avec 为什么 & \zh{你为什么想去城市？} Nǐ wèishénme xiǎng qù chéngshì ? & Descendante vers la fin. « Pourquoi veux-tu aller en ville ? » \\
\hline
Réponse affirmative & \zh{我想去城市找工作。} Wǒ xiǎng qù chéngshì zhǎo gōngzuò. & Descendante, posée. « Je veux aller en ville chercher du travail. » \\
\hline
\end{tabularx}
\end{center}

\begin{propriete}[Deux mélodies de questions à ne pas confondre]
\textbf{Question fermée avec 吗 (oui/non) :} la fin de phrase \textbf{monte} légèrement, comme en français (« Tu viens ? »). Le \zh{吗} lui-même est au ton neutre : court, léger, sans appuyer.\\
\textbf{Question ouverte avec 为什么 / 什么 / 哪儿 (mot interrogatif) :} la fin de phrase \textbf{descend}, car l'information demandée est précise. Ne faites pas monter la voix sur \zh{城市} dans « 你为什么想去城市？ » : cela donnerait une étrange impression d'hésitation.\\
\textbf{Comparaison avec le français :} en français, presque toutes les questions montent ; en chinois, seules les questions en 吗 montent nettement. C'est un réflexe à désapprendre.
\end{propriete}

\begin{savaistu}[Un mauvais ton peut changer le sens : acheter ou vendre ?]
En chinois, \zh{买} mǎi (3\up{e} ton) signifie « acheter » et \zh{卖} mài (4\up{e} ton) signifie « vendre » : même syllabe « mai », deux tons, deux actions exactement contraires ! Au marché de Yaoundé ou de Pékin, confondre les deux peut transformer votre phrase « je veux acheter des légumes » \zh{我想买菜} (wǒ xiǎng mǎi cài) en « je veux vendre des légumes » \zh{我想卖菜} (wǒ xiǎng mài cài). Très utile pour notre thème : celui qui part en ville \emph{achète} (mǎi) des produits, et le commerçant les \emph{vend} (mài). Astuce de mémorisation : le 4\up{e} ton de mài « tombe » comme l'argent qui tombe dans la caisse du vendeur.
\end{savaistu}

\courssub{Exercices de discrimination et choral drill}

\begin{exoresolu}[Discrimination des tons : lève la bonne carte]
\textbf{Énoncé.} Le professeur prépare quatre cartes par élève : 1, 2, 3, 4 (les quatre tons). Il prononce à voix haute, une par une, les syllabes suivantes tirées du vocabulaire de la leçon : nóng, shì, wǒ, jī, ma. Pour chaque syllabe, les élèves lèvent la carte du ton entendu. Donnez les réponses attendues.\\
\tcblower
\textbf{Solution détaillée.}
\begin{itemize}
\item nóng (\zh{农}, dans 农村) : 2\up{e} ton, montant $\rightarrow$ carte \textbf{2}.
\item shì (\zh{市}, dans 城市) : 4\up{e} ton, tombant net $\rightarrow$ carte \textbf{4}.
\item wǒ (\zh{我}) : 3\up{e} ton, la voix descend et reste basse $\rightarrow$ carte \textbf{3}.
\item jī (\zh{机}, dans 机会) : 1\ier{} ton, haut et plat $\rightarrow$ carte \textbf{1}.
\item ma (\zh{吗}) : ton neutre, court et léger $\rightarrow$ aucune carte ! C'est le piège de l'exercice : le ton neutre n'a pas de numéro. Les élèves doivent garder les cartes baissées et dire « ton neutre ».
\end{itemize}
\textbf{Variante plus difficile :} le professeur prononce des mots entiers (gōngzuò / Zhōngguó) et les élèves lèvent la carte du \emph{dernier} ton entendu : 4 pour gōngzuò, 2 pour Zhōngguó.
\end{exoresolu}

\begin{experience}[Choral drill : répéter en tapant les tons]
\textbf{Déroulement (10 minutes).}
\begin{enumerate}
\item Le professeur tape le rythme des tons sur la table : un coup long et plat pour le 1\ier{} ton, un coup glissé vers le haut (main qui monte) pour le 2\up{e}, un coup en creux pour le 3\up{e}, un coup sec pour le 4\up{e}.
\item La classe répète en chœur, en tapant en même temps, les phrases du dialogue de la section précédente :\\
\zh{很多年轻人离开农村。} \emph{Hěn duō niánqīngrén líkāi nóngcūn.}\\
\zh{城市有更多的机会。} \emph{Chéngshì yǒu gèng duō de jīhuì.}\\
\zh{你为什么想去城市？} \emph{Nǐ wèishénme xiǎng qù chéngshì ?}
\item Deuxième passage à vitesse normale, sans le professeur ; un élève « chef d'orchestre » donne le tempo.
\item Troisième passage par groupes : le groupe A pose la question, le groupe B répond, puis on échange.
\end{enumerate}
\textbf{Objectif :} que la mélodie des tons devienne un automatisme du corps, pas un calcul de la tête.
\end{experience}

\begin{exercice}[À vous de jouer]
\begin{enumerate}
\item Indiquez les tons (en chiffres) des mots suivants, puis prononcez-les à voix haute avec la gestuelle : 农村, 城市, 离开, 机会, 工作.
\item Comment se prononcent réellement les groupes suivants ? Écrivez la prononciation effective : 很好 (hěn hǎo), 我想 (wǒ xiǎng), 很少 (hěn shǎo).
\item Lisez à voix haute en respectant l'intonation : \zh{你想去城市吗？} puis \zh{你为什么想去城市？} Quelle différence de mélodie finale ?
\item Vrai ou faux : dans une transcription pinyin, on écrit « ní hǎo » parce que c'est ce qu'on entend. Justifiez.
\end{enumerate}
\end{exercice}

\begin{corrige}[Exercice « À vous de jouer »]
\begin{enumerate}
\item 农村 nóngcūn : 2--1 ; 城市 chéngshì : 2--4 ; 离开 líkāi : 2--1 ; 机会 jīhuì : 1--4 ; 工作 gōngzuò : 1--4.
\item 很好 se prononce hén hǎo (2--3) ; 我想 se prononce wó xiǎng (2--3) ; 很少 se prononce hén shǎo (2--3). Règle : 3 + 3 $\rightarrow$ 2 + 3.
\item La question avec \zh{吗} monte à la fin (吗 léger, ton neutre) ; la question avec \zh{为什么} descend vers la fin, la voix retombe sur \zh{城市}.
\item Faux. Le sandhi est une règle de prononciation, pas d'orthographe : on écrit toujours les tons d'origine, donc nǐ hǎo, même si l'on prononce ní hǎo.
\end{enumerate}
\end{corrige}

\begin{aretenir}[Bilan de la section]
\begin{itemize}
\item Les quatre mots clés du thème ont des tons contrastés : nóngcūn (2--1), chéngshì (2--4), líkāi (2--1), jīhuì (1--4) ; le 4\up{e} ton doit \textbf{tomber net}.
\item Deux 3\up{e} tons qui se suivent : le premier devient un 2\up{e} ton à l'oral (hěn hǎo $\rightarrow$ hén hǎo), mais l'écriture pinyin ne change jamais.
\item Devant un autre ton, le 3\up{e} ton devient un demi-3\up{e} ton (wǒ juéde : la voix reste basse).
\item Question avec \zh{吗} : intonation montante ; question avec \zh{为什么} : intonation descendante.
\item La gestuelle des tons et le choral drill transforment la mélodie en automatisme : la main dessine, la voix suit.
\end{itemize}
Vous êtes maintenant prêts pour la section suivante : les jeux de rôle et le débat guidé, où ces tons travaillés prendront vie dans de vrais échanges sur l'exode rural.
\end{aretenir}

\coursec{Leçon 25 — Expression orale : exode rural}

\courssub{Vocabulaire clé : ville, campagne, exode rural}

\begin{definition}[Lexique de la leçon (HSK 2–3)]
Le tableau ci-dessous regroupe le vocabulaire essentiel pour comprendre les dialogues, participer aux jeux de rôle et construire vos arguments lors du débat. Mémorisez les tons avec soin.
\end{definition}

\begin{center}
\begin{tabularx}{\linewidth}{|X|X|X|X|}
\hline
\rowcolor{popblueL} Caractères & Pinyin & Nature & Traduction \\
\hline
城市 & chéngshì & n. & ville \\
\hline
农村 & nóngcūn & n. & campagne, zone rurale \\
\hline
找工作 & zhǎo gōngzuò & v. & chercher du travail \\
\hline
房租 & fángzū & n. & loyer \\
\hline
后悔 & hòuhuǐ & v. & regretter \\
\hline
机会 & jīhuì & n. & opportunité, occasion \\
\hline
担心 & dānxīn & v. & s'inquiéter, craindre \\
\hline
明白 & míngbai & v. & comprendre \\
\hline
为什么 & wèishénme & adv. & pourquoi \\
\hline
怎么样 & zěnmeyàng & adv. & comment, qu'en est-il ? \\
\hline
还是 & háishì & conj. & ou (question alternative) \\
\hline
或者 & huòzhě & conj. & ou (énoncé affirmatif) \\
\hline
觉得 & juéde & v. & penser, trouver (avis) \\
\hline
因为 & yīnwèi & conj. & parce que \\
\hline
所以 & suǒyǐ & conj. & donc, par conséquent \\
\hline
首先 & shǒuxiān & adv. & d'abord, en premier lieu \\
\hline
然后 & ránhòu & adv. & ensuite, puis \\
\hline
最后 & zuìhòu & adv. & enfin, en dernier lieu \\
\hline
不同意 & bù tóngyì & v. & ne pas être d'accord \\
\hline
总结 & zǒngjié & v./n. & conclure / conclusion \\
\hline
各有各的 & gè yǒu gè de & loc. & chacun a ses propres (avantages) \\
\hline
乡村振兴 & xiāngcūn zhènxīng & n. & revitalisation rurale (programme chinois) \\
\hline
人口外流 & rénkǒu wàiliú & n. & exode rural (litt. écoulement de population) \\
\hline
\end{tabularx}
\end{center}

\courssub{Formules fonctionnelles : présenter, opiner, réagir, relancer}

\begin{propriete}[Boîte à outils communicatifs pour le débat]
Ces formules sont les « briques » de votre expression orale. Elles sont fixes : apprenez-les par cœur comme des blocs indivisibles.
\begin{center}
\begin{tabularx}{\linewidth}{|X|X|X|X|}
\hline
\rowcolor{popblueL} Fonction & Formule chinoise & Pinyin & Équivalent français \\
\hline
Ouvrir / Présenter & \zh{你好！请问……} & \emph{Nǐ hǎo! Qǐngwèn…} & Bonjour ! Puis-je vous demander… \\
\hline
Poser une question (cause) & \zh{你为什么……?} & \emph{Nǐ wèishénme…?} & Pourquoi… ? \\
\hline
Poser une question (état) & \zh{……怎么样?} & \emph{… zěnmeyàng?} & Comment est… ? / Qu'en est-il de… ? \\
\hline
Répondre (cause) & \zh{因为……} & \emph{Yīnwèi…} & Parce que… \\
\hline
Donner son avis & \zh{我觉得……} & \emph{Wǒ juéde…} & Je pense que… / Je trouve que… \\
\hline
Exprimer un désir & \zh{我想……} & \emph{Wǒ xiǎng…} & Je veux… / J'aimerais… \\
\hline
Exprimer une inquiétude & \zh{我担心……} & \emph{Wǒ dānxīn…} & Je m'inquiète pour… / Je crains que… \\
\hline
Nuancer / Contredire & \zh{可是……} / \zh{但是……} & \emph{Kěshì… / Dànshì…} & Mais… / Cependant… \\
\hline
Réfuter poliment & \zh{我不同意。可是……} & \emph{Wǒ bù tóngyì. Kěshì…} & Je ne suis pas d'accord. Mais… \\
\hline
Relancer l'interlocuteur & \zh{你呢？你觉得怎么样？} & \emph{Nǐ ne? Nǐ juéde zěnmeyàng?} & Et toi ? Qu'en penses-tu ? \\
\hline
Structurer (énumérer) & \zh{首先……然后……最后……} & \emph{Shǒuxiān… ránhòu… zuìhòu…} & D'abord… ensuite… enfin… \\
\hline
Conclure & \zh{总之，……} & \emph{Zǒngzhī…} & En bref… / En résumé… \\
\hline
Clore harmonieusement & \zh{各有各的好处。} & \emph{Gè yǒu gè de hǎochù.} & Chacun a ses avantages. \\
\hline
\end{tabularx}
\end{center}
\end{propriete}

\courssub{Dialogues modèles avec pinyin}

\begin{textebox}[Dialogue 1 : À la gare de Yaoundé — Interview CRTV]
\textbf{Contexte :} Un journaliste (\zh{记者}, \emph{jìzhě}) interroge un jeune migrant (\zh{青年}, \emph{qīngnián}) fraîchement arrivé de son village. Repérez \zh{为什么} et \zh{怎么样}.
记者：你好！请问，你为什么来雅温得？\\
\emph{Jìzhě : Nǐ hǎo! Qǐngwèn, nǐ wèishénme lái Yǎwēndé?}\\
« Journaliste : Bonjour ! Puis-je vous demander pourquoi vous êtes venu à Yaoundé ? »\\[4pt]
青年：因为村里没有工作，我想在城里找工作。\\
\emph{Qīngnián : Yīnwèi cūn lǐ méiyǒu gōngzuò, wǒ xiǎng zài chéng lǐ zhǎo gōngzuò.}\\
« Jeune : Parce qu'il n'y a pas de travail au village, je veux en chercher en ville. »\\[4pt]
记者：你在雅温得的生活怎么样？\\
\emph{Jìzhě : Nǐ zài Yǎwēndé de shēnghuó zěnmeyàng?}\\
« Journaliste : Comment se passe votre vie à Yaoundé ? »\\[4pt]
青年：不太容易。房租很贵，但是我不后悔。\\
\emph{Qīngnián : Bú tài róngyì. Fángzū hěn guì, dànshì wǒ bù hòuhuǐ.}\\
« Jeune : Pas facile. Le loyer est cher, mais je ne regrette pas. »
\end{textebox}

\begin{propriete}[Analyse grammaticale : \zh{为什么} et \zh{怎么样}]
\begin{itemize}
\item \cle{为什么} (\emph{wèishénme}) « pourquoi » : \textbf{Sujet + 为什么 + Verbe/Adjectif ?} \\
La réponse attendue utilise \zh{因为} (\emph{yīnwèi}) « parce que » (souvent suivi de \zh{所以} « donc » dans la phrase complète). \\
Ex. \zh{你为什么离开农村？} → \zh{因为我想学习。(所以我来了。)}
\item \cle{怎么样} (\emph{zěnmeyàng}) « comment / qu'en est-il » : \textbf{Sujet/Nom + 怎么样 ?} \\
La réponse est un adjectif qualificatif (souvent précédé de \zh{很}, \zh{不太}, \zh{真}). \textbf{Pas de verbe \zh{是} avant l'adjectif.} \\
Ex. \zh{城市的生活怎么样？} → \zh{很方便。/ 不太安静。}
\end{itemize}
\textbf{Contraste français-chinois :} En français, l'interrogatif monte en tête (\emph{Pourquoi viens-tu ?}). En chinois, l'ordre \textbf{SVO} est intangible : \zh{你 为什么 来} (Sujet - Interrogatif - Verbe).
\end{propriete}

\begin{textebox}[Dialogue 2 : Réunion de famille près de Bafoussam — Partir ou rester ?]
\textbf{Contexte :} \zh{小王} (\emph{Xiǎo Wáng}) annonce son départ pour Douala. Ses parents réagissent. Repérez \zh{我觉得}, \zh{我想}, \zh{担心}, \zh{了} (hypothèse accomplie).
儿子：爸爸，妈妈，我想去杜阿拉。我觉得城里机会多。\\
\emph{Érzi : Bàba, māma, wǒ xiǎng qù Dù'ālā. Wǒ juéde chéng lǐ jīhuì duō.}\\
« Fils : Papa, maman, je veux aller à Douala. Je trouve qu'en ville il y a des opportunités. »\\[4pt]
父亲：可是农村也需要年轻人。你走了，谁帮助我们？\\
\emph{Fùqīn : Kěshì nóngcūn yě xūyào niánqīngrén. Nǐ zǒu le, shéi bāngzhù wǒmen?}\\
« Père : Mais la campagne a besoin de jeunes. Si tu pars, qui nous aidera ? »\\[4pt]
母亲：我担心你。城市的生活很贵，也很难。\\
\emph{Mǔqīn : Wǒ dānxīn nǐ. Chéngshì de shēnghuó hěn guì, yě hěn nán.}\\
« Mère : Je m'inquiète pour toi. La vie en ville est chère et dure. »\\[4pt]
儿子：我明白。我先去城里学习，然后回来帮助村里。\\
\emph{Érzi : Wǒ míngbai. Wǒ xiān qù chéng lǐ xuéxí, ránhòu huílái bāngzhù cūn lǐ.}\\
« Fils : Je comprends. J'irai d'abord étudier en ville, puis je reviendrai aider le village. »
\end{textebox}

\courssub{Travail sur les tons et la prononciation}

\begin{methode}[Points de vigilance phonétique pour cette leçon]
Avant de passer aux jeux de rôle, entraînez-vous sur ces difficultés tonales fréquentes chez les francophones :
\begin{enumerate}
\item \textbf{Ton 3 + Ton 3 (Sandhi) :} \zh{你好} → \emph{ní hǎo} (2\textsuperscript{e} + 3\textsuperscript{e}) ; \zh{我想} → \emph{wó xiǎng} (2\textsuperscript{e} + 3\textsuperscript{e}) ; \zh{农村} → \emph{nóng cūn} (2\textsuperscript{e} + 1\textsuperscript{er}) — \zh{村} est ton 1 ici.
\item \textbf{\zh{不} (bù) devant ton 4 :} \zh{不后悔} → \emph{bú hòuhuǐ} (2\textsuperscript{e} + 4\textsuperscript{e}) ; \zh{不贵} → \emph{bú guì} ; \zh{不容易} → \emph{bú róngyì}.
\item \textbf{\zh{一} (yī) sandhi :} \zh{一个} → \emph{yí ge} (2\textsuperscript{e}) ; \zh{一样} → \emph{yíyàng} ; mais \zh{第一} → \emph{dì yī} (ton 1 car ordinal).
\item \textbf{Tons 2 et 3 (montant vs descendant-montant) :} \zh{为什么} (\emph{wèi shén me} : 4-2-neutre) ; \zh{怎么样} (\emph{zěn me yàng} : 3-neutre-4) ; \zh{机会} (\emph{jī huì} : 1-4) ; \zh{担心} (\emph{dān xīn} : 1-1).
\item \textbf{Neutres (轻声) :} \zh{怎么样} (\emph{yàng} souvent neutre en rapide) ; \zh{明白} (\emph{bai} souvent neutre) ; \zh{里} (\emph{li} dans \zh{城里}, \zh{村里}).
\end{enumerate}
\textbf{Exercice rapide :} Lisez à voix haute la liste de vocabulaire du début de leçon en exagérant les contours tonaux. Enregistrez-vous si possible.
\end{methode}

\courssub{Jeux de rôle et débat guidé}

Maintenant que votre prononciation est plus assurée, passons à l'action : parler, échanger, débattre. Vous allez d'abord jouer des scènes de la vie réelle, puis participer à un débat organisé : \zh{城市好还是农村好？} (\emph{Chéngshì hǎo háishì nóngcūn hǎo?}). L'objectif est de \textbf{communiquer} : exprimer sa pensée, réagir à l'autre, conclure.

\courssub{Jeu de rôle 1 : Le journaliste de CRTV interviewe un jeune migrant}

\begin{experience}[Jeu de rôle : L'interview à la gare]
\textbf{Consigne :} Par groupes de deux. Élève A = Journaliste CRTV. Élève B = Jeune migrant arrivé à Yaoundé (ou Douala).
\begin{enumerate}
\item \textbf{Préparation (3 min) :}
\begin{itemize}
\item Journaliste : Prépare \textbf{4 questions minimum} : 2 avec \zh{为什么}, 2 avec \zh{怎么样}. Ex. \zh{你为什么不留在村里？}, \zh{雅温得的交通怎么样？}
\item Migrant : Prépare les réponses avec \zh{因为……所以……} et au moins 3 adjectifs (\zh{贵}, \zh{忙}, \zh{方便}, \zh{累}, \zh{有意思}, \zh{孤单}).
\end{itemize}
\item \textbf{Jeu (2 min) :} Jouez la scène. Journaliste ouvre par \zh{你好！请问，…} et clôture par \zh{谢谢你！再见！} (\emph{Xièxie nǐ! Zàijiàn!}).
\item \textbf{Inversion :} Échangez les rôles et recommencez avec un autre profil (étudiant, commerçant, footballeur…).
\end{enumerate}
\end{experience}

\courssub{Jeu de rôle 2 : Réunion de famille — Partir ou rester ?}

\begin{experience}[Jeu de rôle : Discussion familiale]
\textbf{Consigne :} Groupes de 4 : Fils/Fille (\zh{儿子/女儿}), Père (\zh{父亲}), Mère (\zh{母亲}), Oncle médiateur (\zh{舅舅} ou \zh{叔叔}).
\begin{enumerate}
\item \textbf{Préparation (3 min) :} Chacun écrit \textbf{2 phrases} minimum utilisant les structures : \zh{我想…}, \zh{我觉得…}, \zh{我担心…}, \zh{可是…}, \zh{先…然后…}.
\item \textbf{Déroulement (3 min) :} Le fils/fille commence. \textbf{Écoutez et réagissez}, ne récitez pas. L'oncle doit relancer avec \zh{你呢？} et conclure avec \zh{我觉得你们都对，但是……} (\emph{Wǒ juéde nǐmen dōu duì, dànshì…}).
\item \textbf{Variante :} Changez le dénouement : le jeune décide de rester ; ou il part mais promet d'envoyer de l'argent (\zh{寄钱}, \emph{jì qián}).
\end{enumerate}
\end{experience}

\courssub{Le débat guidé : \zh{城市好还是农村好？}}

\begin{propriete}[La question alternative avec \zh{还是} vs \zh{或者}]
\textbf{Structure :} \textbf{A + 还是 + B ?} (Choix exclusif dans une \emph{question}). \\
\zh{城市好还是农村好？} \emph{Chéngshì hǎo háishì nóngcūn hǎo?} « La ville ou la campagne : laquelle est mieux ? » \\
\zh{你想住在城市还是农村？} \emph{Nǐ xiǎng zhù zài chéngshì háishì nóngcūn?} \\
\textbf{Attention (piège classique) :} Dans une phrase \emph{affirmative}, « ou » se dit \zh{或者} (\emph{huòzhě}). \\
\zh{我想住在城市或者农村。} \emph{Wǒ xiǎng zhù zài chéngshì huòzhě nóngcūn.} « Je veux habiter en ville ou à la campagne (peu importe laquelle). »
\end{propriete}

\textbf{Organisation du débat :} La classe est divisée en deux équipes. Chaque élève tient un rôle précis avec formules imposées.

\begin{center}
\begin{tabularx}{\linewidth}{|X|X|X|X|}
\hline
\rowcolor{popblueL} Rôle & Formules imposées & Pinyin & Fonction \\
\hline
Animateur & \zh{你呢？你觉得怎么样？} & \emph{Nǐ ne? Nǐ juéde zěnmeyàng?} & Donner la parole, relancer \\
\hline
Argumenteur & \zh{我觉得……因为……} & \emph{Wǒ juéde… yīnwèi…} & Poser un argument justifié \\
\hline
Contradicteur & \zh{我不同意。可是……} & \emph{Wǒ bù tóngyì. Kěshì…} & Réfuter poliment \\
\hline
Rapporteur & \zh{总之，……} & \emph{Zǒngzhī…} & Synthétiser, clore \\
\hline
\end{tabularx}
\end{center}

\begin{exemplebox}[Arguments types fournis (à enrichir)]
\textbf{Équipe A — \zh{城市好} (Pour la ville) :}
\begin{itemize}
\item \zh{我觉得城市好，因为城里工作多，工资高。} (\emph{Wǒ juéde chéngshì hǎo, yīnwèi chéng lǐ gōngzuò duō, gōngzī gāo.})
\item \zh{城市有好的学校和大医院，生活方便。} (\emph{Chéngshì yǒu hǎo de xuéxiào hé dà yīyuàn, shēnghuó fāngbiàn.})
\end{itemize}
\textbf{Équipe B — \zh{农村好} (Pour la campagne) :}
\begin{itemize}
\item \zh{我觉得农村好，因为农村空气好，环境安静。} (\emph{Wǒ juéde nóngcūn hǎo, yīnwèi nóngcūn kōngqì hǎo, huánjìng ānjìng.})
\item \zh{农村的生活不贵，邻里关系好，不孤单。} (\emph{Nóngcūn de shēnghuó bú guì, línlǐ guānxì hǎo, bù gūdān.})
\end{itemize}
\textbf{Réfutation commune (les deux équipes) :}
\zh{我不同意。城市机会多，可是生活太贵了，压力太大。} (\emph{Wǒ bù tóngyì. Chéngshì jīhuì duō, kěshì shēnghuó tài guì le, yālì tài dà.})
\end{exemplebox}

\begin{popfigure}
\begin{tikzpicture}[
node distance=1.0cm,
boite/.style={draw=popblue, fill=popblueL, rounded corners=4pt, align=center, minimum height=1.0cm, inner sep=6pt, font=\small},
fleche/.style={-{Stealth[length=3mm]}, thick, popdark}]
\node[boite, fill=popgoldL, draw=popgold, font=\normalsize\bfseries] (theme) {主题：城市好还是农村好？};
\node[boite, below left=1.0cm and 0.5cm of theme, align=center] (va){观点 A\\ 我觉得城市好，因为……};
\node[boite, below right=1.0cm and 0.5cm of theme, align=center] (vb){观点 B\\ 我觉得农村好，因为……};
\node[boite, fill=poporangeL, draw=poporange, below=3.0cm of theme, align=center] (refut){反驳\\ 我不同意。可是……};
\node[boite, fill=popgreenL, draw=popgreen, below=1.0cm of refut, align=center] (concl){总结\\ 总之，各有各的好处。};
\draw[fleche] (theme) -- (va);
\draw[fleche] (theme) -- (vb);
\draw[fleche] (va) -- (refut);
\draw[fleche] (vb) -- (refut);
\draw[fleche] (refut) -- (concl);
\end{tikzpicture}
\legende{Déroulement type du débat : Thème → Deux arguments → Réfutation croisée → Synthèse finale.}
\end{popfigure}

\begin{methode}[Préparer sa fiche brouillon (3 minutes)]
Un bon orateur ne lit pas, il \textbf{parle} en s'appuyant sur des mots-clés. Sur une fiche bristol (format carte de visite) :
\begin{enumerate}
\item \textbf{3 arguments chainés :} \zh{首先，……} ; \zh{然后，……} ; \zh{最后，……}.
\item \textbf{1 réfutation prête :} \zh{我不同意，因为……可是……}
\item \textbf{1 conclusion :} \zh{总之，各有各的好处。} ou \zh{总之，我认为……}
\end{enumerate}
\textbf{Astuce :} Écrivez en \textbf{pinyin + mots-clés caractères}, pas des phrases complètes. Levez les yeux vers l'auditoire.
\end{methode}

\begin{exemplebox}[Argument complet modèle (niveau visé)]
\zh{我认为年轻人应该先去城市学习技术，然后回农村创业。}\\
\emph{Wǒ rènwéi niánqīngrén yīnggāi xiān qù chéngshì xuéxí jìshù, ránhòu huí nóngcūn chuàngyè.}\\
« Je pense que les jeunes devraient d'abord aller en ville apprendre une technique, puis retourner en campagne entreprendre. »\\
\textbf{Analyse :} \zh{我认为} (opinion formelle) + \zh{应该} (devoir moral/conseil) + structure séquentielle \zh{先……然后……} + verbe \zh{创业} (créer son activité). Phrase nuancée, structurée, HSK 3+.
\end{exemplebox}

\begin{attention}[Erreurs fatales des francophones à l'oral — À bannir]
\begin{itemize}
\item \textbf{Faux ami \zh{是} + adjectif :} « La ville est meilleure » $\neq$ \zh{城市是更好}. \textbf{Correct :} \zh{城市比较好} / \zh{城市好}. L'adjectif \textbf{est} le prédicat.
\item \textbf{Traduction mot à mot « Il y a » :} « Il y a du travail » $\neq$ \zh{有工作} (seul, phrase incomplète). \textbf{Correct :} \zh{工作多} (le travail est abondant) ou \zh{这儿有很多工作} (ici il y a beaucoup de travail).
\item \textbf{Oubli de \zh{了} aspectuel :} « Si tu pars » $\neq$ \zh{如果你走}. \textbf{Correct :} \zh{你走了} (action réalisée, base de l'hypothèse).
\item \textbf{Mauvaise place de \zh{为什么} :} \zh{为什么你来？} $\rightarrow$ \textbf{Faux.} \textbf{Correct :} \zh{你为什么来？} (Sujet + 为什么 + V).
\end{itemize}
\textbf{Stratégie :} Pensez en \textbf{blocs chinois} (\zh{我觉得+S+V}, \zh{因为+S+V}), pas en français.
\end{attention}

\begin{exoresolu}[Construire un argument d'équipe « Campagne »]
Énoncé : Vous êtes dans l'équipe « Campagne ». Construisez un argument complet \zh{我觉得……因为……} à partir de : \zh{农村 / 生活 / 安静 / 空气 / 好}.
\tcblower
\textbf{Solution détaillée :}
1. Ouvrir par la formule d'opinion : \zh{我觉得农村好} (Je trouve la campagne bien). \\
2. Introduire la cause : \zh{因为}. \\
3. Détailler : Sujet \zh{农村的生活} + adjectif \zh{很安静} ; coordination \zh{空气也很好}. \\
\textbf{Résultat :} \\
\zh{我觉得农村好，因为农村的生活很安静，空气也很好。}\\
\emph{Wǒ juéde nóngcūn hǎo, yīnwèi nóngcūn de shēnghuó hěn ānjìng, kōngqì yě hěn hǎo.}\\
\textbf{Vérification tons :} \zh{农村} (2+1), \zh{安静} (1+4), \zh{空气} (1+4), \zh{很好} (3$\rightarrow$2 + 3).
\end{exoresolu}

\begin{exercice}[À vous de jouer : Production écrite préparatoire à l'oral]
\begin{enumerate}
\item \textbf{Équipe Ville :} Rédigez \textbf{deux arguments distincts} avec \zh{我觉得……因为……} en intégrant : \zh{工作}, \zh{学校}, \zh{方便}, \zh{工资}.
\item \textbf{Équipe Campagne :} Réfutez l'argument adverse \zh{城市工作多} en utilisant \zh{我不同意。可是……} et le vocabulaire : \zh{贵}, \zh{累}, \zh{压力大}.
\item \textbf{Rapporteur :} Rédigez la phrase de conclusion officielle avec \zh{总之} et \zh{各有各的}.
\end{enumerate}
\end{exercice}

\begin{corrige}[Exercice : À vous de jouer]
\begin{enumerate}
\item \zh{我觉得城市好，因为城里工作多，工资高，生活很方便。} (\emph{Wǒ juéde chéngshì hǎo, yīnwèi chéng lǐ gōngzuò duō, gōngzī gāo, shēnghuó hěn fāngbiàn.}) \\
\zh{首先城市有好学校，然后医院也好，最后交通方便。} (\emph{Shǒuxiān chéngshì yǒu hǎo xuéxiào, ránhòu yīyuàn yě hǎo, zuìhòu jiāotōng fāngbiàn.})
\item \zh{我不同意。可是城市生活太贵了，房租高，年轻人压力很大，很累。} (\emph{Wǒ bù tóngyì. Kěshì chéngshì shēnghuó tài guì le, fángzū gāo, niánqīngrén yālì hěn dà, hěn lèi.})
\item \zh{总之，城市和农村各有各的好处，关键看个人选择。} (\emph{Zǒngzhī, chéngshì hé nóngcūn gè yǒu gè de hǎochù, guānjiàn kàn gèrén xuǎnzé.})
\end{enumerate}
\end{corrige}

\courssub{Conclure le débat : Synthèse et tour de table}

Le rapporteur de chaque équipe termine par la phrase de synthèse modèle, à connaître par cœur.

\begin{textebox}[Phrase de synthèse modèle — À mémoriser]
总之，城市和农村各有各的好处。\\
\emph{Zǒngzhī, chéngshì hé nóngcūn gè yǒu gè de hǎochù.}\\
« En résumé, la ville et la campagne ont chacune leurs avantages. »
\end{textebox}

La structure \zh{各有各的…} (\emph{gè yǒu gè de…}, « chacun a son/ses propre(s)… ») est la marque d'une conclusion \textbf{harmonieuse} (\zh{和谐}, \emph{héxié}), valeur centrale en culture chinoise : on ne cherche pas à « gagner » le débat mais à montrer qu'on a compris la complexité.

\textbf{Tour de table final (évaluation formative).} Chaque élève se lève, donne son avis personnel en \textbf{une seule phrase} claire, sans lire :
\begin{itemize}
\item \zh{我觉得城市好，因为……}
\item \zh{我觉得农村好，因为……}
\item \zh{我觉得两个都好，因为……}
\end{itemize}

\begin{savaistu}[Jeunesse et monde rural : regards croisés Cameroun — Chine]
Au Cameroun, les politiques d'\cle{insertion des jeunes} (programmes \emph{PIAASI}, \emph{PADER}, formations en agropastoral) visent à freiner l'exode vers Yaoundé et Douala en valorisant les terroirs. En Chine, la stratégie nationale de \cle{revitalisation rurale} (\zh{乡村振兴}, \emph{xiāngcūn zhènxīng}), lancée en 2017, injecte des milliards dans les infrastructures, l'e-commerce agricole et le tourisme vert pour rendre les villages attractifs. Deux contextes différents, un même défi : \textbf{comment garder la jeunesse au village sans l'enfermer ?} Votre débat d'aujourd'hui rejoint une préoccupation planétaire.
\end{savaistu}

\courssub{Grille d'auto-évaluation (à remplir après le débat)}

Soyez honnête. Cochez : \zh{很好} (très bien) / \zh{好} (bien) / \zh{还要努力} (à retravailler).

\begin{center}
\begin{tabularx}{\linewidth}{|X|X|X|X|}
\hline
\rowcolor{popblueL} Critère d'évaluation & 很好 & 好 & 还要努力 \\
\hline
Prononciation des tons (focus 2\textsuperscript{e}/3\textsuperscript{e} tons, sandhi \zh{不}/\zh{一}) & & & \\
\hline
Usage correct des formules (\zh{我觉得}, \zh{因为}, \zh{总之}, \zh{还是}) & & & \\
\hline
Fluidité (parole continue, sans longs silences, sans lecture) & & & \\
\hline
Respect du rôle assigné (animateur, argumenteur, contradicteur, rapporteur) & & & \\
\hline
Écoute active et réaction pertinente aux camarades & & & \\
\hline
Richesse lexicale (adjectifs variés : \zh{贵/便宜}, \zh{方便/麻烦}, \zh{安静/吵}, \zh{孤单/热闹}) & & & \\
\hline
\end{tabularx}
\end{center}

\begin{aretenir}[L'essentiel de la Leçon 25 — Expression orale : Exode rural]
\begin{itemize}
\item \textbf{Interview :} Questions \zh{为什么} (réponse \zh{因为}) et \zh{怎么样} (réponse = adjectif). Ordre SVO strict.
\item \textbf{Discussion :} Avis avec \zh{觉得 + phrase complète} (jamais \zh{是} + adj). Inquiétude \zh{担心}. Concession \zh{可是/但是}.
\item \textbf{Débat :} Question alternative \zh{A 还是 B ?} (vs \zh{或者} en affirmatif). Réfutation polie \zh{我不同意。可是…}. Conclusion \zh{总之，各有各的好处。}
\item \textbf{Méthode :} Fiche brouillon \zh{首先 / 然后 / 最后} + 1 réfutation + 1 conclusion. Parler \textbf{les yeux levés}.
\item \textbf{Phonétique :} Sandhi \zh{不} (→ \zh{bu} 2\textsuperscript{e} ton devant 4\textsuperscript{e}), sandhi 3\textsuperscript{e} ton, neutres (\zh{里}, \zh{怎么样}).
\end{itemize}
\end{aretenir}

\coursec{Bilan et préparation de l'exposé évalué}

Après nos jeux de rôle et notre débat guidé sur l'exode rural, vous savez désormais échanger, réagir et relancer une conversation en chinois. Il est temps de faire le \cle{bilan} de tout ce que nous avons appris dans cette leçon et de préparer sérieusement l'\cle{exposé évalué} qui clôturera ce module. Cette dernière étape est décisive : c'est elle qui transforme des connaissances éparses en une compétence orale solide et notée.

\courssub{Synthèse du vocabulaire : reclasser pour mémoriser}

Un bon orateur ne récite pas une liste de mots : il les organise par \cle{thèmes}. Reprenons tout le vocabulaire de la leçon et classons-le en trois familles : la \cle{ville}, la \cle{campagne} et la \cle{migration}.

\begin{center}
\begin{tabularx}{\linewidth}{|X|X|X|X|}
\hline
\rowcolor{popblueL} Caractères & Pinyin & Nature & Traduction \\
\hline
城市 & chéngshì & nom & la ville \\
\hline
农村 & nóngcūn & nom & la campagne, les zones rurales \\
\hline
乡下 & xiāngxia & nom & la campagne (familier) \\
\hline
家乡 & jiāxiāng & nom & le pays natal, la ville natale \\
\hline
人口 & rénkǒu & nom & la population \\
\hline
外流 & wàiliú & nom / verbe & l'exode, sortir (d'une population) \\
\hline
打工 & dǎgōng & verbe & travailler (comme ouvrier migrant) \\
\hline
机会 & jīhuì & nom & l'occasion, l'opportunité \\
\hline
工作 & gōngzuò & nom / verbe & le travail ; travailler \\
\hline
生活 & shēnghuó & nom / verbe & la vie ; vivre \\
\hline
方便 & fāngbiàn & adjectif & pratique, commode \\
\hline
空气 & kōngqì & nom & l'air \\
\hline
环境 & huánjìng & nom & l'environnement \\
\hline
安静 & ānjìng & adjectif & calme, tranquille \\
\hline
热闹 & rènao & adjectif & animé, bruyant (joyeux) \\
\hline
交通 & jiāotōng & nom & les transports \\
\hline
医院 & yīyuàn & nom & l'hôpital \\
\hline
学校 & xuéxiào & nom & l'école \\
\hline
年轻人 & niánqīngrén & nom & les jeunes \\
\hline
老人 & lǎorén & nom & les personnes âgées \\
\hline
\end{tabularx}
\end{center}

\begin{exercice}[Reclassement express]
Reclassez les mots suivants en trois colonnes « ville », « campagne », « migration » : \zh{医院}、\zh{空气}、\zh{打工}、\zh{交通}、\zh{安静}、\zh{外流}、\zh{热闹}、\zh{环境}、\zh{机会}。
\end{exercice}

\begin{corrige}[Reclassement express]
\textbf{Ville} : 医院 \emph{yīyuàn} (hôpital), 交通 \emph{jiāotōng} (transports), 热闹 \emph{rènao} (animé). \textbf{Campagne} : 空气 \emph{kōngqì} (air), 安静 \emph{ānjìng} (calme), 环境 \emph{huánjìng} (environnement). \textbf{Migration} : 打工 \emph{dǎgōng} (travailler comme migrant), 外流 \emph{wàiliú} (exode), 机会 \emph{jīhuì} (opportunité). Astuce : certains mots peuvent appartenir à deux colonnes selon l'argument ; à l'oral, c'est votre argumentation qui compte.
\end{corrige}

\courssub{Carte mentale récapitulative de la leçon}

\begin{popfigure}
\begin{tikzpicture}[
  mindmap,
  grow cyclic,
  every node/.style={concept, align=center, text=white, font=\small},
  root concept/.append style={concept color=popdark, font=\small\bfseries},
  level 1 concept/.append style={sibling angle=90, level distance=3.4cm}
]
\node[root concept, align=center]{农村人口外流\\ nóngcūn rénkǒu wàiliú\\ \emph{l'exode rural}}
  child[concept color=popblue]{ node[align=center]{词汇\\ vocabulaire\\ 城市 / 农村\\ 机会 / 环境} }
  child[concept color=poporange]{ node[align=center]{句型\\ formules\\ 我认为...\\ 我同意...} }
  child[concept color=popgreen]{ node[align=center]{声调\\ les tons\\ lent et juste\\ mieux que vite} }
  child[concept color=poppurple]{ node[align=center]{辩论\\ le débat\\ 你呢？\\ 你说得对。} };
\end{tikzpicture}
\legende{Carte mentale récapitulative : les quatre piliers de la leçon sur l'exode rural.}
\end{popfigure}

\courssub{Synthèse des formules fonctionnelles}

Voici le tableau de référence à garder sous les yeux pendant votre exposé et vos débats. Apprenez au moins \textbf{une formule par fonction} par cœur.

\begin{center}
\begin{tabularx}{\linewidth}{|X|X|X|}
\hline
\rowcolor{popblueL} Fonction & Formule (caractères + pinyin) & Traduction \\
\hline
Présenter & 大家好，我今天讲…… \emph{Dàjiā hǎo, wǒ jīntiān jiǎng...} & Bonjour à tous, aujourd'hui je parle de... \\
\hline
Opiner & 我认为…… \emph{Wǒ rènwéi...} & Je pense que... \\
\hline
Opiner & 我觉得…… \emph{Wǒ juéde...} & Je trouve que... \\
\hline
Réagir (accord) & 我同意你的看法。 \emph{Wǒ tóngyì nǐ de kànfǎ.} & Je suis d'accord avec ton point de vue. \\
\hline
Réagir (désaccord) & 我不太同意。 \emph{Wǒ bú tài tóngyì.} & Je ne suis pas tout à fait d'accord. \\
\hline
Relancer & 你呢？ \emph{Nǐ ne ?} & Et toi ? \\
\hline
Relancer & 你怎么看？ \emph{Nǐ zěnme kàn ?} & Qu'en penses-tu ? \\
\hline
Conclure & 我的讲话完了，谢谢大家！ \emph{Wǒ de jiǎnghuà wán le, xièxie dàjiā !} & Mon exposé est terminé, merci à tous ! \\
\hline
\end{tabularx}
\end{center}

\begin{exemplebox}[La formule de conclusion modèle]
\zh{我的讲话完了，谢谢大家！}\\
\emph{Wǒ de jiǎnghuà wán le, xièxie dàjiā !}\\
« Mon exposé est terminé, merci à tous ! »\\
C'est LA phrase de fin à utiliser à l'examen. Notez le \zh{了} qui marque l'accomplissement : l'exposé est \emph{fini}. Ne dites jamais \zh{我说完了没有} : c'est incorrect (la négation de \zh{了} se fait avec \zh{没有} seul, sans \zh{了} : \zh{我还没说完。} « Je n'ai pas encore fini. »).
\end{exemplebox}

\begin{propriete}[La structure valorisée : 各有各的 + nom]
\textbf{Structure} : A 和 B 各有各的 + nom.\\
\textbf{Sens} : « chacun a ses propres... », « chacun a ses... à lui ».\\
\textbf{Emploi} : idéale pour conclure un débat équilibré, sans prendre parti de façon brutale. C'est une structure de niveau supérieur, très appréciée à l'examen.\\[4pt]
\zh{城市和农村各有各的好处。}\\
\emph{Chéngshì hé nóngcūn gè yǒu gè de hǎochù.}\\
« La ville et la campagne ont chacune leurs avantages. »\\[4pt]
\zh{他们俩各有各的想法。}\\
\emph{Tāmen liǎ gè yǒu gè de xiǎngfǎ.}\\
« Ils ont chacun leur propre opinion. »\\[4pt]
\textbf{Attention} : le nom après \zh{的} doit être un nom simple (好处, 想法, 特点), pas une phrase entière. On ne dit pas \zh{各有各的很好}.
\end{propriete}

\begin{attention}[还是 ou 或者 ? Le « ou » chinois a deux visages]
Le français « ou » se traduit de deux façons différentes :
\begin{itemize}
\item \zh{还是} \emph{háishi} dans les \textbf{questions} : \zh{城市好还是农村好？} « La ville ou la campagne : qu'est-ce qui est mieux ? »
\item \zh{或者} \emph{huòzhě} dans les \textbf{phrases affirmatives} : \zh{我想住在城市或者农村。} « Je veux habiter en ville ou à la campagne. »
\end{itemize}
Erreur fréquente des francophones : utiliser \zh{或者} dans une question. Retenez : question → \zh{还是} ; affirmation → \zh{或者}.
\end{attention}

\courssub{Consigne de l'exposé évalué}

\begin{definition}[Sujet de l'exposé]
Vous préparerez un exposé oral d'\textbf{une minute} au choix :
\begin{itemize}
\item \zh{我的家乡} \emph{Wǒ de jiāxiāng} — « Mon pays natal » (décrire votre ville ou village, dire si les jeunes y restent ou partent) ;
\item \zh{城市好还是农村好？} \emph{Chéngshì hǎo háishi nóngcūn hǎo ?} — « La ville ou la campagne : qu'est-ce qui est mieux ? » (donner votre avis argumenté).
\end{itemize}
\end{definition}

\begin{methode}[Fiche de préparation d'exposé : la structure en 5 étapes]
\begin{enumerate}
\item \textbf{Salutation} : \zh{大家好！} \emph{Dàjiā hǎo !} — « Bonjour à tous ! »
\item \textbf{Annonce du sujet} : \zh{我今天讲我的家乡。} \emph{Wǒ jīntiān jiǎng wǒ de jiāxiāng.} — « Aujourd'hui je parle de mon pays natal. »
\item \textbf{Trois arguments} : utilisez \zh{第一……第二……第三……} \emph{dì-yī... dì-èr... dì-sān...} — « premièrement... deuxièmement... troisièmement... ». Exemple : \zh{第一，农村的空气很好。} \emph{Dì-yī, nóngcūn de kōngqì hěn hǎo.} « Premièrement, l'air de la campagne est bon. »
\item \textbf{Avis personnel} : \zh{我认为……} ou \zh{我觉得……}, éventuellement avec la structure \zh{各有各的}.
\item \textbf{Remerciement} : \zh{我的讲话完了，谢谢大家！}
\end{enumerate}
\end{methode}

\begin{aretenir}[Critères d'évaluation de l'exposé]
Votre note reposera sur quatre critères :
\begin{itemize}
\item \textbf{Tons corrects} : la prononciation des tons est le premier critère ;
\item \textbf{Au moins 3 formules fonctionnelles} du tableau (présenter, opiner, conclure...) ;
\item \textbf{Structure en 5 étapes} respectée (salutation → sujet → arguments → avis → remerciement) ;
\item \textbf{Durée} : environ une minute, sans lire intégralement un texte.
\end{itemize}
\end{aretenir}

\begin{attention}[Parler lentement et juste vaut mieux que parler vite et faux]
L'erreur classique des francophones à l'oral : vouloir parler vite pour « faire chinois ». Résultat : les tons disparaissent et \zh{mǎ} (cheval) devient \zh{mā} (maman) ! Un évaluateur préfère toujours un débit \textbf{lent, posé, avec des tons nets}. Respirez, articulez chaque syllabe, marquez le 4\textsuperscript{e} ton en descendant franchement et le 3\textsuperscript{e} ton en creusant la voix. La clarté est une compétence, pas une faiblesse.
\end{attention}

\courssub{Gagner du temps à l'oral : les mots de secours}

Même les meilleurs orateurs ont des trous de mémoire. En chinois, il existe des \cle{mots de remplissage} naturels qui vous permettent de réfléchir sans rester muet :

\begin{center}
\begin{tabularx}{\linewidth}{|X|X|X|}
\hline
\rowcolor{popblueL} Expression & Pinyin & Équivalent français \\
\hline
这个…… & zhège... & « euh... », « ce truc... » \\
\hline
那个…… & nàge... & « euh... », « comment dire... » \\
\hline
怎么说？ & zěnme shuō ? & « comment dire ? » (se relancer soi-même) \\
\hline
就是说…… & jiùshì shuō... & « c'est-à-dire... » \\
\hline
\end{tabularx}
\end{center}

\begin{exemplebox}[Un trou de mémoire bien géré]
\zh{我觉得农村很……那个……很安静。}\\
\emph{Wǒ juéde nóngcūn hěn... nàge... hěn ānjìng.}\\
« Je trouve que la campagne est très... euh... très calme. »\\[4pt]
\zh{很多年轻人去城市打工，这个……怎么说？……为了找到好工作。}\\
\emph{Hěn duō niánqīngrén qù chéngshì dǎgōng, zhège... zěnme shuō ? ... wèile zhǎodào hǎo gōngzuò.}\\
« Beaucoup de jeunes vont travailler en ville, euh... comment dire ?... pour trouver un bon travail. »\\
Notez \zh{找到} \emph{zhǎodào} : \zh{找} (chercher) + \zh{到} (résultat atteint) = « trouver ». Le verbe de résultat rend la phrase plus naturelle.
\end{exemplebox}

\begin{attention}[Ne pas abuser des mots de remplissage]
\zh{这个} et \zh{那个} sont des béquilles : une ou deux fois dans un exposé d'une minute, c'est naturel ; à chaque phrase, cela devient un tic pénalisant. Préférez une courte pause silencieuse si vous hésitez longuement.
\end{attention}

\begin{savaistu}[Ce que l'évaluateur note vraiment]
À l'examen d'expression orale, l'évaluateur n'écoute pas seulement votre exposé préparé : il note surtout votre capacité à \textbf{réagir spontanément}. Après votre minute, il pourra vous poser une question simple comme \zh{你喜欢农村吗？} \emph{Nǐ xǐhuan nóngcūn ma ?} « Aimes-tu la campagne ? ». Gardez donc sous la main deux phrases magiques : \zh{你呢？} \emph{Nǐ ne ?} « Et toi ? » pour relancer, et \zh{你说得对。} \emph{Nǐ shuō de duì.} « Tu as raison. » pour réagir avec élégance. Un candidat qui réagit naturellement obtient toujours une meilleure note qu'un candidat qui récite parfaitement mais reste bloqué à la première question.
\end{savaistu}

\courssub{Modèle complet d'exposé}

\begin{textebox}[Exposé modèle : 我的家乡]
大家好！我今天讲我的家乡。\\
\emph{Dàjiā hǎo ! Wǒ jīntiān jiǎng wǒ de jiāxiāng.}\\
« Bonjour à tous ! Aujourd'hui je parle de mon pays natal. »\\[4pt]
我的家乡是雅温得附近的一个小村子。那里空气很好，环境很安静。\\
\emph{Wǒ de jiāxiāng shì Yǎwēndé fùjìn de yí ge xiǎo cūnzi. Nàlǐ kōngqì hěn hǎo, huánjìng hěn ānjìng.}\\
« Mon village natal est un petit village près de Yaoundé. Là-bas, l'air est bon et l'environnement est calme. »\\[4pt]
第一，农村的生活很舒服；第二，人们很友好；第三，吃的东西很新鲜。\\
\emph{Dì-yī, nóngcūn de shēnghuó hěn shūfu ; dì-èr, rénmen hěn yǒuhǎo ; dì-sān, chī de dōngxi hěn xīnxiān.}\\
« Premièrement, la vie à la campagne est agréable ; deuxièmement, les gens sont amicaux ; troisièmement, la nourriture est fraîche. »\\[4pt]
可是，很多年轻人去城市打工，因为城市有更多的工作机会。\\
\emph{Kěshì, hěn duō niánqīngrén qù chéngshì dǎgōng, yīnwèi chéngshì yǒu gèng duō de gōngzuò jīhuì.}\\
« Cependant, beaucoup de jeunes partent travailler en ville, car la ville offre plus d'opportunités de travail. »\\[4pt]
我认为城市和农村各有各的好处。我的讲话完了，谢谢大家！\\
\emph{Wǒ rènwéi chéngshì hé nóngcūn gè yǒu gè de hǎochù. Wǒ de jiǎnghuà wán le, xièxie dàjiā !}\\
« Je pense que la ville et la campagne ont chacune leurs avantages. Mon exposé est terminé, merci à tous ! »
\end{textebox}

\textbf{Glossaire} : 附近 \emph{fùjìn} (à proximité) ; 村子 \emph{cūnzi} (village) ; 舒服 \emph{shūfu} (confortable, agréable) ; 友好 \emph{yǒuhǎo} (amical) ; 新鲜 \emph{xīnxiān} (frais) ; 可是 \emph{kěshì} (mais, cependant) ; 为了 \emph{wèile} (pour, afin de) ; 更 \emph{gèng} (davantage, plus).

\textbf{Questions de compréhension} :
\begin{enumerate}
\item Où se trouve le village natal du locuteur ?
\item Quels sont les trois arguments donnés en faveur de la campagne ?
\item Pourquoi les jeunes partent-ils en ville selon ce texte ?
\end{enumerate}

\begin{corrige}[Questions de compréhension]
\begin{enumerate}
\item C'est un petit village près de Yaoundé : \zh{雅温得附近的一个小村子。}
\item La vie y est agréable (\zh{生活很舒服}), les gens sont amicaux (\zh{人们很友好}) et la nourriture est fraîche (\zh{吃的东西很新鲜}).
\item Parce que la ville offre plus d'opportunités de travail : \zh{因为城市有更多的工作机会。}
\end{enumerate}
\end{corrige}

\courssub{Entraînement en binôme avec feedback croisé}

\begin{experience}[Feedback croisé en binôme]
\begin{enumerate}
\item Préparez individuellement votre exposé d'une minute (5 minutes de préparation, plan en 5 étapes sur une fiche).
\item L'élève A présente son exposé ; l'élève B écoute et remplit la grille : tons corrects (oui/non), nombre de formules fonctionnelles entendues, structure en 5 étapes respectée (oui/non).
\item L'élève B donne un feedback en chinois si possible : \zh{你说得很好！} \emph{Nǐ shuō de hěn hǎo !} « Tu parles très bien ! » ou \zh{我觉得你可以说得慢一点。} \emph{Wǒ juéde nǐ kěyǐ shuō de màn yìdiǎn.} « Je pense que tu peux parler plus lentement. » (notez le \zh{得} qui introduit le complément de manière : \zh{说得慢一点} « parler plus lentement »).
\item Échangez les rôles. Puis l'élève B pose une question surprise à l'élève A pour tester sa réaction spontanée.
\item Recommencez une fois en corrigeant le point faible signalé par votre binôme.
\end{enumerate}
\end{experience}

\begin{exoresolu}[Construire un exposé en 5 étapes]
\textbf{Énoncé} : À partir des éléments suivants, construisez un exposé complet respectant la structure en 5 étapes : sujet = \zh{城市好还是农村好} ; arguments = ville : transports pratiques, hôpitaux, écoles ; avis = chacun a ses avantages.
\tcblower
\textbf{Solution détaillée} :\\
Étape 1 (salutation) : \zh{大家好！} \emph{Dàjiā hǎo !}\\
Étape 2 (sujet) : \zh{我今天讲：城市好还是农村好？} \emph{Wǒ jīntiān jiǎng : chéngshì hǎo háishi nóngcūn hǎo ?}\\
Étape 3 (arguments) : \zh{第一，城市的交通很方便；第二，城市有大医院；第三，城市有很多好学校。} \emph{Dì-yī, chéngshì de jiāotōng hěn fāngbiàn ; dì-èr, chéngshì yǒu dà yīyuàn ; dì-sān, chéngshì yǒu hěn duō hǎo xuéxiào.}\\
Étape 4 (avis) : \zh{我认为城市和农村各有各的好处。} \emph{Wǒ rènwéi chéngshì hé nóngcūn gè yǒu gè de hǎochù.}\\
Étape 5 (remerciement) : \zh{我的讲话完了，谢谢大家！} \emph{Wǒ de jiǎnghuà wán le, xièxie dàjiā !}\\
Vérification : 3 formules fonctionnelles (\zh{我今天讲}, \zh{我认为}, \zh{谢谢大家}), structure complète, tons à soigner sur \zh{hǎo} (3\textsuperscript{e} ton) et \zh{háishi} (4\textsuperscript{e} ton puis ton léger).
\end{exoresolu}

\begin{exercice}[À vous de jouer]
Rédigez puis présentez à votre binôme un exposé d'une minute sur \zh{我的家乡} en utilisant : la structure en 5 étapes, au moins 3 formules fonctionnelles, la structure \zh{各有各的}, et un mot de secours (\zh{这个} ou \zh{怎么说？}) utilisé naturellement une seule fois.
\end{exercice}

\begin{aretenir}[L'essentiel avant l'évaluation]
\begin{itemize}
\item Vocabulaire classé en trois familles : ville, campagne, migration ;
\item Une formule par fonction mémorisée : présenter, opiner, réagir, relancer, conclure ;
\item Structure en 5 étapes : salutation → sujet → 3 arguments → avis → remerciement ;
\item Parler lentement avec des tons nets ; utiliser \zh{各有各的} pour briller ;
\item Question avec « ou » = \zh{还是} ; affirmation avec « ou » = \zh{或者} ;
\item En cas de trou : \zh{这个……} ou \zh{怎么说？}, puis repartir calmement.
\end{itemize}
\end{aretenir}

\newpage

\coursec{Activité d'intégration}

\courssub{Objectifs de l'activité}

À l'issue de cette activité d'intégration, tu seras capable de :
\begin{itemize}
\item participer à un débat oral structuré en chinois sur le thème de l'exode rural ;
\item présenter ton opinion avec les formules \zh{我觉得} et \zh{我认为}, la justifier avec \zh{因为} et réagir à un avis contraire avec \zh{我不同意} ;
\item relancer la parole d'un camarade avec \zh{你呢？} et \zh{你觉得怎么样？} ;
\item conclure un échange avec \zh{总之} ;
\item prononcer correctement les tons dans les mots clés du thème (\zh{城市} chéngshì, \zh{农村} nóngcūn, \zh{工作} gōngzuò, \zh{离开} líkāi, \zh{机会} jīhuì, \zh{好处} hǎochù) ;
\item écouter activement tes camarades et repérer, sur une grille, les formules et les tons utilisés ;
\item rédiger un brouillon structuré pour ton propre mini-exposé évalué.
\end{itemize}

\courssub{Situation}

Le club de chinois de ton lycée de Yaoundé a été invité à simuler une émission de débat télévisé, comme celles de la CRTV ou de la CCTV. Le sujet du jour concerne tout le Cameroun : chaque année, beaucoup de jeunes quittent leur village pour aller étudier ou travailler à Douala ou à Yaoundé. Est-ce une bonne chose ? Faut-il quitter le village ?

La classe est organisée en plateau de télévision : un animateur pose la grande question, deux équipes d'invités défendent leur point de vue, un journaliste interroge le public, et un rapporteur conclut l'émission. Voici le texte d'ouverture lu par l'animatrice :

\begin{textebox}[Ouverture de l'émission — 开场白 (kāichǎngbái)]
大家好！欢迎收看今天的节目。\\
Dàjiā hǎo! Huānyíng shōukàn jīntiān de jiémù.\\
« Bonjour à tous ! Bienvenue dans l'émission d'aujourd'hui. »\\[4pt]
今天的问题是：城市好还是农村好？\\
Jīntiān de wèntí shì: chéngshì hǎo háishi nóngcūn hǎo?\\
« La question du jour est : la ville, c'est mieux, ou la campagne ? »\\[4pt]
很多年轻人离开农村，去城市学习和工作。\\
Hěn duō niánqīngrén líkāi nóngcūn, qù chéngshì xuéxí hé gōngzuò.\\
« Beaucoup de jeunes quittent la campagne pour aller étudier et travailler en ville. »\\[4pt]
你觉得怎么样？请你谈谈你的看法。\\
Nǐ juéde zěnmeyàng? Qǐng nǐ tántan nǐ de kànfǎ.\\
« Qu'en penses-tu ? Donne-nous ton avis, s'il te plaît. »
\end{textebox}

\courssub{Consignes}

\begin{enumerate}
\item \textbf{Préparation du lexique (5 min).} Relis le tableau de vocabulaire ci-dessous. À voix haute, prononce chaque mot en insistant sur les tons : \zh{城市} (2\textsuperscript{e} + 4\textsuperscript{e} ton), \zh{农村} (2\textsuperscript{e} + 1\textsuperscript{er} ton), \zh{离开} (2\textsuperscript{e} + 1\textsuperscript{er} ton), \zh{工作} (1\textsuperscript{er} + 4\textsuperscript{e} ton), \zh{机会} (1\textsuperscript{er} + 4\textsuperscript{e} ton), \zh{好处} (3\textsuperscript{e} + 4\textsuperscript{e} ton). Ton voisin vérifie ta prononciation à l'aide de la fiche « Tons ».
\item \textbf{Formation des rôles (5 min).} La classe se divise en cinq groupes : l'animateur ou l'animatrice, l'équipe « jeunes ruraux » (\zh{农村队}), l'équipe « jeunes urbains » (\zh{城市队}), le journaliste, le rapporteur. Chaque équipe prépare deux arguments avec la structure \zh{我觉得……因为……} ou \zh{我认为……因为……}.
\item \textbf{Préparation des arguments (10 min).} Dans ton équipe, rédige au brouillon deux phrases d'opinion complètes en caractères, pinyin et français. Utilise le vocabulaire du tableau et les connecteurs \zh{也} (aussi), \zh{但是} (mais). Exemple de trame : \zh{我认为农村很好，因为空气很新鲜，也很安静。} (Wǒ rènwéi nóngcūn hěn hǎo, yīnwèi kōngqì hěn xīnxiān, yě hěn ānjìng.)
\item \textbf{Débat (20 min).} L'animateur pose la question \zh{城市好还是农村好？}. Chaque équipe présente ses arguments. Pour réagir, utilise \zh{我不同意，因为……}. Le journaliste interviewe deux ou trois élèves du public avec \zh{你呢？你觉得怎么样？}. Chaque élève doit parler au moins deux fois.
\item \textbf{Écoute active (pendant le débat).} Les élèves qui ne parlent pas remplissent la \textbf{grille d'écoute} (fournie ci-dessous) : coche chaque formule entendue et note un mot dont les tons ont été bien prononcés.
\item \textbf{Conclusion (5 min).} Le rapporteur résume le débat en commençant par \zh{总之} (zǒngzhī, « en résumé ») et donne la conclusion de la classe.
\end{enumerate}

\courssub{Ressources à mobiliser}

\textbf{Vocabulaire clé (HSK 2-3)}

\begin{center}\begin{tabularx}{\linewidth}{|X|X|X|X|}\hline
\rowcolor{popblueL} Caractères & Pinyin & Nature & Traduction \\ \hline
城市 & chéngshì & nom & ville \\ \hline
农村 & nóngcūn & nom & campagne, village \\ \hline
离开 & líkāi & verbe & quitter \\ \hline
工作 & gōngzuò & verbe / nom & travailler ; travail, emploi \\ \hline
学习 & xuéxí & verbe & étudier, apprendre \\ \hline
年轻人 & niánqīngrén & nom & jeunes (personnes) \\ \hline
机会 & jīhuì & nom & occasion, opportunité \\ \hline
空气 & kōngqì & nom & air \\ \hline
方便 & fāngbiàn & adjectif & pratique, commode \\ \hline
安静 & ānjìng & adjectif & calme, tranquille \\ \hline
生活 & shēnghuó & nom / verbe & vie, quotidien ; vivre \\ \hline
好处 & hǎochù & nom & avantage, côté positif \\ \hline
回 & huí & verbe & retourner, revenir (chez soi) \\ \hline
两 & liǎng & numéral & deux (devant classif.) \\ \hline
地方 & dìfang & nom & endroit, lieu \\ \hline
\end{tabularx}\end{center}

\textbf{Formules fonctionnelles pour le débat}

\begin{center}\begin{tabularx}{\linewidth}{|X|X|X|}\hline
\rowcolor{popblueL} Fonction & Formule (caractères + pinyin) & Traduction \\ \hline
Donner son avis (standard) & 我觉得…… (Wǒ juéde...) & Je trouve que... / J'ai l'impression que... \\ \hline
Donner son avis (soutenu) & 我认为…… (Wǒ rènwéi...) & Je pense que... / J'estime que... \\ \hline
Justifier (cause) & ……因为…… (...yīnwèi...) & ...parce que... \\ \hline
Ajouter un argument & 也 (yě) & aussi, également \\ \hline
Nuancer / Opposer & 但是 (dànshì) & mais \\ \hline
Réagir (désaccord) & 我不同意 (Wǒ bù tóngyì) & Je ne suis pas d'accord \\ \hline
Relancer (retour) & 你呢？ (Nǐ ne?) & Et toi ? \\ \hline
Demander l'avis & 你觉得怎么样？ (Nǐ juéde zěnmeyàng?) & Qu'en penses-tu ? / Comment trouves-tu ça ? \\ \hline
Conclure & 总之…… (Zǒngzhī...) & En résumé... / Bref... \\ \hline
\end{tabularx}\end{center}

\begin{propriete}[La question alternative avec 还是]
Structure : \cle{Sujet + A 还是 B？} (A \emph{háishi} B) pour proposer un choix entre deux options dans une question.\\
\zh{城市好还是农村好？} \emph{Chéngshì hǎo háishi nóngcūn hǎo?} « La ville c'est bien, ou la campagne ? »\\
\zh{你去城市还是农村？} \emph{Nǐ qù chéngshì háishi nóngcūn?} « Tu vas en ville ou à la campagne ? »\\
\cle{Règle d'or :} On n'utilise \textbf{jamais} \zh{吗} (ma) dans une question contenant \zh{还是}.\\
\cle{Distinction :} \zh{或者} (huòzhě) s'utilise dans les phrases affirmatives ou négatives (« ou bien »), \zh{还是} dans les questions.
\end{propriete}

\begin{attention}[Les tons des mots du débat — Pièges pour francophones]
\begin{itemize}
\item \zh{城市} \emph{chéngshì} : 2\textsuperscript{e} ton montant (\emph{chéng}) + 4\textsuperscript{e} ton descendant (\emph{shì}). \textbf{Piège :} ne pas dire \emph{chēngshì} (1\textsuperscript{e} ton) ni \emph{chéngshí} (2\textsuperscript{e} ton sur shi).
\item \zh{农村} \emph{nóngcūn} : 2\textsuperscript{e} ton (\emph{nóng}) + 1\textsuperscript{er} ton haut et plat (\emph{cūn}). \textbf{Piège :} \emph{cūn} ne doit pas descendre.
\item \zh{工作} \emph{gōngzuò} : 1\textsuperscript{er} ton (\emph{gōng}) + 4\textsuperscript{e} ton (\emph{zuò}). \textbf{Piège :} \emph{zuò} (travail) $\neq$ \emph{zuǒ} (gauche, 3\textsuperscript{e} ton).
\item \zh{离开} \emph{líkāi} : 2\textsuperscript{e} ton + 1\textsuperscript{er} ton. \textbf{Piège :} ne pas confondre \emph{kāi} (ouvrir/partir, 1\textsuperscript{er} ton) et \emph{kǎi} (triomphant, 3\textsuperscript{e} ton).
\item \zh{好处} \emph{hǎochù} : 3\textsuperscript{e} ton (\emph{hǎo}) + 4\textsuperscript{e} ton (\emph{chù}). \textbf{Sandhi :} devant un 4\textsuperscript{e} ton, \emph{hǎo} se prononce souvent demi-3\textsuperscript{e} ton (bas) ou 2\textsuperscript{e} ton (\emph{háo}) à l'oral rapide.
\item \zh{两} \emph{liǎng} (3\textsuperscript{e} ton) + classif. \zh{个} \emph{ge} (ton neutre). \textbf{Rappel :} On écrit \zh{两个}, jamais \zh{二个}.
\end{itemize}
\end{attention}

\begin{savaistu}[Exode rural : regards croisés Cameroun — Chine]
\textbf{Cameroun :} L'exode rural massif vers Yaoundé et Douala (mais aussi Bafoussam, Garoua) concerne surtout les 15-35 ans. Causes : recherche d'emploi (secteur informel), études secondaires/supérieures, accès aux soins. Conséquences : vieillissement des villages, pression urbaine (logement, transport), mais aussi envois de fonds (\emph{transferts d'argent}) qui soutiennent les familles restées au village. Politiques : programmes « Retour à la terre », incubation agricole (ex. PEA-Jeunes).

\textbf{Chine :} Plus grande migration interne de l'histoire (\emph{农民工} \emph{nóngmíngōng}, « ouvriers paysans ») depuis les années 1980 : 290 millions de migrants (chiffres ~2020). Système \emph{户口} (\emph{hùkǒu}) : registre de résidence limitant l'accès aux services publics (école, hôpital) en ville pour les ruraux. Réformes récentes (Nouvelle urbanisation) pour accorder la citoyenneté urbaine aux migrants. Similitude : quête de meilleurs revenus/éducation. Différence : échelle, rôle de l'État, système \emph{hùkǒu} spécifique.
\end{savaistu}

\begin{methode}[Préparer et réussir son exposé oral (évalué)]
\begin{enumerate}
\item \textbf{Structurer :} Introduction (sujet + accroche) $\rightarrow$ 2 arguments (\zh{首先}..., \zh{其次}...) + justifications (\zh{因为}...) $\rightarrow$ Conclusion (\zh{总之}...).
\item \textbf{Lexique \& Tons :} Sélectionne 8-10 mots-clés. Vérifie chaque ton sur dictionnaire (Pleco, Hanping). Entraîne les paires tonales difficiles (\emph{chéng-shì}, \emph{nóng-cūn}, \emph{hǎo-chù}).
\item \textbf{Connecteurs :} Utilise \zh{因为} (cause), \zh{所以} (conséquence), \zh{也} (addition), \zh{但是} (concession).
\item \textbf{Répétition :} Lis ton texte 3 fois : 1. lentement (tons exacts), 2. vitesse normale (fluidité), 3. sans papier (regard vers public).
\item \textbf{Posture :} Debout, regard balayant la classe, voix projetée, pas de lecture mot à mot.
\end{enumerate}
\end{methode}

\textbf{Grille d'écoute active (à photocopier ou recopier)}

\begin{center}\begin{tabularx}{\linewidth}{|X|c|X|X|}\hline
\rowcolor{popblueL} Formule / Élément attendu & Cocher (✓) & Exemple entendu (caractères / pinyin) & Tons bien prononcés (mot noté) \\ \hline
我觉得 (Wǒ juéde) & & & \\ \hline
我认为 (Wǒ rènwéi) & & & \\ \hline
因为 (yīnwèi) & & & \\ \hline
我不同意 (Wǒ bù tóngyì) & & & \\ \hline
你呢？ (Nǐ ne?) & & & \\ \hline
你觉得怎么样？ (Nǐ juéde zěnmeyàng?) & & & \\ \hline
总之 (Zǒngzhī) & & & \\ \hline
Autre connecteur (也, 但是, 所以...) & & & \\ \hline
\end{tabularx}\end{center}

\courssub{Proposition de production et corrigé commenté}

\begin{exoresolu}[Extrait modèle du débat télévisé]
Reconstitue le déroulement du débat à partir des formules de la leçon, puis compare ta production avec le modèle ci-dessous.
\tcblower
\textbf{Animatrice :} 大家好！今天的问题是：城市好还是农村好？\\
\emph{Dàjiā hǎo! Jīntiān de wèntí shì: chéngshì hǎo háishi nóngcūn hǎo?}\\
« Bonjour à tous ! La question du jour : la ville ou la campagne ? »\\[3pt]
\textbf{Équipe ville :} 我觉得城市很好，因为城市有很多工作机会，生活也很方便。\\
\emph{Wǒ juéde chéngshì hěn hǎo, yīnwèi chéngshì yǒu hěn duō gōngzuò jīhuì, shēnghuó yě hěn fāngbiàn.}\\
« Je trouve que la ville est très bien, parce qu'il y a beaucoup d'opportunités de travail et que la vie est pratique. »\\[3pt]
\textbf{Équipe village :} 我不同意。我认为农村很好，因为空气很新鲜，也很安静。\\
\emph{Wǒ bù tóngyì. Wǒ rènwéi nóngcūn hěn hǎo, yīnwèi kōngqì hěn xīnxiān, yě hěn ānjìng.}\\
« Je ne suis pas d'accord. Je pense que la campagne est très bien, parce que l'air est frais et que c'est calme. »\\[3pt]
\textbf{Journaliste (au public) :} 你呢？你觉得怎么样？\\
\emph{Nǐ ne? Nǐ juéde zěnmeyàng?}\\
« Et toi ? Qu'en penses-tu ? »\\[3pt]
\textbf{Élève du public :} 我觉得两个地方都很好。年轻人可以去城市学习，也可以回农村工作。\\
\emph{Wǒ juéde liǎng ge dìfang dōu hěn hǎo. Niánqīngrén kěyǐ qù chéngshì xuéxí, yě kěyǐ huí nóngcūn gōngzuò.}\\
« Je trouve que les deux endroits sont bien. Les jeunes peuvent aller étudier en ville, et aussi revenir travailler au village. »\\[3pt]
\textbf{Rapporteur :} 总之，城市有城市的好处，农村有农村的好处。\\
\emph{Zǒngzhī, chéngshì yǒu chéngshì de hǎochù, nóngcūn yǒu nóngcūn de hǎochù.}\\
« En résumé, la ville a ses avantages, et la campagne a les siens. »\\[6pt]
\textbf{Commentaire des choix de langue :}
\begin{itemize}
\item \zh{我觉得} introduit une opinion personnelle simple, courante à l'oral ; \zh{我认为} marque un avis plus réfléchi, soutenu, adapté au débat formel.
\item La justification suit l'ordre canonique : \cle{Opinion + 因为 + Raison}. En chinois, la cause (\zh{因为}) suit souvent l'opinion, alors qu'en français « parce que » peut précéder ou suivre.
\item \zh{也} (yě) place \textbf{avant} le verbe/adjectif qu'il modifie : \zh{生活也很方便} (la vie \textbf{aussi} est pratique), \zh{也很安静} (\textbf{aussi} calme). Ne pas mettre \zh{也} en fin de phrase.
\item \zh{两个地方都很好} : le numéral \zh{两} (liǎng) s'impose devant classificateur (\zh{个}), jamais \zh{二} (èr). \zh{都} (dōu, « tous/deux ») se place \textbf{après} le sujet et \textbf{avant} le verbe/adjectif.
\item \zh{回农村工作} : \zh{回} (huí) implique un retour vers l'origine (le village natal), nuance culturelle forte en Chine comme au Cameroun.
\item La conclusion utilise la \textbf{structure parallèle} (对仗 \emph{duìzhàng}) : \zh{城市有……，农村有……}. Figure de style prisée à l'oral et à l'écrit, elle montre la maîtrise rhétorique.
\end{itemize}
\end{exoresolu}

\courssub{Bilan de l'activité}

Grâce à ce débat télévisé simulé, tu as mobilisé en situation réelle les trois compétences de la leçon :
\begin{itemize}
\item \textbf{Lexique :} \zh{城市}, \zh{农村}, \zh{离开}, \zh{工作}, \zh{学习}, \zh{机会}, \zh{空气}, \zh{方便}, \zh{安静}, \zh{生活}, \zh{好处}, \zh{回}, \zh{两}, \zh{地方}.
\item \textbf{Grammaire / Fonctionnel :} Question alternative \zh{还是}, opinion \zh{觉得}/\zh{认为}, justification \zh{因为}, désaccord \zh{不同意}, relance \zh{你呢？}/\zh{觉得怎么样？}, conclusion \zh{总之}, classificateur \zh{两个}, structure parallèle.
\item \textbf{Phonologie :} Entraînement ciblé sur les paires tonales 2-4, 2-1, 1-4, 3-4 et sandhi du 3\textsuperscript{e} ton (\zh{好处}).
\item \textbf{Compétence socio-culturelle :} Comparaison des dynamiques migratoires Cameroun/Chine (vocabulaire \zh{农民工}, \zh{户口}).
\end{itemize}
La grille d'écoute t'a appris à écouter activement en chinois : identifier les marqueurs discursifs et contrôler la justesse tonale chez l'autre. Pour préparer l'exposé évalué (« Mon village et ma ville »), reprends tes deux arguments, enrichis-les d'un connecteur \zh{但是} ou \zh{所以}, vérifie les tons de chaque mot clé sur ton dictionnaire et entraîne-toi à les dire à voix haute, seul puis devant un camarade, en respectant la \textbf{Méthode} ci-dessus. La prochaine étape : passage à l'oral individuel noté.

\newpage

\coursec{Méthodes et exercices résolus}

\courssub{Méthodes et exercices résolus}

\begin{methode}[Préparer et structurer un exposé oral simple sur l'exode rural]
\textbf{Objectif :} Construire un exposé clair de 1 à 2 minutes, en respectant la structure \cle{Introduction -- Développement -- Conclusion} et en utilisant le vocabulaire thématique.

\textbf{Étapes :}
\begin{enumerate}
    \item \textbf{Analyser le sujet :} Identifier les mots-clés : \zh{农村} (campagne), \zh{城市} (ville), \zh{外出务工} (aller travailler loin), \zh{留守儿童} (enfants restés au village), \zh{发展} (développement).
    \item \textbf{Rédiger un plan en 3 parties (mots-clés seulement) :}
    \begin{itemize}
        \item \textit{Introduction :} Phénomène général + accroche (chiffre ou observation).
        \item \textit{Développement (2 arguments) :} 1. Causes (revenus, éducation, santé). 2. Conséquences (village vieillissant, pression urbaine).
        \item \textit{Conclusion :} Bilan personnel + ouverture (politiques \zh{乡村振兴} revitalisation rurale).
    \end{itemize}
    \item \textbf{Écrire les phrases-clés en chinois (brouillon) :} Utiliser les structures \zh{因为……所以……} (parce que... donc...), \zh{不仅……而且……} (non seulement... mais aussi...), \zh{随着……} (avec... / au fur et à mesure que...).
    \item \textbf{S'entraîner à l'oral :} Lire 3 fois en se chronométrant. Vérifier les tons des mots nouveaux (\zh{外出} wàichū, \zh{务工} wùgōng, \zh{留守} liúshǒu). S'enregistrer si possible.
\end{enumerate}

\textbf{Structures à retenir :}
\begin{center}
\begin{tabularx}{\linewidth}{|X|X|X|}
\hline
\rowcolor{popblueL} Structure & Exemple (Caractères + Pinyin) & Traduction \\ \hline
\zh{随着 + Nom/Verbe，...} & \zh{随着经济发展，农村变了。} (Suízhe jīngjì fāzhǎn, nóngcūn biàn le.) & Avec le développement économique, la campagne a changé. \\ \hline
\zh{越来越 + Adj/Verbe} & \zh{年轻人越来越多去城市。} (Niánqīng rén yuè lái yuè duō qù chéngshì.) & De plus en plus de jeunes vont en ville. \\ \hline
\zh{一方面……，另一方面……} & \zh{一方面收入高，另一方面想念家。} (Yī fāngmiàn shōurù gāo, lìng yī fāngmiàn xiǎngniàn jiā.) & D'un côté le revenu est haut, de l'autre on regrette la maison. \\ \hline
\end{tabularx}
\end{center}
\end{methode}

\begin{exoresolu}[Rédiger l'introduction et le plan détaillé d'un exposé]
\textbf{Énoncé :} Vous devez présenter un exposé de 90 secondes sur le thème \zh{农村青年外出务工} (Les jeunes ruraux partent travailler en ville). 
\begin{enumerate}
    \item Rédigez l'\textbf{introduction complète} en chinois (3--4 phrases) : accroche, définition du phénomène, annonce du plan.
    \item Proposez un \textbf{plan détaillé} (arguments + connecteurs logiques en chinois) pour la partie développement.
\end{enumerate}

\tcblower
\textbf{Solution détaillée :}

\textbf{1. Introduction rédigée :}
\begin{textebox}[Introduction type exposé]
\zh{大家好。今天我想谈谈农村青年外出务工的现象。} \\
\zh{随着中国经济的快速发展，越来越多的年轻人离开家乡去城市打工。} \\
\zh{这个现象既有好处，也有坏处。下面我从原因和影响两个方面来介绍。} \\
\textit{Pinyin :} Dàjiā hǎo. Jīntiān wǒ xiǎng tántan nóngcūn qīngnián wàichū wùgōng de xiànxiàng. Suízhe Zhōngguó jīngjì de kuàsù fāzhǎn, yuè lái yuè duō de niánqīng rén líkāi jiāxiāng qù chéngshì dǎgōng. Zhège xiànxiàng jì yǒu hǎochù, yě yǒu huàichù. Xiàmiàn wǒ cóng yuányīn hé yǐngxiǎng liǎng gè fāngmiàn lái jièshào. \\
\textit{Traduction :} Bonjour à tous. Aujourd'hui je veux parler du phénomène des jeunes ruraux partant travailler en ville. Avec le développement économique rapide de la Chine, de plus en plus de jeunes quittent leur ville natale pour aller travailler en ville. Ce phénomène a des avantages et des inconvénients. Ci-dessous, je vais l'introduire sous deux aspects : les causes et les impacts.
\end{textebox}

\textbf{Analyse grammaticale :}
\begin{itemize}
    \item \zh{谈谈} (tántan) : Reduplication du verbe = « parler un peu de », adoucit le ton.
    \item \zh{离开 + lieu} (líkāi) : Structure verbe + complément de lieu (quitter).
    \item \zh{既有……也有……} (jì yǒu... yě yǒu...) : Structure corrélative « à la fois... et... ».
    \item \zh{从……两个方面来介绍} (cóng... liǎng gè fāngmiàn lái jièshào) : Locution prépositive « sous deux aspects ».
\end{itemize}

\textbf{2. Plan détaillé (Développement) :}

\begin{center}
\begin{tabularx}{\linewidth}{|X|X|X|}
\hline
\rowcolor{poporangeL} Partie & Argument principal (Chinois + Pinyin) & Connecteurs / Structures utiles \\ \hline
\textbf{Argument 1 : Causes} & \zh{主要原因是收入低，农村工作机会少。} (Zhǔyào yuányīn shì shōurù dī, nóngcūn gōngzuò jīhuì shǎo.) & \zh{因为……所以……} (Cause/Effet) \\
& \zh{而且城市教育、医疗条件更好。} (Érqiě chéngshì jiàoyù, yīliáo tiáojiàn gèng hǎo.) & \zh{而且} (De plus), \zh{比较} (Comparatif) \\ \hline
\textbf{Argument 2 : Conséquences} & \zh{好处：农民收入增加，生活水平提高。} (Hǎochù: Nóngmín shōurù zēngjiā, shēnghuó shuǐpíng tígāo.) & \zh{好处 / 坏处} (Avantages/Inconvénients) \\
& \zh{坏处：留守老人和儿童多，村庄空心化。} (Huàichù: Liúshǒu lǎorén hé értóng duō, cūnzhuāng kōngxīnhuà.) & \zh{导致} (Entraîner), \zh{空心化} (Vidage/Creusement) \\ \hline
\end{tabularx}
\end{center}

\textbf{Conclusion :} \zh{总之，国家正在推行乡村振兴战略，希望农村越来越好。} (Zǒngzhī, guójiā zhèngzài tuīxíng xiāngcūn zhènxīng zhànlüè, xīwàng nóngcūn yuè lái yuè hǎo.) -- « En résumé, le pays met en œuvre la stratégie de revitalisation rurale, on espère que la campagne ira de mieux en mieux. »
\end{exoresolu}

\begin{methode}[Utiliser les formules fonctionnelles pour donner son avis, nuancer et relancer]
\textbf{Objectif :} Maîtriser les « briques » du débat oral : \cle{Exprimer son opinion}, \cle{Être d'accord / Pas d'accord}, \cle{Nuancer}, \cle{Relancer l'interlocuteur}.

\textbf{Étapes :}
\begin{enumerate}
    \item \textbf{Classer les formules par fonction} (voir tableau ci-dessous). Mémoriser 2 formules par catégorie.
    \item \textbf{Attention à l'ordre des mots :} En chinois, l'adverbe de degré (\zh{很}, \zh{非常}, \zh{特别}) se place \textbf{avant} l'adjectif/verbe, jamais à la fin comme en français (« C'est bien \emph{très} » $\rightarrow$ Faux).
    \item \textbf{Intonation :} \zh{吧} (ba) adoucit (suggestion/accord) ; \zh{呢} (ne) relance (question retour) ; \zh{啊} (a) exprime la surprise/réaction.
    \item \textbf{S'entraîner en binôme :} L'un tire une carte « Pour », l'autre « Contre », on débat 1 minute en utilisant \textbf{au moins 3 formules différentes} chacun.
\end{enumerate}

\textbf{Tableau des formules fonctionnelles (Niveau HSK 2--3) :}
\begin{center}
\begin{tabularx}{\linewidth}{|X|X|X|}
\hline
\rowcolor{popgreenL} Fonction & Formules chinoises (Caractères + Pinyin) & Nuance / Contexte \\ \hline
\textbf{Donner son avis} & \zh{我觉得……} (Wǒ juéde...) / \zh{我认为……} (Wǒ rènwéi...) & Standard / Un peu plus formel \\ \hline
\textbf{Accord total} & \zh{我完全同意。} (Wǒ wánquán tóngyì.) / \zh{没错！} (Méi cuò!) & Fort \\ \hline
\textbf{Accord partiel / Nuance} & \zh{有道理，但是……} (Yǒu dàoli, dànshì...) / \zh{一方面……，不过……} (Yī fāngmiàn..., bùguò...) & Indispensable pour le Bac \\ \hline
\textbf{Désaccord poli} & \zh{我怕不太对。} (Wǒ pà bù tài duì.) / \zh{不一定吧？} (Bù yídìng ba?) & « J'ai peur que ce ne soit pas tout à fait ça » / « Pas sûr, non ? » \\ \hline
\textbf{Relancer / Demander avis} & \zh{你怎么看？} (Nǐ zěnme kàn?) / \zh{你觉得呢？} (Nǐ juéde ne?) / \zh{还有什么想法？} (Hái yǒu shénme xiǎngfǎ?) & Ouvert / Fermé / Élargir \\ \hline
\textbf{Demander précision} & \zh{具体说说？} (Jùtǐ shuōshuo?) / \zh{为什么？} (Wèishénme?) & Inciter à développer \\ \hline
\end{tabularx}
\end{center}
\end{methode}

\begin{exoresolu}[Transformer et compléter un échange de débat]
\textbf{Énoncé :} Complétez le dialogue suivant entre deux élèves (A et B) qui débattent sur \zh{城市生活好还是农村生活好} (La vie en ville est-elle meilleure qu'à la campagne ?). Utilisez les formules de la méthode. Chaque trait \_\_\_\_\_ correspond à une phrase complète.

\begin{textebox}[Dialogue à compléter]
\zh{A: 我觉得城市生活更方便，医院、学校都近。\textbf{你怎么看？}} \\
\zh{B: \textbf{我有不同意见。} 城市虽然方便，但是房租太贵，压力很大。} \\
\zh{A: \textbf{你的观点很有道理。} 不过农村年轻人都走了，只剩老人和孩子。} \\
\zh{B: \textbf{这是事实。} 现在国家搞乡村振兴，以后农村会更好。} \\
\zh{A: \textbf{我也这么希望。} 希望如此吧。}
\end{textebox}

\tcblower
\textbf{Solution détaillée :}

\textbf{Proposition de complétion (une version possible parmi d'autres correctes) :}

\begin{textebox}[Dialogue complété]
\zh{A: 我觉得城市生活更方便，医院、学校都近。\textbf{你怎么看？} (Nǐ zěnme kàn?)} \\
\zh{B: \textbf{我有不同意见。} (Wǒ yǒu bùtóng yìjiàn.) 城市虽然方便，但是房租太贵，压力很大。} \\
\zh{A: \textbf{你的观点很有道理。} (Nǐ de guāndiǎn hěn yǒu dàoli.) 不过农村年轻人都走了，只剩老人和孩子。} \\
\zh{B: \textbf{这是事实。} (Zhè shì shìshí.) 现在国家搞乡村振兴，以后农村会更好。} \\
\zh{A: \textbf{我也这么希望。} (Wǒ yě zhème xīwàng.) 希望如此吧。 (Xīwàng rúcǐ ba.)}
\end{textebox}

\textbf{Justifications grammaticales et lexicales :}
\begin{enumerate}
    \item \zh{你怎么看？} : Formule standard pour « Qu'en penses-tu ? » (Relance).
    \item \zh{我有不同意见} : Manière polie de dire « Je ne suis pas d'accord » (Désaccord). \zh{虽然……但是……} (Suīrán... dànshì...) : Concession (Bien que... cependant...). \zh{压力} (yālì) : Pression/Stress (vocabulaire clé).
    \item \zh{很有道理} (hěn yǒu dàoli) : « Très sensé / Tu as raison sur ce point » (Accord partiel / Nuance). \zh{只剩} (zhǐ shèng) : « Ne rester que » (Structure Verbe + Complément de résultat).
    \item \zh{这是事实} : Constat neutre (Accord sur le fait). \zh{搞} (gǎo) : Verbe familier pour « mener/mettre en œuvre » (politique, projet). \zh{乡村振兴} (xiāngcūn zhènxīng) : Terme officiel « Revitalisation rurale ».
    \item \zh{我也这么希望} : Réaction empathique. \zh{希望如此吧} : « Espérons-le » (\zh{吧} adoucit l'incertitude).
\end{enumerate}

\textbf{Point de vigilance :} Ne pas confondre \zh{觉得} (sentir/penser subjectif) et \zh{认为} (juger/estimer, plus intellectuel). À l'oral, \zh{觉得} est plus naturel.
\end{exoresolu}

\begin{methode}[Maîtriser la prononciation des tons et de l'intonation dans le débat]
\textbf{Objectif :} Éviter les erreurs de tons qui changent le sens (ex: \zh{买} mǎi acheter vs \zh{卖} mài vendre) et adopter l'intonation montante pour les questions ouvertes, descendante pour les affirmations.

\textbf{Étapes :}
\begin{enumerate}
    \item \textbf{Réviser les paires minimales du thème :} Travailler les tons 2 (montant) et 4 (descendant) très fréquents dans le vocabulaire « Ville/Campagne ».
    \item \textbf{Règle du \zh{不} (bù) :} Devant un ton 4, \zh{不} devient ton 2 (\zh{不是} bú shì). Devant autres tons (1, 2, 3), reste ton 4 (\zh{不好} bù hǎo, \zh{不想} bù xiǎng).
    \item \textbf{Règle du \zh{一} (yī) :} Seul / Chiffre = T1 (\zh{一} yī). Devant T4 = T2 (\zh{一个} yí ge). Devant T1/2/3 = T4 (\zh{一样} yíyàng).
    \item \textbf{Intonation de phrase :} 
    \begin{itemize}
        \item Affirmation / Ordre : Courbe descendante globale.
        \item Question ouverte (\zh{吗}, \zh{呢}, \zh{怎么}) : Montée finale sur la dernière syllabe.
        \item Énumération : Montée sur chaque élément sauf le dernier (descendant).
    \end{itemize}
    \item \textbf{Exercice « Ombre » (Shadowing) :} Écouter le modèle audio (ou prof), répéter avec 0,5s de décalage, imiter la mélodie exacte.
\end{enumerate}

\textbf{Paires minimales à travailler (Vocabulaire leçon 25) :}
\begin{center}
\begin{tabularx}{\linewidth}{|X|X|X|X|}
\hline
\rowcolor{poppurpleL} Caractères & Pinyin & Sens & Piège tonal \\ \hline
\zh{农村 / 弄船} & nóngcūn / nòngchuán & Campagne / Ramer un bateau & T2+T1 vs T4+T2 \\ \hline
\zh{城市 / 诚实} & chéngshì / chéngshí & Ville / Honnête & T2+T4 vs T2+T2 \\ \hline
\zh{外出 / 歪斜} & wàichū / wāixié & Sortir / De travers & T4+T1 vs T1+T2 \\ \hline
\zh{务工 / 无锅} & wùgōng / wúguō & Travailler (migrant) / Sans marmite & T4+T1 vs T2+T1 \\ \hline
\zh{留守 / 流手} & liúshǒu / liúshǒu & Rester garder / Main qui glisse & \textbf{Mêmes tons !} Contexte seul distingue. \\ \hline
\end{tabularx}
\end{center}
\end{methode}

\begin{exoresolu}[Discrimination auditive et correction phonétique]
\textbf{Énoncé :} 
\begin{enumerate}
    \item \textbf{Dictée tonale :} L'enseignant lit 5 phrases. L'élève note les tons (chiffres 1-4-5) sur les syllabes en gras.
    \item \textbf{Correction :} L'élève lit la phrase corrigée ci-dessous. Identifiez l'erreur de ton dans la lecture fautive de l'élève \textbf{Élève X} et corrigez-la.
\end{enumerate}

\textbf{Phrases du prof (Modèle) :}
1. \zh{农村的环境很好。} (Nóngcūn de huánjìng hěn hǎo.)
2. \zh{他不想去城市打工。} (Tā bù xiǎng qù chéngshì dǎgōng.)
3. \zh{留守儿童需要关心。} (Liúshǒu értóng xūyào guānxīn.)
4. \zh{一个农民工回家了。} (Yí ge nóngmíngōng huí jiā le.)
5. \zh{你觉得乡村振兴有用吗？} (Nǐ juéde xiāngcūn zhènxīng yǒuyòng ma?)

\textbf{Lecture de l'Élève X (avec erreurs) :}
1. \zh{农村} (nòngcūn) \zh{环境} (huánjǐng) \zh{很好} (hěn hào).
2. \zh{他} (tā) \zh{不} (bù) \zh{想} (xiǎng) \zh{去} (qù) \zh{城市} (chéngshí) \zh{打工} (dǎgòng).
3. \zh{留守} (liúshòu) \zh{儿童} (értóng) \zh{需要} (xūyào) \zh{关心} (guānxīn).
4. \zh{一个} (yī ge) \zh{农民工} (nóngmíngōng) \zh{回家} (huí jiā) \zh{了} (le).
5. \zh{你} (nǐ) \zh{觉得} (juéde) \zh{乡村振兴} (xiāngcūn zhènxīng) \zh{有用} (yǒuyòng) \zh{吗} (ma).

\tcblower
\textbf{Solution détaillée :}

\begin{center}
\begin{tabularx}{\linewidth}{|X|X|X|X|}
\hline
\rowcolor{popredL} Phrase & Erreur(s) de l'Élève X & Correction (Pinyin standard) & Règle / Explication \\ \hline
1 & \zh{农村} $\rightarrow$ \textbf{nòngcūn} (T4 T1) \newline \zh{环境} $\rightarrow$ \textbf{huánjǐng} (T2 T3) \newline \zh{很好} $\rightarrow$ \textbf{hěn hào} (T3 T4) & \zh{农村} \textbf{nóngcūn} (T2 T1) \newline \zh{环境} \textbf{huánjìng} (T2 T4) \newline \zh{很好} \textbf{hěn hǎo} (T3 T3 $\rightarrow$ T2 T3) & \zh{农} T2 (montant), pas T4. \newline \zh{境} T4 (descendant), pas T3. \newline \textbf{Règle du 3e ton :} Deux T3 consécutifs $\rightarrow$ le 1er devient T2 (\zh{hén hǎo}). \\ \hline
2 & \zh{城市} $\rightarrow$ \textbf{chéngshí} (T2 T2) \newline \zh{打工} $\rightarrow$ \textbf{dǎgòng} (T3 T4) & \zh{城市} \textbf{chéngshì} (T2 T4) \newline \zh{打工} \textbf{dǎgōng} (T3 T1) & \zh{不想} (bù xiǎng) : \zh{不} devant T3 reste T4 (correct). \newline \textbf{Erreurs réelles :} \zh{市} \zh{shì} (T4) pas \zh{shí} (T2). \zh{工} \zh{gōng} (T1) pas \zh{gòng} (T4). \\ \hline
3 & \zh{留守} $\rightarrow$ \textbf{liúshòu} (T2 T4) & \zh{留守} \textbf{liúshǒu} (T2 T3) & \zh{守} se lit T3 (shǒu) ici (verbe garder), pas T4 (shòu - garder/obéir). Contexte : \zh{留守儿童} (enfants laissés en garde). \\ \hline
4 & \zh{一个} $\rightarrow$ \textbf{yī ge} (T1) & \zh{一个} \textbf{yí ge} (T2 T0) & \textbf{Règle \zh{一} :} Devant classifieur T4 (\zh{ge} ton neutre mais origine T4) $\rightarrow$ \zh{一} devient T2 (\zh{yí}). \\ \hline
5 & \zh{吗} $\rightarrow$ \textbf{ma} (T3 ou T4 audible) & \zh{吗} \textbf{ma} (Ton neutre / léger) & \zh{吗} particule interrogative = \textbf{ton neutre} (court, léger), jamais ton plein. L'intonation montante de la phrase porte sur \zh{用} (yòng). \\ \hline
\end{tabularx}
\end{center}

\textbf{Bilan pour l'élève :} Réviser la règle du \zh{不} (devant T4 $\rightarrow$ T2), la règle du \zh{一} (devant T4 $\rightarrow$ T2), le changement de T3+T3, et le ton neutre des particules (\zh{了}, \zh{吗}, \zh{呢}, \zh{的}, \zh{地}, \zh{得}).
\end{exoresolu}

\begin{methode}[Conduire un jeu de rôle / débat guidé : réagir et relancer]
\textbf{Objectif :} Simuler une situation de communication réelle (ex: réunion de village, émission radio) pour développer la fluidité et la réactivité.

\textbf{Étapes :}
\begin{enumerate}
    \item \textbf{Attribuer les rôles :} \cle{Animateur} (guide, relance, synthétise), \cle{Témoin 1} (Jeune parti en ville), \cle{Témoin 2} (Parent resté au village), \cle{Expert/Autorité} (Maire, ONG, \zh{村干部} cadre de village).
    \item \textbf{Préparer des « cartes arguments » (3 min) :} Chaque rôle note 2 arguments clés + 1 vocabulaire imposé à placer.
    \item \textbf{Déroulement type (10 min) :}
    \begin{enumerate}
        \item Animateur : Lance le sujet (\zh{今天我们讨论……}).
        \item Tour de table : Chaque rôle s'exprime 1 min (\zh{首先，我认为……}).
        \item Débat libre : Animateur relance (\zh{XX，你怎么看？}, \zh{有人不同意吗？}).
        \item Synthèse : Animateur conclut (\zh{总的来说……}).
    \end{enumerate}
    \item \textbf{Auto-évaluation / Pair-évaluation :} Grille critériée : Prononciation (Tons), Vocabulaire (Précision), Structures (Variété), Interaction (Relance), Contenu (Cohérence).
\end{enumerate}

\textbf{Expressions de « gestion de tour de parole » (Turn-taking) :}
\begin{center}
\begin{tabularx}{\linewidth}{|X|X|}
\hline
\rowcolor{poppinkL} Situation & Formule chinoise \\ \hline
Prendre la parole & \zh{请允许我说一下。} (Qǐng yǔnxǔ wǒ shuō yíxià.) / \zh{我想补充一点。} (Wǒ xiǎng bǔchōng yìdiǎn.) \\ \hline
Couper poliment / Revenir & \zh{对不起，打断一下。} (Duìbuqǐ, dǎduàn yíxià.) / \zh{回到刚才的话题……} (Huídào gāngcái de huàtí...) \\ \hline
Gagner du temps / Réfléchir & \zh{这个问题嘛……} (Zhège wèntí ma...) / \zh{让我想想。} (Ràng wǒ xiǎngxiǎng.) \\ \hline
Reformuler pour vérifier & \zh{你的意思是……对吗？} (Nǐ de yìsi shì... duì ma?) \\ \hline
Clore son intervention & \zh{我就说到这里。} (Wǒ jiù shuō dào zhèlǐ.) \\ \hline
\end{tabularx}
\end{center}
\end{methode}

\begin{exoresolu}[Réorganiser et enrichir un débat désordonné]
\textbf{Énoncé :} Les répliques ci-dessous proviennent d'un débat sur \zh{年轻人应该留在农村发展吗？} (Les jeunes doivent-ils rester développer la campagne ?). Elles sont mélangées.
\begin{enumerate}
    \item Remettez-les dans l'ordre logique (Animateur $\rightarrow$ Témoins $\rightarrow$ Débat $\rightarrow$ Conclusion).
    \item Ajoutez \textbf{une formule de relance} de l'animateur entre chaque intervention de témoin.
    \item Ajoutez \textbf{une formule de nuance} dans la bouche du Témoin 2.
\end{enumerate}

\textbf{Répliques brutes :}
A. \zh{我觉得留在农村也好，现在电商发达，卖农产品赚钱。} (Témoin 1)
B. \zh{大家好，今天我们讨论：年轻人该不该留在农村？} (Animateur - Ouverture)
C. \zh{但是网络信号不好，物流也慢，不方便。} (Témoin 2 - Objection brute)
D. \zh{总之，各有优劣，关键看个人选择。谢谢大家。} (Animateur - Conclusion)
E. \zh{我出去打工十年，累得慌，想回老家养鸡。} (Témoin 1 - Expérience)
F. \zh{你觉得呢，老张？} (Animateur - Relance brute)

\tcblower
\textbf{Solution détaillée :}

\textbf{1. Ordre logique reconstitué :} B $\rightarrow$ F (vers Témoin 1) $\rightarrow$ E $\rightarrow$ A $\rightarrow$ [Relance vers Témoin 2] $\rightarrow$ C (modifié) $\rightarrow$ D.

\textbf{2. Dialogue final enrichi (Version corrigée) :}

\begin{textebox}[Débat structuré et enrichi]
\zh{\textbf{Animateur :} 大家好，今天我们讨论：年轻人该不该留在农村？} \\
\zh{\textbf{Animateur :} 先请小李谈谈，\textbf{你怎么看？} (Nǐ zěnme kàn?)} \\
\zh{\textbf{Témoin 1 (Petit Li) :} 我出去打工十年，累得慌，想回老家养鸡。} \\
\zh{\textbf{Témoin 1 :} 我觉得留在农村也好，现在电商发达，卖农产品赚钱。} \\
\zh{\textbf{Animateur :} \textbf{老王，你有什么不同看法？} (Nǐ yǒu shénme bùtóng kànfǎ?)} \\
\zh{\textbf{Témoin 2 (Vieux Wang) :} \textbf{有道理，但是}网络信号不好，物流也慢，不方便。 (Yǒu dàoli, dànshì wǎngluò xìnhào bù hǎo, wùliú yě màn, bù fāngbiàn.)} \\
\zh{\textbf{Animateur :} 确实，基建还要加强。\textbf{大家觉得呢？} (Dàjiā juéde ne?)} \\
\zh{\textbf{Animateur :} 总之，各有优劣，关键看个人选择。谢谢大家。}
\end{textebox}

\textbf{Analyse des ajouts :}
\begin{itemize}
    \item \zh{你怎么看？} / \zh{你有什么不同看法？} / \zh{大家觉得呢？} : Trois formules de relance distinctes (ouvert, ciblé, collectif).
    \item \zh{有道理，但是……} : Formule de nuance (Accord partiel) insérée au début de la réplique C. Transforme une opposition frontale en débat constructif.
    \item \zh{电商} (diànshāng - e-commerce), \zh{农产品} (nóngchǎnpǐn - produits agricoles), \zh{物流} (wùliú - logistique), \zh{基建} (jījàn - infrastructures de base) : Vocabulaire précis HSK 3-4 valorisé.
\end{itemize}
\end{exoresolu}

\begin{methode}[Rédiger une synthèse écrite de son exposé (Expression écrite d'appui)]
\textbf{Objectif :} Transposer l'oral en écrit formel (style \zh{书面语} vs \zh{口语}) pour l'évaluation écrite ou le portfolio. Respecter la structure \cle{Thèse -- Arguments -- Synthèse}.

\textbf{Étapes :}
\begin{enumerate}
    \item \textbf{Passer de l'oral à l'écrit :} Remplacer les marqueurs oraux par des connecteurs écrits.
    \begin{center}
    \begin{tabularx}{\linewidth}{|X|X|}
    \hline
    \rowcolor{popgoldL} Oral (口语) & Écrit (书面语) \\ \hline
    \zh{我觉得 / 我觉得} & \zh{本文认为 / 笔者认为} (Cet article pense / L'auteur pense) \\ \hline
    \zh{因为……所以……} & \zh{由于……因此……} (En raison de... par conséquent...) \\ \hline
    \zh{而且 / 还有} & \zh{此外 / 另外} (De plus / En outre) \\ \hline
    \zh{但是 / 不过} & \zh{然而 / 不过 / 尽管如此} (Cependant / Néanmoins) \\ \hline
    \zh{最后 / 总之} & \zh{综上所述 / 综合来看} (En résumé / Globalement) \\ \hline
    \end{tabularx}
    \end{center}
    \item \textbf{Structurer le paragraphe type « Argument » :} \cle{Phrase sujet} (Idée force) + \cle{Explication} (Cause/Effet) + \cle{Exemple/Donnée} (Concret) + \cle{Mini-conclusion} (Lien thèse).
    \item \textbf{Soigner les caractères :} Écrire les 20 mots-clés du thème 3 fois chacun (ordre des traits). Vérifier les homophones : \zh{外出} vs \zh{外chu}, \zh{务工} vs \zh{无工}, \zh{留守} vs \zh{流守}.
    \item \textbf{Relire :} Check-list : Accords (pas d'accords en chinois !), Particules \zh{的/得/地}, Compléments d'état (\zh{得}), Aspects verbaux (\zh{了/过/着}).
\end{enumerate}
\end{methode}

\begin{exoresolu}[Traduction Thème / Version : Paragraphe argumentatif]
\textbf{Énoncé :} Traduisez le paragraphe suivant en chinois (Thème). Utilisez le style écrit (\zh{书面语}).

\textbf{Texte source (Français) :}
« Avec le développement rapide de l'économie, le phénomène de l'exode rural des jeunes s'intensifie. D'une part, travailler en ville apporte des revenus élevés et élargit les horizons. D'autre part, cela entraîne le vieillissement des villages et le problème des enfants "laissés-pour-compte". Par conséquent, le gouvernement met en œuvre la stratégie de revitalisation rurale pour améliorer les conditions de vie à la campagne. À mon avis, que l'on parte ou que l'on reste, l'essentiel est de contribuer au développement de sa ville natale. »

\tcblower
\textbf{Solution détaillée :}

\textbf{Traduction proposée (Style écrit) :}
\begin{textebox}[Version chinoise - Style écrit]
\zh{随着经济的快速发展，农村青年外出务工的现象日益加剧。} \\
\zh{一方面，进城务工带来了较高收入，开阔了眼界；} \\
\zh{另一方面，这导致村庄老龄化严重，留守儿童问题凸显。} \\
\zh{因此，政府实施乡村振兴战略，旨在改善农村生活条件。} \\
\zh{笔者认为，无论外出还是留守，关键在于为家乡发展贡献力量。}
\end{textebox}

\textbf{Pinyin :}
Suízhe jīngjì de kuàsù fāzhǎn, nóngcūn qīngnián wàichū wùgōng de xiànxiàng rìyì jiājí. Yī fāngmiàn, jìnchéng wùgōng dàilái le jiào gāo shōurù, kāikuò le yǎnjiè; lìng yī fāngmiàn, zhè dǎozhì cūnzhuāng lǎolínghuà yánzhòng, liúshǒu értóng wèntí tūxiǎn. Yīncǐ, zhèngfǔ shíshī xiāngcūn zhènxīng zhànlüè, zhǐzài gǎishàn nóngcūn shēnghuó tiáojiàn. Bǐzhě rènwéi, wúlùn wàichū háishì liúshǒu, guānjiàn zài yú wèi jiāxiāng fāzhǎn gòngxiàn lìliàng.

\textbf{Justifications grammaticales et lexicales (Points clés du Bac) :}
\begin{enumerate}
    \item \zh{随着……发展} (Suízhe... fāzhǎn) : Structure prépositive standard « Avec le développement de... ». Éviter \zh{因为经济发展} (trop oral).
    \item \zh{日益加剧} (rìyì jiājí) : « S'intensifie de jour en jour » (Chengyu / Structure formelle 4 caractères). Remplace \zh{越来越严重} (oral).
    \item \zh{一方面……另一方面……} (Yī fāngmiàn... lìng yī fāngmiàn...) : Structure binaire formelle. Ponctuation : point-virgule ; entre les deux arguments.
    \item \zh{进城务工} (jìnchéng wùgōng) : Verbe composé « aller en ville travailler » (plus concis que \zh{去城市打工}).
    \item \zh{开阔眼界} (kāikuò yǎnjiè) : « Élargir les horizons » (Expression idiomatique). Éviter \zh{看见更多东西}.
    \item \zh{导致……严重 / 凸显} (dǎozhì... yánzhòng / tūxiǎn) : \zh{导致} (entraîner conséquence négative) + adjectif \zh{严重} / verbe \zh{凸显} (mettre en relief / devenir saillant). Vocabulaire précis HSK 4.
    \item \zh{实施战略} (shíshī zhànlüè) : Collocation « mettre en œuvre une stratégie » (verbe \zh{实施} pour politique/plan/loi).
    \item \zh{旨在} (zhǐzài) : « Viser à / avoir pour but » (Formel, suivi d'un verbe d'action). Remplace \zh{想要} / \zh{为了}.
    \item \zh{笔者认为} (bǐzhě rènwéi) : « L'auteur pense » (Impersonnel, style rédactionnel). Remplace \zh{我觉得}.
    \item \zh{无论……还是……} (wúlùn... háishì...) : « Quoi que... / Que... ou... » (Concession). Structure formelle.
    \item \zh{关键在于} (guānjiàn zài yú) : « La clé réside dans ». \zh{在于} (résider dans) + Nom/Verbe.
    \item \zh{贡献力量} (gòngxiàn lìliàng) : « Contribuer par ses forces » (Collocation standard).
\end{enumerate}

\textbf{Erreurs fréquentes à éviter (Pièges Thème) :}
\begin{itemize}
    \item \textbf{Ordre des mots :} \zh{在农村生活条件改善} (Faux : PP avant verbe) $\rightarrow$ \zh{改善农村生活条件} (V + OBJ) ou \zh{在农村改善生活条件} (PP + V + OBJ).
    \item \textbf{Particules :} \zh{严重的问题} (Adj + \zh{的} + N) vs \zh{问题严重} (Sujet + Adj verbe). Ici \zh{老龄化严重} (Sujet + Verbe d'état) $\rightarrow$ PAS de \zh{的}.
    \item \textbf{Aspect \zh{了} :} \zh{带来了} (Action accomplie/résultative) nécessaire car conséquence passée/constatée. \zh{实施了} possible mais \zh{实施} seul suffit (verbe d'action unique passé).
\end{itemize}
\end{exoresolu}

\begin{attention}[Les 5 erreurs qui coûtent des points au Bac (Leçon 25)]
\begin{enumerate}
    \item \textbf{Tons : Confusion \zh{务} (wù, T4) / \zh{无} (wú, T2) et \zh{守} (shǒu, T3) / \zh{手} (shǒu, T3).}
    \textit{Risque :} Dire \zh{无工} (sans travail) au lieu de \zh{务工} (travailler) ; dire \zh{流手} (main qui coule) au lieu de \zh{留守} (rester garder). \textbf{Reflexe :} Associer \zh{务} à \zh{服务} (service T4) et \zh{守} à \zh{守住} (garder T3).
    
    \item \textbf{Caractères confondus (Graphie) : \zh{外} (extérieur) vs \zh{处} (lieu/endroit) ; \zh{振} (secouer/revitaliser) vs \zh{震} (tremblement de terre/se plaindre).}
    \textit{Risque :} Écrire \zh{处出务工} ou \zh{乡村震兴}. \textbf{Reflexe :} \zh{外} = \zh{夕} (soir) + \zh{卜} (divination) $\rightarrow$ dehors le soir. \zh{振} = \zh{振奋} (revigorer) main \zh{扌} qui secoue. \zh{震} = pluie \zh{雨} au-dessus \zh{陳} $\rightarrow$ tonnerre/tremblement.
    
    \item \textbf{Ordre des mots / Place des compléments : \zh{在农村} (PP Lieu) placé APRÈS le verbe.}
    \textit{Erreur :} \zh{我去了在农村工作。} (Français : « Je suis allé À la campagne travailler »).
    \textit{Correct :} \zh{我去农村工作了。} ou \zh{我在农村工作过。} (Chinois : Sujet + Lieu (Préposé) + Verbe). \textbf{Règle d'or :} Complément de lieu \textbf{avant} le verbe d'action (sauf verbes de placement \zh{放/把/住}).
    
    \item \textbf{Mots-outils : \zh{的 / 得 / 地} dans les phrases complexes du débat.}
    \textit{Erreur :} \zh{他说得很流利的中文} (Mélange \zh{得} et \zh{的}). \zh{农村发展得很快的} (Double marque).
    \textit{Correct :} \zh{他说中文说得很流利} (V \zh{得} Complément d'état) / \zh{流利的中文} (Adj \zh{的} N). \zh{农村发展很快} (Sujet + V + Adv) $\rightarrow$ PAS de \zh{得} ni \zh{的}.
    
    \item \textbf{Calque du français : « C'est \emph{moi qui} pars » $\rightarrow$ \zh{是我去} (Maladroit).}
    \textit{Correct :} \zh{我去就行了} (C'est bon si j'y vais) / \zh{由我去吧} (Laisse-moi y aller). 
    \textit{Autre calque :} « Il \emph{a} 20 ans » $\rightarrow$ \zh{他有20岁} (Faux). \textbf{Correct :} \zh{他20岁了} (Âge = Verbe d'état, pas \zh{有}).
\end{enumerate}
\end{attention}

\newpage

\coursec{Exercices}

\courssub{Je vérifie mes connaissances}

\begin{exercice}[QCM : le bon choix]
Pour chaque question, entoure la bonne réponse.
\begin{enumerate}
\item \zh{农村} signifie : a) la ville \quad b) la campagne \quad c) le marché
\item \zh{农民工} désigne : a) un ouvrier agricole migrant \quad b) un professeur \quad c) un étudiant
\item Pour donner ton avis, tu dis : a) \zh{我认为……} \quad b) \zh{对不起……} \quad c) \zh{没关系……}
\item \zh{因为……所以……} exprime : a) le temps \quad b) la cause et la conséquence \quad c) la comparaison
\item Pour relancer la parole dans un débat, tu dis : a) \zh{你呢？} \quad b) \zh{谢谢！} \quad c) \zh{再见！}
\end{enumerate}
\end{exercice}

\begin{exercice}[Vrai ou faux ? Justifie]
Réponds par vrai ou faux, puis justifie ta réponse en une phrase en français.
\begin{enumerate}
\item \zh{城市} (chéngshì) signifie « la campagne ».
\item \zh{我同意你的看法。} (Wǒ tóngyì nǐ de kànfǎ.) veut dire « Je ne suis pas d'accord avec toi. »
\item La structure pour expliquer un choix est : \zh{因为 + cause，所以 + conséquence。}
\item \zh{机会} (jīhuì) signifie « problème ».
\item \zh{怎么样？} (Zěnmeyàng ?) sert à demander l'avis de quelqu'un.
\end{enumerate}
\end{exercice}

\begin{exercice}[Vocabulaire : complète le tableau]
Complète le tableau avec la traduction française des mots suivants, puis écris la nature de chaque mot (nom, verbe, adjectif).

\begin{center}
\begin{tabularx}{\linewidth}{|X|X|X|X|}
\hline
\rowcolor{popblueL} Caractères & Pinyin & Nature & Traduction \\
\hline
\zh{家乡} & jiāxiāng & & \\
\hline
\zh{离开} & líkāi & & \\
\hline
\zh{工作} & gōngzuò & & \\
\hline
\zh{生活} & shēnghuó & & \\
\hline
\zh{方便} & fāngbiàn & & \\
\hline
\zh{希望} & xīwàng & & \\
\hline
\end{tabularx}
\end{center}
\end{exercice}

\begin{exercice}[Pinyin et caractères : associe]
Associe chaque mot en caractères à son pinyin. Attention aux tons !

\begin{center}
\begin{tabularx}{\linewidth}{|X|X|}
\hline
\rowcolor{popblueL} Caractères & Pinyin \\
\hline
1. \zh{农村} & a. gōngzuò \\
\hline
2. \zh{城市} & b. nóngcūn \\
\hline
3. \zh{工作} & c. chéngshì \\
\hline
4. \zh{机会} & d. shēnghuó \\
\hline
5. \zh{生活} & e. jīhuì \\
\hline
6. \zh{问题} & f. wèntí \\
\hline
\end{tabularx}
\end{center}
\end{exercice}

\courssub{Je m'entraîne}

\begin{exercice}[Traduction vers le français]
Traduis les phrases suivantes en français.
\begin{enumerate}
\item \zh{很多年轻人离开农村，去城市找工作。}\\(Hěn duō niánqīngrén líkāi nóngcūn, qù chéngshì zhǎo gōngzuò.)
\item \zh{因为农村的工作机会很少，所以他们去了大城市。}\\(Yīnwèi nóngcūn de gōngzuò jīhuì hěn shǎo, suǒyǐ tāmen qùle dà chéngshì.)
\item \zh{我认为城市的生活很方便，但是也很贵。}\\(Wǒ rènwéi chéngshì de shēnghuó hěn fāngbiàn, dànshì yě hěn guì.)
\item \zh{我希望我的家乡越来越好。}\\(Wǒ xīwàng wǒ de jiāxiāng yuè lái yuè hǎo.)
\end{enumerate}
\end{exercice}

\begin{exercice}[Dialogue à compléter puis à jouer]
Complète le dialogue avec les mots de la liste : \zh{因为}、\zh{所以}、\zh{机会}、\zh{同意}、\zh{怎么样}、\zh{呢}. Puis joue le dialogue à deux, en faisant attention aux tons.

\begin{textebox}[Dialogue : 你为什么离开农村？]
A：你为什么离开农村？\\
(Nǐ wèishénme líkāi nóngcūn ?)\\
B：\_\_\_\_\_\_ 农村的工作 \_\_\_\_\_\_ 很少，\_\_\_\_\_\_ 我来杜阿拉找工作。\\
(\_\_\_\_\_\_ nóngcūn de gōngzuò \_\_\_\_\_\_ hěn shǎo, \_\_\_\_\_\_ wǒ lái Dù'ālā zhǎo gōngzuò.)\\
A：你觉得城市的生活 \_\_\_\_\_\_？\\
(Nǐ juéde chéngshì de shēnghuó \_\_\_\_\_\_ ?)\\
B：很方便，但是太贵了。你 \_\_\_\_\_\_？\\
(Hěn fāngbiàn, dànshì tài guì le. Nǐ \_\_\_\_\_\_ ?)\\
A：我不 \_\_\_\_\_\_。我觉得农村的生活也很好。\\
(Wǒ bù \_\_\_\_\_\_. Wǒ juéde nóngcūn de shēnghuó yě hěn hǎo.)
\end{textebox}
\end{exercice}

\begin{exercice}[Discrimination des tons]
Ton professeur (ou ton camarade) lit un mot de chaque paire. Entoure le mot entendu. Puis lis toi-même chaque paire à voix haute en marquant bien la différence de tons.
\begin{enumerate}
\item \zh{农村} nóngcūn \quad / \quad nòngcūn
\item \zh{工作} gōngzuò \quad / \quad gǒngzuò
\item \zh{城市} chéngshì \quad / \quad chěngshì
\item \zh{机会} jīhuì \quad / \quad jǐhuì
\item \zh{同意} tóngyì \quad / \quad tǒngyì
\item \zh{生活} shēnghuó \quad / \quad shènghuó
\end{enumerate}
\end{exercice}

\begin{exercice}[Réagir et relancer]
\textbf{A.} Réagis à chaque affirmation avec \zh{我同意} ou \zh{我不同意} et justifie ta réponse avec \zh{因为……}.
\begin{enumerate}
\item \zh{城市的生活很舒服。}(Chéngshì de shēnghuó hěn shūfu.)
\item \zh{农村没有工作机会。}(Nóngcūn méiyǒu gōngzuò jīhuì.)
\item \zh{年轻人都应该去大城市。}(Niánqīngrén dōu yīnggāi qù dà chéngshì.)
\end{enumerate}
\textbf{B.} Transforme chaque affirmation en question de relance avec \zh{呢} ou \zh{怎么样}.
\begin{enumerate}
\item \zh{我喜欢城市。} → \zh{你呢？} (modèle)
\item \zh{他觉得农村的生活很好。} → \ldots
\item \zh{我想去杜阿拉工作。} → \ldots
\end{enumerate}
\end{exercice}

\courssub{Je me prépare au Bac}

\begin{exercice}[Compréhension de texte et questions de langue]
Lis le texte suivant, puis réponds aux questions.

\begin{textebox}[从农村到城市]
我叫阿里，二十岁。我的家乡在喀麦隆的北方。三年前，我离开了家乡，来到杜阿拉。因为农村的工作机会很少，所以我想在城市找工作。现在我在一个公司工作，每天工作八个小时。城市的生活很方便，有医院，有学校，也有很多商店。但是城市的生活很贵，房租也很高。我常常想家。我希望以后能回家乡，帮助家乡发展。\\
(Wǒ jiào Ālǐ, èrshí suì. Wǒ de jiāxiāng zài Kāmàilóng de běifāng. Sān nián qián, wǒ líkāile jiāxiāng, láidào Dù'ālā. Yīnwèi nóngcūn de gōngzuò jīhuì hěn shǎo, suǒyǐ wǒ xiǎng zài chéngshì zhǎo gōngzuò. Xiànzài wǒ zài yí ge gōngsī gōngzuò, měi tiān gōngzuò bā ge xiǎoshí. Chéngshì de shēnghuó hěn fāngbiàn, yǒu yīyuàn, yǒu xuéxiào, yě yǒu hěn duō shāngdiàn. Dànshì chéngshì de shēnghuó hěn guì, fángzū yě hěn gāo. Wǒ chángcháng xiǎng jiā. Wǒ xīwàng yǐhòu néng huí jiāxiāng, bāngzhù jiāxiāng fāzhǎn.)
\end{textebox}

\textbf{Questions de compréhension} (réponds en français) :
\begin{enumerate}
\item D'où vient Ali et où habite-t-il maintenant ?
\item Pourquoi a-t-il quitté son village ?
\item Quels sont, selon le texte, les avantages et les inconvénients de la vie en ville ?
\item Quel est le projet d'Ali pour l'avenir ?
\end{enumerate}
\textbf{Questions de langue} :
\begin{enumerate}
\item Releve dans le texte une structure \zh{因为……所以……} et traduis-la.
\item Trouve dans le texte deux mots du champ lexical de la ville et deux mots du champ lexical de la campagne.
\item Traduis en chinois : « Je pense que la vie à la campagne est aussi très bonne. »
\end{enumerate}
\end{exercice}

\begin{exercice}[Thème et version]
\textbf{A. Thème.} Traduis en chinois (caractères obligatoires, pinyin en plus si tu veux) :
\begin{enumerate}
\item Beaucoup de jeunes quittent la campagne pour aller en ville.
\item Parce qu'il n'y a pas de travail au village, il est allé à Douala.
\item Je ne suis pas d'accord : la vie à la campagne est très agréable aussi.
\end{enumerate}
\textbf{B. Version.} Traduis en français :
\begin{enumerate}
\item \zh{城市的机会很多，但是生活太贵了。}\\(Chéngshì de jīhuì hěn duō, dànshì shēnghuó tài guì le.)
\item \zh{我希望我的家乡发展得越来越好。}\\(Wǒ xīwàng wǒ de jiāxiāng fāzhǎn de yuè lái yuè hǎo.)
\end{enumerate}
\end{exercice}

\begin{exercice}[Expression orale et écrite type Bac]
\textbf{A. Exposé oral (préparation écrite).} Prépare un exposé d'une minute sur le thème \zh{我的家乡} (« Mon village natal »). Rédige ta fiche de préparation en respectant les critères suivants :
\begin{itemize}
\item une structure en 5 étapes : salutation, présentation du sujet, deux arguments, avis personnel, conclusion ;
\item au moins 3 formules fonctionnelles vues en classe (\zh{我认为……}、\zh{我同意／不同意……}、\zh{因为……所以……}、\zh{大家觉得怎么样？}) ;
\item 8 à 10 phrases en caractères chinois, avec le pinyin au brouillon ;
\item attention particulière aux tons lors de la présentation orale.
\end{itemize}
\textbf{B. Jeu de rôle.} À deux : l'un est journaliste, l'autre est un jeune migrant arrivé à Douala. Préparez une interview d'au moins 5 échanges (questions sur les raisons du départ, la vie en ville, les difficultés, les projets). Le pinyin est autorisé au brouillon, mais la présentation finale se fait en caractères et à l'oral.
\end{exercice}

\newpage

\coursec{Corrigés des exercices}

\begin{corrige}[Exercice 1 — QCM : le bon choix]
\begin{enumerate}
\item \textbf{b) la campagne.} \zh{农村} (nóngcūn) = « la campagne, les zones rurales ». La ville se dit \zh{城市} (chéngshì).
\item \textbf{a) un ouvrier agricole migrant.} \zh{农民工} (nóngmíngōng) : \zh{农民} (paysan) + \zh{工} (travailleur) ; désigne les paysans partis travailler en ville.
\item \textbf{a) \zh{我认为……}} (Wǒ rènwéi…) = « Je pense que… », formule pour donner son avis. \zh{对不起} = « pardon », \zh{没关系} = « ce n'est pas grave ».
\item \textbf{b) la cause et la conséquence.} \zh{因为……所以……} (yīnwèi… suǒyǐ…) = « parce que… donc… ».
\item \textbf{a) \zh{你呢？}} (Nǐ ne ?) = « Et toi ? », question de relance. \zh{谢谢} et \zh{再见} ne servent pas à relancer la parole.
\end{enumerate}
\end{corrige}

\begin{corrige}[Exercice 2 — Vrai ou faux ? Justifie]
\begin{enumerate}
\item \textbf{Faux.} \zh{城市} (chéngshì) signifie « la ville » ; la campagne se dit \zh{农村} (nóngcūn).
\item \textbf{Faux.} \zh{我同意你的看法。} (Wǒ tóngyì nǐ de kànfǎ.) signifie « Je suis d'accord avec ton point de vue ». Le désaccord se dit \zh{我不同意} (wǒ bù tóngyì).
\item \textbf{Vrai.} La structure est bien : \zh{因为 + cause，所以 + conséquence。} Exemple : \zh{因为工作少，所以我离开了农村。} (Yīnwèi gōngzuò shǎo, suǒyǐ wǒ líkāile nóngcūn.) « Parce qu'il y a peu de travail, j'ai quitté la campagne. »
\item \textbf{Faux.} \zh{机会} (jīhuì) signifie « occasion, opportunité » ; « problème » se dit \zh{问题} (wèntí).
\item \textbf{Vrai.} \zh{怎么样？} (Zěnmeyàng ?) = « Comment ? Qu'en penses-tu ? », sert à demander l'avis ou l'état de quelque chose.
\end{enumerate}
\end{corrige}

\begin{corrige}[Exercice 3 — Vocabulaire : complète le tableau]
\begin{center}
\begin{tabularx}{\linewidth}{|X|X|X|X|}
\hline
\rowcolor{popblueL} Caractères & Pinyin & Nature & Traduction \\
\hline
\zh{家乡} & jiāxiāng & nom & village (pays) natal \\
\hline
\zh{离开} & líkāi & verbe & quitter, partir de \\
\hline
\zh{工作} & gōngzuò & nom / verbe & travail ; travailler \\
\hline
\zh{生活} & shēnghuó & nom / verbe & vie ; vivre \\
\hline
\zh{方便} & fāngbiàn & adjectif & pratique, commode \\
\hline
\zh{希望} & xīwàng & verbe / nom & espérer ; espoir \\
\hline
\end{tabularx}
\end{center}
\textbf{Points de vigilance :} \zh{工作} et \zh{生活} peuvent être noms ou verbes selon le contexte : \zh{找工作} (zhǎo gōngzuò, chercher du travail — nom) ; \zh{我在杜阿拉工作} (wǒ zài Dù'ālā gōngzuò, je travaille à Douala — verbe). \zh{方便} est un adjectif : \zh{很方便} (hěn fāngbiàn, très pratique).
\end{corrige}

\begin{corrige}[Exercice 4 — Pinyin et caractères : associe]
\begin{center}
\begin{tabularx}{\linewidth}{|X|X|}
\hline
\rowcolor{popblueL} Caractères & Pinyin \\
\hline
1. \zh{农村} & b. nóngcūn \\
\hline
2. \zh{城市} & c. chéngshì \\
\hline
3. \zh{工作} & a. gōngzuò \\
\hline
4. \zh{机会} & e. jīhuì \\
\hline
5. \zh{生活} & d. shēnghuó \\
\hline
6. \zh{问题} & f. wèntí \\
\hline
\end{tabularx}
\end{center}
\textbf{Points de vigilance :} attention aux tons : \zh{农村} nóngcūn (2\textsuperscript{e} + 1\textsuperscript{er} ton), \zh{城市} chéngshì (2\textsuperscript{e} + 4\textsuperscript{e} ton), \zh{机会} jīhuì (1\textsuperscript{er} + 4\textsuperscript{e} ton). Une erreur de ton change le sens ou rend le mot incompréhensible.
\end{corrige}

\begin{corrige}[Exercice 5 — Traduction vers le français]
\begin{enumerate}
\item \zh{很多年轻人离开农村，去城市找工作。} (Hěn duō niánqīngrén líkāi nóngcūn, qù chéngshì zhǎo gōngzuò.) — « Beaucoup de jeunes quittent la campagne pour aller chercher du travail en ville. »
\item \zh{因为农村的工作机会很少，所以他们去了大城市。} (Yīnwèi nóngcūn de gōngzuò jīhuì hěn shǎo, suǒyǐ tāmen qùle dà chéngshì.) — « Parce qu'il y a très peu d'opportunités de travail à la campagne, ils sont allés dans les grandes villes. »
\item \zh{我认为城市的生活很方便，但是也很贵。} (Wǒ rènwéi chéngshì de shēnghuó hěn fāngbiàn, dànshì yě hěn guì.) — « Je pense que la vie en ville est très pratique, mais aussi très chère. »
\item \zh{我希望我的家乡越来越好。} (Wǒ xīwàng wǒ de jiāxiāng yuè lái yuè hǎo.) — « J'espère que mon village natal ira de mieux en mieux. »
\end{enumerate}
\textbf{Points de vigilance :} \zh{越来越 + adjectif} (yuè lái yuè…) = « de plus en plus… » ; \zh{去 + lieu + verbe} exprime le but du déplacement (« aller en ville \emph{pour} chercher du travail »).
\end{corrige}

\begin{corrige}[Exercice 6 — Dialogue à compléter puis à jouer]
Dialogue complété :
\begin{itemize}
\item B：\zh{因为} 农村的工作 \zh{机会} 很少，\zh{所以} 我来杜阿拉找工作。\\(Yīnwèi nóngcūn de gōngzuò jīhuì hěn shǎo, suǒyǐ wǒ lái Dù'ālā zhǎo gōngzuò.)\\« Parce qu'il y a très peu d'opportunités de travail à la campagne, je suis venu à Douala chercher du travail. »
\item A：你觉得城市的生活 \zh{怎么样}？\\(Nǐ juéde chéngshì de shēnghuó zěnmeyàng ?)\\« Comment trouves-tu la vie en ville ? »
\item B：很方便，但是太贵了。你 \zh{呢}？\\(Hěn fāngbiàn, dànshì tài guì le. Nǐ ne ?)\\« Très pratique, mais trop chère. Et toi ? »
\item A：我不 \zh{同意}。我觉得农村的生活也很好。\\(Wǒ bù tóngyì. Wǒ juéde nóngcūn de shēnghuó yě hěn hǎo.)\\« Je ne suis pas d'accord. Je trouve que la vie à la campagne est très bien aussi. »
\end{itemize}
\textbf{Justification :} la première réponse exige la paire \zh{因为……所以……} ; \zh{机会} complète \zh{工作机会} (« opportunités de travail ») ; \zh{怎么样} forme la question d'avis ; \zh{呢} relance la parole ; \zh{不同意} exprime le désaccord (négation \zh{不} devant le verbe).
\textbf{Conseil oral :} soigner le 3\textsuperscript{e} ton de \zh{很} (hěn), le 4\textsuperscript{e} ton de \zh{贵} (guì) et le ton montant de \zh{同} (tóng).
\end{corrige}

\begin{corrige}[Exercice 7 — Discrimination des tons]
Réponses attendues (le mot correct de la leçon est le premier de chaque paire) :
\begin{enumerate}
\item \zh{农村} \textbf{nóngcūn} (2\textsuperscript{e} ton) — nòngcūn (4\textsuperscript{e} ton) est la forme fautive.
\item \zh{工作} \textbf{gōngzuò} (1\textsuperscript{er} ton) — gǒngzuò (3\textsuperscript{e} ton) est fautif.
\item \zh{城市} \textbf{chéngshì} (2\textsuperscript{e} ton) — chěngshì (3\textsuperscript{e} ton) est fautif.
\item \zh{机会} \textbf{jīhuì} (1\textsuperscript{er} ton) — jǐhuì (3\textsuperscript{e} ton) est fautif.
\item \zh{同意} \textbf{tóngyì} (2\textsuperscript{e} ton) — tǒngyì (3\textsuperscript{e} ton) est fautif.
\item \zh{生活} \textbf{shēnghuó} (1\textsuperscript{er} ton) — shènghuó (4\textsuperscript{e} ton) est fautif.
\end{enumerate}
\textbf{Méthode :} pour chaque paire, la seule différence est le ton de la première syllabe. À l'oral, exagère les tons : ton 1 haut et plat, ton 2 qui monte, ton 3 qui descend puis remonte, ton 4 qui tombe net. Un ton faux peut faire comprendre un autre mot : c'est l'erreur n° 1 des francophones.
\end{corrige}

\begin{corrige}[Exercice 8 — Réagir et relancer]
\textbf{A.} Réponses possibles (toute réaction justifiée par \zh{因为} est acceptée) :
\begin{enumerate}
\item \zh{我同意。因为城市有医院、学校和商店，生活很方便。}\\(Wǒ tóngyì. Yīnwèi chéngshì yǒu yīyuàn, xuéxiào hé shāngdiàn, shēnghuó hěn fāngbiàn.)\\« Je suis d'accord. Parce qu'en ville il y a des hôpitaux, des écoles et des magasins, la vie est très pratique. »\\(ou : \zh{我不同意。因为城市的生活太贵了。} « Je ne suis pas d'accord, parce que la vie en ville est trop chère. »)
\item \zh{我不同意。因为农村也有工作机会，比如农业工作。}\\(Wǒ bù tóngyì. Yīnwèi nóngcūn yě yǒu gōngzuò jīhuì, bǐrú nóngyè gōngzuò.)\\« Je ne suis pas d'accord. Parce qu'à la campagne aussi il y a du travail, par exemple dans l'agriculture. »
\item \zh{我不同意。因为农村的生活也很好，年轻人也可以在家乡发展。}\\(Wǒ bù tóngyì. Yīnwèi nóngcūn de shēnghuó yě hěn hǎo, niánqīngrén yě kěyǐ zài jiāxiāng fāzhǎn.)\\« Je ne suis pas d'accord. Parce que la vie à la campagne est bien aussi, les jeunes peuvent aussi se développer dans leur village. »
\end{enumerate}
\textbf{B.} Questions de relance :
\begin{enumerate}
\item (modèle) \zh{我喜欢城市。你呢？}
\item \zh{他觉得农村的生活很好。你呢？} ou \zh{你觉得怎么样？}\\(Tā juéde nóngcūn de shēnghuó hěn hǎo. Nǐ ne ? / Nǐ juéde zěnmeyàng ?)\\« Il trouve que la vie à la campagne est très bien. Et toi ? / Qu'en penses-tu ? »
\item \zh{我想去杜阿拉工作。你呢？} ou \zh{你想去哪里工作呢？}\\(Wǒ xiǎng qù Dù'ālā gōngzuò. Nǐ ne ?)\\« Je voudrais aller travailler à Douala. Et toi ? »
\end{enumerate}
\textbf{Points de vigilance :} la négation de \zh{同意} est \zh{不同意} (et non \zh{没同意}) ; \zh{呢} se place en fin de phrase pour relancer.
\end{corrige}

\begin{corrige}[Exercice 9 — Compréhension de texte et questions de langue]
\textbf{Questions de compréhension :}
\begin{enumerate}
\item Ali vient du nord du Cameroun (\zh{我的家乡在喀麦隆的北方}) ; il habite maintenant à Douala, où il est arrivé il y a trois ans.
\item Il a quitté son village parce qu'il y avait très peu d'opportunités de travail à la campagne (\zh{因为农村的工作机会很少}) et qu'il voulait trouver du travail en ville.
\item Avantages : la vie en ville est pratique — il y a des hôpitaux, des écoles et beaucoup de magasins. Inconvénients : la vie est chère et les loyers sont élevés (\zh{房租也很高}) ; de plus, son village lui manque (\zh{我常常想家}).
\item Pour l'avenir, Ali espère retourner dans son village natal et l'aider à se développer (\zh{我希望以后能回家乡，帮助家乡发展}).
\end{enumerate}
\textbf{Questions de langue :}
\begin{enumerate}
\item Structure relevée : \zh{因为农村的工作机会很少，所以我想在城市找工作。}\\(Yīnwèi nóngcūn de gōngzuò jīhuì hěn shǎo, suǒyǐ wǒ xiǎng zài chéngshì zhǎo gōngzuò.)\\« Parce qu'il y a très peu d'opportunités de travail à la campagne, je voulais trouver du travail en ville. »
\item Champ lexical de la ville : \zh{城市} (ville), \zh{公司} (entreprise), \zh{医院} (hôpital), \zh{商店} (magasin) — deux suffisent. Champ lexical de la campagne : \zh{农村} (campagne), \zh{家乡} (village natal).
\item « Je pense que la vie à la campagne est aussi très bonne. » → \zh{我认为农村的生活也很好。}\\(Wǒ rènwéi nóngcūn de shēnghuó yě hěn hǎo.)
\end{enumerate}
\textbf{Points de vigilance :} \zh{来到} (láidào, arriver à) et \zh{离开} (líkāi, quitter) sont des verbes directionnels ; \zh{能} (néng) exprime la possibilité (« pouvoir » rentrer).
\end{corrige}

\begin{corrige}[Exercice 10 — Thème et version]
\textbf{A. Thème :}
\begin{enumerate}
\item \zh{很多年轻人离开农村，去城市。}\\(Hěn duō niánqīngrén líkāi nóngcūn, qù chéngshì.)\\(variante plus complète : \zh{很多年轻人离开农村，去城市找工作。})
\item \zh{因为农村没有工作，所以他去了杜阿拉。}\\(Yīnwèi nóngcūn méiyǒu gōngzuò, suǒyǐ tā qùle Dù'ālā.)
\item \zh{我不同意：农村的生活也很舒服。}\\(Wǒ bù tóngyì : nóngcūn de shēnghuó yě hěn shūfu.)
\end{enumerate}
\textbf{B. Version :}
\begin{enumerate}
\item \zh{城市的机会很多，但是生活太贵了。} — « En ville, les opportunités sont nombreuses, mais la vie est trop chère. »
\item \zh{我希望我的家乡发展得越来越好。} — « J'espère que mon village natal se développera de mieux en mieux. »
\end{enumerate}
\textbf{Points de vigilance :} la négation de \zh{有} est toujours \zh{没有} (jamais \zh{不有}) ; dans \zh{发展得越来越好}, le complément de degré exige \zh{得} (et non \zh{的}) ; \zh{去了} marque l'accomplissement du déplacement.
\end{corrige}

\begin{corrige}[Exercice 11 — Expression orale et écrite type Bac]
\textbf{Barème indicatif (sur 20) :} respect de la structure en 5 étapes (4 pts) ; au moins 3 formules fonctionnelles correctement employées (4 pts) ; correction grammaticale et caractères (4 pts) ; richesse du vocabulaire du thème (4 pts) ; prononciation et tons à l'oral (4 pts).

\textbf{A. Proposition de fiche modèle — \zh{我的家乡}}
\begin{itemize}
\item Salutation : \zh{大家好！我叫玛丽。}(Dàjiā hǎo ! Wǒ jiào Mǎlì.) « Bonjour à tous ! Je m'appelle Marie. »
\item Présentation du sujet : \zh{今天我要谈一谈我的家乡。}(Jīntiān wǒ yào tán yi tán wǒ de jiāxiāng.) « Aujourd'hui, je vais vous parler de mon village natal. »
\item Argument 1 : \zh{我的家乡在喀麦隆的西部。农村的生活很安静，也很舒服。}(Wǒ de jiāxiāng zài Kāmàilóng de xībù. Nóngcūn de shēnghuó hěn ānjìng, yě hěn shūfu.) « Mon village est dans l'ouest du Cameroun. La vie à la campagne est calme et agréable. »
\item Argument 2 : \zh{但是，因为工作机会很少，所以很多年轻人离开家乡，去大城市找工作。}(Dànshì, yīnwèi gōngzuò jīhuì hěn shǎo, suǒyǐ hěn duō niánqīngrén líkāi jiāxiāng, qù dà chéngshì zhǎo gōngzuò.) « Mais, parce que les opportunités de travail sont rares, beaucoup de jeunes quittent le village pour chercher du travail dans les grandes villes. »
\item Avis personnel : \zh{我认为年轻人也可以在家乡发展。}(Wǒ rènwéi niánqīngrén yě kěyǐ zài jiāxiāng fāzhǎn.) « Je pense que les jeunes peuvent aussi se développer dans leur village. »
\item Conclusion : \zh{我希望我的家乡越来越好。大家觉得怎么样？谢谢！}(Wǒ xīwàng wǒ de jiāxiāng yuè lái yuè hǎo. Dàjiā juéde zěnmeyàng ? Xièxie !) « J'espère que mon village ira de mieux en mieux. Qu'en pensez-vous ? Merci ! »
\end{itemize}
Formules fonctionnelles utilisées : \zh{我认为……}、\zh{因为……所以……}、\zh{大家觉得怎么样？} — le critère des 3 formules est rempli.

\textbf{B. Proposition de jeu de rôle (interview, 5 échanges)}
\begin{itemize}
\item Journaliste : \zh{你好！请问，你叫什么名字？}(Nǐ hǎo ! Qǐngwèn, nǐ jiào shénme míngzi ?) « Bonjour ! Comment t'appelles-tu, s'il te plaît ? »
\item Migrant : \zh{我叫保罗，我来自喀麦隆的北方。}(Wǒ jiào Bǎoluó, wǒ láizì Kāmàilóng de běifāng.) « Je m'appelle Paul, je viens du nord du Cameroun. »
\item Journaliste : \zh{你为什么离开家乡？}(Nǐ wèishénme líkāi jiāxiāng ?) « Pourquoi as-tu quitté ton village ? »
\item Migrant : \zh{因为农村的工作机会很少，所以我来杜阿拉找工作。}(Yīnwèi nóngcūn de gōngzuò jīhuì hěn shǎo, suǒyǐ wǒ lái Dù'ālā zhǎo gōngzuò.) « Parce qu'il y a peu de travail à la campagne, je suis venu à Douala chercher du travail. »
\item Journaliste : \zh{你觉得城市的生活怎么样？}(Nǐ juéde chéngshì de shēnghuó zěnmeyàng ?) « Comment trouves-tu la vie en ville ? »
\item Migrant : \zh{很方便，但是太贵了，房租也很高。}(Hěn fāngbiàn, dànshì tài guì le, fángzū yě hěn gāo.) « Très pratique, mais trop chère, et les loyers sont élevés. »
\item Journaliste : \zh{你以后想回家乡吗？}(Nǐ yǐhòu xiǎng huí jiāxiāng ma ?) « Veux-tu retourner dans ton village plus tard ? »
\item Migrant : \zh{想。我希望回家乡，帮助家乡发展。}(Xiǎng. Wǒ xīwàng huí jiāxiāng, bāngzhù jiāxiāng fāzhǎn.) « Oui. J'espère retourner au village et l'aider à se développer. »
\item Journaliste : \zh{谢谢你！再见！}(Xièxie nǐ ! Zàijiàn !) « Merci ! Au revoir ! »
\end{itemize}
\textbf{Points de vigilance :} (1) ne pas oublier la paire complète \zh{因为……所以……} ; (2) \zh{为什么} pour « pourquoi », réponse avec \zh{因为} ; (3) soigner les tons : \zh{杜阿拉} Dù'ālā, \zh{房租} fángzū (2\textsuperscript{e} + 1\textsuperscript{er} ton) ; (4) à l'oral, parler lentement, phrase par phrase, et regarder l'auditoire ; (5) la présentation finale se fait sans pinyin : seuls les caractères et la mémoire des tons comptent.
\end{corrige}

\newpage

\coursec{Fiche bilan}

\begin{popfigure}
\begin{tikzpicture}[mindmap, grow cyclic, every node/.style={concept, font=\small},
  root concept/.append style={concept color=popblue, font=\bfseries, text=white},
  level 1/.append style={level distance=3.2cm, sibling angle=51.4},
  level 1 concept/.append style={font=\footnotesize}]
\node[root concept, align=center]{Exode rural\\ 农村人口外流}
  child[concept color=poporange]{ node[align=center]{Ville / campagne\\ 城市 chéngshì\\ 农村 nóngcūn} }
  child[concept color=popgreen]{ node[align=center]{Causes\\ 找工作 zhǎo gōngzuò\\ 上学 shàngxué} }
  child[concept color=poppurple]{ node[align=center]{Conséquences\\ 老人、孩子留下\\ 城市人太多} }
  child[concept color=popgold]{ node[align=center]{Solutions\\ 发展农村\\ 修路、建学校} }
  child[concept color=poppink]{ node[align=center]{Donner son avis\\ 我觉得……\\ 我认为……} }
  child[concept color=popblue!70]{ node[align=center]{Relancer le débat\\ 你呢？\\ 你同意吗？} }
  child[concept color=popgreen!70]{ node[align=center]{Prononciation\\ 4 tons + ton neutre} };
\end{tikzpicture}
\legende{Carte mentale de la leçon : parler de l'exode rural en chinois}
\end{popfigure}

\begin{bilan}
\textbf{1. Lexique clé à connaître par cœur}
\begin{itemize}
  \item \zh{农村} nóngcūn : la campagne \quad \zh{城市} chéngshì : la ville \quad \zh{农村人口外流} nóngcūn rénkǒu wàiliú : l'exode rural
  \item \zh{离开} líkāi : quitter \quad \zh{搬到} bāndào : déménager à \quad \zh{工作} gōngzuò : travail \quad \zh{机会} jīhuì : occasion, opportunité
  \item \zh{生活} shēnghuó : vie \quad \zh{方便} fāngbiàn : pratique \quad \zh{发展} fāzhǎn : développer \quad \zh{问题} wèntí : problème
  \item \zh{越来越多} yuè lái yuè duō : de plus en plus nombreux \quad \zh{老人} lǎorén : personne âgée \quad \zh{孩子} háizi : enfant
\end{itemize}

\textbf{2. Structures et formules fonctionnelles}
\begin{center}
\begin{tabularx}{\linewidth}{|X|X|X|}
\hline
\rowcolor{popblueL} Fonction & Structure chinoise + pinyin & Traduction \\
\hline
Présenter son sujet & \zh{今天我要谈谈……} Jīntiān wǒ yào tántan…… & Aujourd'hui, je vais parler de…… \\
\hline
Donner son avis & \zh{我觉得 / 我认为……} Wǒ juéde / Wǒ rènwéi…… & Je pense que…… \\
\hline
Exprimer une cause & \zh{因为……所以……} yīnwèi…… suǒyǐ…… & Parce que…… donc…… \\
\hline
De plus en plus & \zh{越来越 + adjectif} yuè lái yuè duō & De plus en plus + adj. \\
\hline
Être d'accord & \zh{我同意你的看法。} Wǒ tóngyì nǐ de kànfǎ. & Je suis d'accord avec toi. \\
\hline
Ne pas être d'accord & \zh{我不太同意，因为……} Wǒ bú tài tóngyì, yīnwèi…… & Je ne suis pas tout à fait d'accord, car…… \\
\hline
Relancer le débat & \zh{你呢？你怎么看？} Nǐ ne ? Nǐ zěnme kàn ? & Et toi ? Qu'en penses-tu ? \\
\hline
Conclure & \zh{总之，……} Zǒngzhī,…… & En résumé,…… \\
\hline
\end{tabularx}
\end{center}

\textbf{3. Méthode express pour l'exposé oral}
\begin{enumerate}
  \item \textbf{Introduction} : saluer (\zh{大家好！}) puis annoncer le sujet avec \zh{今天我要谈谈……}.
  \item \textbf{Développement} : 2 ou 3 idées simples (causes, conséquences, solutions), reliées par \zh{因为……所以……}, \zh{但是}, \zh{还有}.
  \item \textbf{Conclusion} : donner son avis avec \zh{我觉得……} et terminer par \zh{谢谢大家！}.
  \item \textbf{Prononciation} : soigner les 4 tons, parler lentement, phrase par phrase.
\end{enumerate}
\end{bilan}

\begin{aretenir}[Checklist avant le Bac]
\begin{itemize}
  \item[$\square$] Je connais le lexique ville / campagne : \zh{城市}, \zh{农村}, \zh{农村人口外流}.
  \item[$\square$] Je sais annoncer mon sujet : \zh{今天我要谈谈……}.
  \item[$\square$] Je sais donner mon avis : \zh{我觉得……} et \zh{我认为……}.
  \item[$\square$] Je sais exprimer une cause avec \zh{因为……所以……} (jamais « suǒyǐ » seul au début).
  \item[$\square$] Je sais utiliser \zh{越来越 + adjectif} : \zh{越来越多人搬到城市。}.
  \item[$\square$] Je sais réagir : \zh{我同意。} / \zh{我不太同意，因为……}.
  \item[$\square$] Je sais relancer le débat : \zh{你呢？你怎么看？}.
  \item[$\square$] Je prononce correctement les 4 tons, surtout \zh{nóngcūn} (2\textsuperscript{e} + 1\textsuperscript{er} ton).
  \item[$\square$] Je place le complément de lieu avant le verbe : \zh{从农村搬到城市}.
  \item[$\square$] Je sais conclure : \zh{总之，……谢谢大家！}.
  \item[$\square$] Je peux parler 2 minutes sans lire mes notes.
  \item[$\square$] Je donne un exemple camerounais : \zh{很多年轻人去雅加达、杜阿拉……} non — \zh{很多年轻人去雅温得、杜阿拉。} Hěn duō niánqīngrén qù Yǎwēndé, Dù'ālā.
\end{itemize}
\end{aretenir}

\findecours

\end{document}
